%% ---------------------------------------------------------------
%% $URL: https://repository.cs.ru.is/svn/thesis-template/trunk/ruthesis/latex/DEGREE-NAME-YEAR.tex $
%% $Id: DEGREE-NAME-YEAR.tex 289 2016-02-17 12:49:46Z foley $
%% This is a template LaTeX file for theses or reports at Reykjavík University
%% Particularly the Computer Science or Science and Engineering departments
%% The template and ruthesis.cls file have been designed to be compatible with
%%   LaTeX, PDFTeX, and XeTeX.
%%
%% Comments and questions can be sent to
%% General: the RU LaTeX group (latex AT list.ru.is)
%% Computer Science: the Graduate Study Council (cs-grad AT ru.is)
%% Science and Engineering: graduate administrator: Kristín Kristinsdóttir (kristinkri AT ru.is)
%% See the README for information on final preparation for PhD Dissertations
%%
%% Original Author: Bjórn Þór Jónsson (bjorn AT ru.is)
%% Modified by: Eyjólfur Ingi Ásgeirsson <eyjo AT ru.is>
%% Modified and Maintained by: Joseph Timothy Foley <foley AT ru.is>
%% ---------------------------------------------------------------

%% METHOD:
%% 0) Read README.TXT
%% 0.1) !!!NO REALLY, READ README.TXT!!!
%% 0.2) Subscribe to the announcements email list at
%%    https://list.ru.is/mailman/listinfo/latex-announcements
%% 1 LaTeX instructions.tex or goto http://afs.rnd.ru.is/project/thesis-template/trunk/ruthesis/latex/instructions.pdf
%% 2) Copy the template (or unzip) to your working area
%% 3) Rename this file (if needed) with your information e.g. MSC-FOLEY-2007.tex
%% 4) Modify this file to fit your needs (please follow all comments below in the text)
%% 5) For making bibliographies, run "biber".  You can also change
%% this back to "bibtex".  See below in "Bibliography options".

%% Options for ruthesis class:

%%%%%%% CHOOSE ONE OF THESE %%%%%%%%%%%%%%%
%% projectreport: Project report (CS)
%% bachelors: Bachelor of Science thesis
%% masters: Master of Science thesis
%% doctorate: Doctor of Philosophy dissertation
%%
%%%%%%% CHOOSE ONE OF THESE %%%%%%
%%
%% draft: speed up processing by skipping graphics and adding useful
%%     information for editing.  Also sets spacing to double so that it is easier to
%%     write editing marks on paper copy.
%% proof:  proofreading version (final formatting with warnings)
%% final: generate document for submission, removing FIXMEs, and
%%     other markup.  Throw error if any fatal FIXMEs still in document.
%%
%%%%%%%% CHOOSE ANY COMBINATION OF THESE %%%%%%%%%%%%
%%
%% booklet: set the stock/font size for booklets (experimental, DO NOT USE)
%% forcegraphics: force graphics, etc. to be included, even in draft mode
%% debug:  writes more messages to the log file, adds debugging output
%%     and sizing boxes
%% miktex: disable some packages that cause problems for (older) MikTeX
%% icelandic: thesis is in Icelandic
%% oldstyle:  use the PhD headers and footers from the old CS template
%% online: for online versions (skip additional signature pages)
\documentclass[masters,forcegraphics,proof]{ruthesis}

%%%%%%% HACKS:  USE THESE ONLY IF INSTRUCTED TO DO SO!!!! %%%%%%%%%%%%%%%
%% yngvihack:  Proposed changes to the template from Dean Yngvi
%\setboolean{yngvihack}{true}

%% Note, if you are writing a PhD thesis, it will be reduced at the
%% printers when it is bound into a book: B5 240x170mm (aka Programme
%% or Book economy) See the README.TXT for details.  Plan your figures
%% and diagrams for this size.  In a future release of this template,
%% setting booklet will allow you to see what it would look like.

%% For other thesis, print them out in A4, then they will be cut down
%% to 293x206mm by the binder.  Make sure you use "acid-free" paper or
%% the book will yellow and fall apart in a few years.
%% TODO: instructions to the printer about bleed

%%%%%%%%%%%%%%%%%%%% Bibliography options %%%%%%%%%%%%%%%%%%%%%
%% We suggest switching from bibtex to biblatex/biber because it is better able
%% to deal with Icelandic characters and other bibliography issues
%% As long as you use biblatex instead of bibtex by itself, it will at least
%%  generate a document without errors.
\usepackage[backend=biber,bibencoding=utf8,style=ieee, sorting=nty]{biblatex}
%\DeclareLanguageMapping{american}{american-apa}
% need to declare mapping for style=apa to alphabetize properly
% If you set backend=bibtex, it will use bibtex for processing (old way)
%    this can work with Icelandic characters, but you may get weird results.
%    bibtex does not know how to sort Þ and ð
% if you set backend=biber, you can use UTF8 characters such as Þ and
%     ð  but you will have to remember to switch from using bibtex to
%     biber in your client
% If you use JabRef, make sure the file is encoded in UTF-8 which is
%    not the default.

% Where is your reference library?
\addbibresource{references.bib}

%%%%%%%%%%%%%%%%%%% CUSTOMIZATIONS %%%%%%%%%%%%%%%%%%%%%%%%%%%%%
%% It is not recommended that you customize this file
%%   nor ruthesis.cls.   You should put your macros and packages into
%%   a separate file so that it is easier to use updates to the template.
%% The custom.tex file was created for this reason.
%% We load this much later so that it can overrite any existing settings
\InputIfExists{custom.tex}

%%%%%%%%%%%%%%% INFORMATION %%%%%%%%%%%%%%%%%%5
%% University information must be multilingual to deal with the
%%  required cover pages and abstract on thesis
%% NOTE: This may not be required for other reports!!!

%% Babel Icelandic macros are setup  on RedHat at
%% /usr/share/texlive/texmf-dist/tex/generic/babel/icelandic.sty
%% /usr/share/texlive/texmf-dist/tex/generic/babel-icelandic/icelandic.ldf


%% Multilingual macros
%\newML{macroname}{englishword}{icelandicword}
%  creates \macronameML
%    \MLmacroname[english] - returns the english word
%    \MLmacroname[icelandic] - returns the icelandic word
%    \MLmacroname  - uses the current language setting
% Some useful ones have already been defined, but can be redefined
%% Predefined: \MLIceland \MLReykjavikUniversity \MLUniversityIceland

%% What institute?  Default is RU.
%\setInstitution{\MLReykjavikUniversity}
% \newML{InstitutionAddress}{Menntavegur 1\\101 Reykjavík, Iceland}
% {Menntavegi 1\\101 Reykjavík, Ísland}
% \setInstitutionAddress{\MLInstitutionAddress}
% \newML{Tel}{Tel.}{Sími}
% \setInstitutionPhone{\MLTel{} +354 599 6200\\
% Fax +354 599 6201}
% \setInstitutionURL{www.ru.is}


%% Department and degree program

%% School of Computer Science: \MLCS
% Programs:  Computer Science, Software Engineering, Language Technology
%\program{Computer Science}

%% School of Science and Engineering \MLSSE
%% Programs:
%% Bioinformatics, Biomedical Engineering
%% Civil Engineering with specialization in Concrete Technology
%% Civil Engineering with specialization in Construction Management
%% Civil Engineering with specialization in Structural Design
%% Civil Engineering with specialization in Transport and Urban Planning
%% Construction Management
%% Electrical Engineering
%% Mechatronics
%% Engineering Management
%% Exercise Science and Coaching
%% Financial Engineering
%% Mechanical Engineering
%% Sustainable Energy Engineering - REYST
%% Sustainable Science - REYST
%% Urban Planning and Transport

\setSchool{\MLSSE}
\newML{program}{Computer Science}{tölvunarfræði}
\program{\MLprogram}

%% Degree long name.  If not already defined, you can create a macro
%\newML{DEGREE}{English Degree Name}{Icelandic Degree Name}
%% Default is set based upon doctorate vs masters option
%% Predefined: \MLMSc \MLPhd
%\setDegreelong{\MLMSc}

%% Degree abb, change if default is not right
%% Default is set based upon doctorate vs masters option
%\degreeabbrv{Sc.D.}


%\setFrontLogo{reyst-logo}
%% Use this if you need a different front logo on the first page
%% e.g. reyst-logo

%% Date in english and icelandic
%% NOTE: THIS IS THE DATE OF THE SUPERVISOR'S SIGNATURE!!!!!!
%% Predefined: \MLjan, \MLfeb, \MLmar, ... \MLdec
%\whensigned{day}{month}{year} %day is only used on some formats, but you must put something.
\whensigned{10}{\MLapr}{2016}

%% Title in english and icelandic
%% This is required on Thesis, but not necessarily elsewhere.
%%
%% Note that title length is limited to what can fit in three lines in the inside page
%% Also note that the capitalization of the text on the front page is a potential source
%% of latex errors as it may not deal well with math, international letters,
%% and other latex constructs
%% Warning: capitalization does not always work anyway
\newML{Title}{Equivalence classes of mesh patterns with a dominating pattern}{Jafngildisflokkar möskvamynstra með ríkjandi mynstri}
\newML{TITLE}{EQUIVALENCE CLASSES OF MESH PATTERNS WITH A DOMINATING PATTERN}{JAFNGILDISFLOKKAR MÖSKVAMYNSTRA MEÐ RÍKJANDI MYNSTRI}
% YFIRGNÆFANDI
%% If you have the title in just one language, put it into \setTitle{} directly e.g.
%%  \setTitle{Working Title}
%%
\setTitle{\MLTitle}
% \setTitle{Equivalence Classes of Mesh Patterns with a Dominating Pattern}
%% If you need to adjust the main title later, you can use
%% \setMainTitle{}
%% ***** Special Titles ******
%% If the title must be formatted specifically for the cover page or internal pages
%% (typically via line-breaks using the \newline command) then the following commands must be used
%%
%% This one for the online/PhD front page:
%%  and should usually be kept short per line
%% The template specified that a PhD cover should have the title in capitals, but
%%  the physical copies have not done this.  This macro is in case that
%%  turns out to be important
%\setTitleCover{\MLTITLE}
%% These two for the internal cover pages, usually not needed
%\newML{TitleInternal}{Internal Title}{Icelandic Internal Title}
%\setTitleInternal{\MLTitleInternal}

%% Author name (should be the same in any language, if not use \newML)
%% If you are writing a Project report with multiple authors, separate them with \\:
%% To keep the names typeset together, you want to use non-breaking spaces
%\author{Firstname1~Lastname1\\Firstname2~Lastname2}
\author{Murray~Tannock}

%% If the name must be formatted specifically for the signature page
%% (typically via line-breaks) then the following command must be used
%\setAuthorSignature{Student\\Name}
%% This macro adjusts the author name in the headers of the oldstyle formatting
%\setAuthorHeader{StudentLast}

%%%%%%%%%%%%%%%%%%%%%%%%%% Project Report or Bachelor's Only!!! %%%%%%%%%%%%%%%%%%%%%%%%%%%%%%%%%%%%%%%%%%%
% \setCourse{VT LOK 1012}

%%%%%%%%%%%%%%%%%%%%%%%%%% Project Report Only!!! %%%%%%%%%%%%%%%%%%%%%%%%%%%%%%%%%%%%%%%%%%%


%%%%%%%%%%%%%%%%%%%%%%%%%% Bachelors Only!!! %%%%%%%%%%%%%%%%%%%%%%%%%%%%%%%%%%%%%%%%%%%
% \setID{010105--9999}%kennitala
% \setSemester{2015--1}

% \setOrganization{Marel ehf.\\Austurhrauni 9\\210 Garðabær}
% \setSubProgram{Tæknifræði}

%%%%%%%%%%%%%%%%%%%%%%%%%%% Masters Only!! %%%%%%%%%%%%%%%%%%%%%%%%%%%%%%%%%%%%%%%%%%%%
%% How many credits (ECTS) on Master's degree
%% Usually 30 or 60
\ects{60}

%%%%%%%%%%%%%%%%%%%%%%%%%%% Doctorate Only!! %%%%%%%%%%%%%%%%%%%%%%%%%%%%%%%%%%%%%%%%%%
%% Some Computer Science Thesis have an ISSN number.
%% Most other documents do not.
%\bookidnumber{ISSN: 1670-8539}
%% ID numbers are optional, but nice for sorting in libraries

%% International Standard Book Number (ISBN)
%% This is what most people should use if the thesis is being published.

%% International Standard Serial Number (ISSN)
%% This is usually only for a PhD dissertation as part of a series when published
%%   Computer Science: 1670-8539

%% Additional degrees?  (optional, usually not needed)
%\adddegree{(list of degrees in appendix)}{(sjá lista yfir prófgraður í viðauka)}
%%%%%%%%%%%%%%%%%%%%%%%%%%%%%%%%%%%%%%%%%%%%%%%%%%%%%%%%%%%%%%%%%%%%%%%%%%%%%%%%%%%%%%%%


%% List the entire committee.  Each member has a name (degree should be omitted, unless it is not PhD),
%% Supervisor(s) must appear first
%% On a Bachelors, there is usually only one supervisor and one examiner.

%% Format for each entry:
%%  \personinfo{Name}{Role}{Job Title}{Company/institution}{Country}
%% Predefined macros: \MLSupervisor \MLSupervisors \MLExaminer \MLExaminers

%% Change these to singular/plural as needed.
%% Just uncomment and change the plurality of the macro.
%\setSupervisorHeading{\MLSupervisors}
%\setExaminerHeading{\MLExaminer}

%% Predefined macros:
%% \MLSeniorProfessor \MLProfessor \MLAssociateProfessor \MLAdjunctProfessor \MLEmeritusProfessor \Iceland
%% \MLReykjavikUniversity \MLUniversityIceland

%% Bachelors: primary advisor (Umsjónarkennari), ONLY ONE!
%% All others: As many as you want
\supervisors{
  \personinfo{Henning Ulfarsson}{\MLSupervisor}{\MLProfessor}{\MLReykjavikUniversity}{\MLIceland}
%  \personinfo{Helpful A. Teacher}{Co-advisor}{\MLAssistantProfessor}{\MLUniversityIceland}{\MLIceland}
%  \personinfo{Ian M. Great}{Co-advisor}{\MLProfessor}{Hochschule Düsseldorf}{Germany}
}

%% Bachelors: secondary advisor (Leiðbeinandi), ONLY ONE
%% All others: As many as you want
\examiners{
  \personinfo{Michael H. Albert}{\MLExaminer}{\MLProfessor}{University of Otago}{New Zealand}
  \personinfo{Anders Claesson}{\MLExaminer}{Docent}{\MLUniversityIceland}{\MLIceland}

}

%% An abstract is required to be in both Icelandic and English for most degrees.
%% It is considered good form to limit the abstract to a single paragraph in each language,
%%   around 250 words or so.  Refer to your degree's instructions.
%% Note: Icelandic quotation marks cannot be typeset using "` and "'.  You should use \enquote{}
%%   this is probably due to interactions with the MultiLingual macros.
\newML{AbstractText}{A permutation is an arrangement of \(n\) objects.
Avoidance classes of classical permutation patterns
capture many interesting properties, such as
stack-sortability, and have links to many different combinatorial objects.
Mesh patterns are an extension of classical patterns that allow additional
restrictions to be placed on occurrences of the pattern.
Two mesh patterns are coincident if they are avoided by the same set of
permutations. In this thesis, the author provides sufficient conditions for
coincidence among mesh patterns, whilst also avoiding a longer classical
pattern. These conditions, along with two special cases are used to completely
classify coincidence amongst families containing a mesh pattern of length 2
and a classical pattern of length 3. Two patterns are Wilf-equivalent if they
have the same number of avoiders at every length, the author completely
Wilf-classifies mesh patterns of length 2 when avoiding the classical pattern 231.
% Finally we attempt to show some non-trivial Wilf-equivalences between avoiders
% of sets of the form \(231,m_1\) and \(231,m_2\)
}
% ipsum generaes text text
{\lipsum[1]} % Icelandic abstract goes here
\setAbstract{\MLAbstractText}


%%%%%%%%%%%%%%INDEX SETUP %%%%%%%%%%%%%%%%%%%%%%%%%%%%%%%%%%%%%%%%%%%%%%%%%%%%
%% Indexes, and other auto-generated material
%% You will need two components:
%% 1) a program to extract the index entries and sort them
%% 2) LaTeX package to parse them

%% 1) Index generation programs: xindy, makeindex
%% Warning: makeindex and other traditional latex indexing
%%   does not know how to deal with accented characters
%% You will need to run this when you want to generate/update the index.
%%
%% For abbreviations, you may want to try
%%\usepackage[intoc]{nomencl} % intoc: In Table of Contents
%% remember to run:
%% makeindex filename.nlo  -s nomencl.ist -o filename.nls

%% 2) PICK ONE OF THESE INDEX PACKAGES: imakeidx, makeidx, index (is default)
%%  imakeidx does not need makeindex to be run seperately but
%%  it means that generating your document is slower when you are not updating the index
%\RequirePackage{imakeidx}
%\makeindex[intoc]  % add index to toc

\finalifforcegraphics{hyperref} %hyperlinks even in draft mode
\usepackage[hidelinks]{hyperref}
%% !!!Must be the last package loaded except otherwise mentioned!!!!
%% \usepackage{hypcap}  %% puts link at top of figure, must be after hyperref

%%%%%%%%%%%%%GLOSSARY SETUP%%%%%%%%%%%%%%%%%%%%
%% A glossary is a nice thing to have in larger documents.
%% Unfortunately, it seems to cause problems on many TeX installations
%%  so it is disabled by default.
%% Uncomment the lines that start with \% in order to re-enable it

% % add an autogenerated glossary (separate from main index)
%% Warning: package glossaries doesn't like it when they are defined in the document
%% http://mirror.ctan.org/macros/latex/contrib/glossaries/glossaries-user.html#sec%3adocdefs
%%
%% WARNING!  Memoir has an internal "index" option which conflicts with glossaries
%% Loading the package must go after hyperref (it is special).
% \RequirePackage[acronym,toc]{glossaries}
%% nomain: no main glossary
%% acronym: acronym glossary
%% toc: add entry in Table of Contents
%% automake: try to run makeindex using \write18 shell escape
%% without automake, run "makeglossaries" to process them
%% or the right makeindex options
% \newglossary[slg]{symbols}{sym}{sbl}{List of Symbols}
% \makeglossaries{}  % intitialize the glossary.  Must be in the main file!!!
% %% $URL$
%% $Id$
%% Glossary term definitions
%% Note, you will need to reference them in the text after defining them

\newglossaryentry{glossary}{name=glossary, description={a place to put all the terms that confuse people}}
\newglossaryentry{ohm}{name=ohm,symbol={\ensuremath{\Omega}}, description=unit of electrical resistance}  % Do not break this with newlines!
\newglossaryentry{dissertation}{name=dissertation, description={A large pile of paper that professors use to prop up shelves.}}
\newacronym{rfid}{RFID}{Radio Frequency IDentification}

\newacronym{led}{LED}{light-emitting diodes}

\newacronym{auv}{AUV}{Autonomous Undersea Vehicle}
%% Reference manual: http://ctan.uib.no/macros/latex/contrib/glossaries/glossaries-user.pdf

%% You can also define most entries near where they are used, but this may cause problems.  See manual chapter 4.8
%% Be careful about putting entries on Chapter, Section, etc.  See manual Chapter 6.



%%% Local Variables: 
%%% mode: latex
%%% TeX-master: "msc-tannock-2016"
%%% End: 
  % load glossary definitions in preamble

%%%%%%%%%%%%%%%%%%%%%%%%%%%%%%%%%%%%%%%%%%%%%%%%%%%%%%%%%%%%%%%%%%%%%%%%%%
%%%%%%%%%%%%%%%%%%%%%%% DOCUMENT START %%%%%%%%%%%%%%%%%%%%%%%%%%%%%%%%%%%

\begin{document}
%% Some elements have different names on the RU Masters rules
%% They will be annotated with RUM: "name"
\frontmatter{} % setup formatting at beginning

\frontcover{}%% The MSc/BSc cover can be gotten from the front desk and has a window in it
%%RUM: Front cover (standard)
\frontrequiredpages{}  % the various signature pages

% \ifdraft{\coverchapter{Important!!!  Read the Instructions!!!} If you
%   have not already done so, \LaTeX{} the \path{instructions.tex} to
%   learn how to setup your document and use some of the features.  You
%   can see a (somewhat recent) rendered PDF of the instructions at
%   \url{http://afs.rnd.ru.is/project/thesis-template/trunk/ruthesis/latex/instructions.pdf}.
%   There is also more information on working with \LaTeX{} at
%   \url{http://afs.rnd.ru.is/project/htgaru/trunk/how-to-get-around-projects.pdf}.
%   This includes common problems and fixes.

%   This page will disappear in anything other than draft mode.}{}



% \newML{child}{child}{barn}
% %% Dedication is optional, comment out if it is absent
% %% RUM: Not mentioned
% \begin{dedications}
%   I dedicate this to my spouse/\MLchild{}/pet/power animal.
% \end{dedications}

\enableindents{}% turn on/off paragraph indents
% RUM: "Acknowledgements (optional)"
\coverchapter{Acknowledgements}
\begin{quotation}
So long, and thanks for all the fish.
\end{quotation}\sourceatright{Douglas Adams\cite{adams84fish}}
\vspace{\baselineskip}

I would like to thank my advisor, Henning Ulfarsson, for his time, support,
guidance and introducing me to a research topic that I greatly enjoy.

I would like to thank Michael Albert and Anders Claesson for the discussions that
prompted the topic of this thesis, as well as their valuable input and suggestions.

I would like to thank Christian Bean for proofreading my thesis, providing useful
feedback, and for the time spent staring at a whiteboard while I tried to check
my mathematical reasoning.

I would like to thank Stefanía Andersen Aradóttir for company and sanity checks
during long working weekends and for making sure I didn't look like a fool in my
attempt at an Icelandic title.


\draftnote{Acknowledgements are optional; comment this chapter out if they are absent
  Note that it is important to acknowledge any funding that helped in the work}

% This work was funded by \the\year~RANNIS grant ``Survey of man-eating Minke whales'' 1415550.
% Additional equipment was generously donated by the Icelandic Tourism Board.

% \coverchapter{Preface}
% % RUM: "Preface (optional)"
% This dissertation is original work by the author, Firstname~Lastname.
% Portions of the introductory text are used with permission from
% Student et al.\cite{student2015awesome} of which I am an author.


% \draftnote{The preface is an optional element
%   explaining a little who performed what work.  See
%   \url{https://www.grad.ubc.ca/sites/default/files/materials/thesis_sample_prefaces.pdf}
%   for suggestions.

%   List of publications as part of the preface is
%   optional unless elements of the work have already been published.
%   It should be a comprehensive list of all publications in which
%   material in the thesis has appeared, preferably with references to
%   sections as appropriate.  This is also a good place to state
%   contribution of student and contribution of others to the work
%   represented in the thesis.}


%\coverchapter{Publications}
%% RUM: Not mentioned, this was found in the CS thesis template.
%% Maybe more applicable to PhD dissertations?
%%% Probably a duplication from before Preface became standard.


\starttables{} % setup formatting
%% TOC, list of figures and list of tables are required
\tableofcontents{}%%RUM: "Table of contents"
\thesislistoffigures{}%%RUM: "List of figures"
\thesislistoftables{}%%RUM: "List of tables"

%\coverchapter{List of drawings and enclosed material}
%RUM: "List of drawings and enclosed material, e.g. CD(as appropriate)"

\listoffixmes{}
% if using fixme package, lists what needs to be done

%% The list of abbreviations is an example of a special list
%% Other lists may be added, such as lists of algorithms, symbols, theorems, etc.
%% IN CS PhD, this is sometimes centered.
% \newtoclist{Abbreviations}%%RUM: Not mentioned
% \begin{tabular}{ll}
% MSc &Masters of Science\\
% PhD &Doctor of Philosophy\\
% \end{tabular}

% \newtoclist{Symbols}%%RUM: Not mentioned
% \begin{tabular}{lll}
% Symbol &Description &Value/Units\\
% $E$ &Energy &\si{\joule}\\
% $m$ &Mass &\si{gram}\\
% $c$ &Speed of Light &\SI{2.99E8}{\meter\per\second\square}\\
% \end{tabular}

%% This command prepares for the actual text, e.g. by
%% calling \mainmatter{}
\starttext{}

%% ---------------------------------------------------------------
%% From this point on, it is standard Latex, except the very end.
%% This is a "report"-based template, so the top-level heading
%% is \chapter{}

%% WARNING: Make sure that all of these files (and any new ones)
%% are UTF-8 otherwise you will get weird encoding errors.
%\part{The First Part} % Parts optional but useful in longer documents

%% The default division is IMRAD, you may want to divide differently
%% See the introduction for guidance.

%% This file only is included if we are draft mode.
%% The instructions.tex file contains instructions and examples.
%\ifdraft{\documentclass[projectreport,online,forcegraphics,miktex,draft]{ruthesis}
\usepackage[backend=biber,bibencoding=utf8,style=ieee]{biblatex}
\addbibresource{references.bib}
\InputIfExists{custom.tex}
\setTitle{Writing a Thesis in \LaTeX{}}
%% Because datetime is broken, we always have to set a \whensigned{day}{MLmonth}{year}
\whensigned{12}{\MLdec}{2016}
\author{Joseph Timothy Foley}
\setCourse{Final Project Reports, Thesis, and Dissertations}
\setSchool{\MLSSE{} \& \MLSCS{}}

%% Making glossaries needs perl installed
%% Usually only Windows needs this installed separately
%% http://www.activestate.com/activeperl
\RequirePackage[acronym,toc]{glossaries}
%% nomain: no main glossary
%% acronym: acronym glossary
%% toc: add entry in Table of Contents
%% automake: try to run makeindex using \write18 shell escape
%% without automake, run "makeglossaries" to process them
%% or the right makeindex options
\newglossary[slg]{symbols}{sym}{sbl}{List of Symbols}
\makeglossaries{}  % intitialize the glossary.  Must be in the main file!!!
%% $URL$
%% $Id$
%% Glossary term definitions
%% Note, you will need to reference them in the text after defining them

\newglossaryentry{glossary}{name=glossary, description={a place to put all the terms that confuse people}}
\newglossaryentry{ohm}{name=ohm,symbol={\ensuremath{\Omega}}, description=unit of electrical resistance}  % Do not break this with newlines!
\newglossaryentry{dissertation}{name=dissertation, description={A large pile of paper that professors use to prop up shelves.}}
\newacronym{rfid}{RFID}{Radio Frequency IDentification}

\newacronym{led}{LED}{light-emitting diodes}

\newacronym{auv}{AUV}{Autonomous Undersea Vehicle}
%% Reference manual: http://ctan.uib.no/macros/latex/contrib/glossaries/glossaries-user.pdf

%% You can also define most entries near where they are used, but this may cause problems.  See manual chapter 4.8
%% Be careful about putting entries on Chapter, Section, etc.  See manual Chapter 6.



%%% Local Variables: 
%%% mode: latex
%%% TeX-master: "msc-tannock-2016"
%%% End: 
  % load glossary definitions in preamble

%\RequirePackage{imakeidx}
%\makeindex[intoc]  % add index to toc

\finalifforcegraphics{hyperref} %hyperlinks even in draft mode
\usepackage[hidelinks]{hyperref} 


\begin{document}
\frontmatter{}
\frontcover{}
\enableindents{}
\starttables{} % setup formatting
%% TOC, list of figures and list of tables are required
\tableofcontents{}%%RUM: "Table of contents"
\thesislistoffigures{}%%RUM: "List of figures"
\thesislistoftables{}%%RUM: "List of tables"
\starttext{}
%\coverchapter{List of drawings and enclosed material}
%RUM: "List of drawings and enclosed material, e.g. CD(as appropriate)"

\listoffixmes{}
% if using fixme package, lists what needs to be done

\chapter{Getting started}
\section{Introduction}
Before using the template in \path{DEGREE-NAME-YEAR.tex}, the author
recommends reading the (in progress) guide to successful projects,
including \LaTeX{} at
\url{http://afs.rnd.ru.is/project/htgaru/trunk/how-to-get-around-projects.pdf}.
Full instructions on setting up \LaTeX{} and SubVersion are there.
Additionally, as issues are discovered, they are added to the
Frequently Asked Questions section in that document.

Text will be double spaced when in draft mode so that there is room
for reviewers to write comments.  Extra blank pages are also removed.
If you want to check formatting in a more final form, switch to
proofing mode.  This text in gray will disappear when you are not in
draft mode.  It is better to use FiXme's as a reminder mechanism,
though.  See Section~\ref{sec:fixme}

\textbf{WARNING} It is very very important that you print the generated PDF at 100\% scaling.   
On many computers, printing is set to adjust the size to fit the page.
Do not do this, because it will mess up the sizes of the font and the placement of the letters.
This becomes very evident when you print it out and the cover window does not line up with your title!

\section{Citing properly}
Any content from external sources must be properly cited.
If any part of this document is plagiarized (from another group,
the internet, or anywhere else) or your references are not properly
cited, you will be in a lot of trouble.

This is in compliance with Reykjavík University's Code of
Ethics\footnote{\url{http://en.ru.is/the-university/ru-code-of-ethics/}
  -- see Item 8 and General Rules on Study and Examinations}
\footnote{\url{http://en.ru.is/studies/study--and-examination-rules/}
  -- see Article 4.5} \footnote{see also the document "Your Work -
  Your Contribution" \url{http://www.ru.is/hugverk}, available only in
  Icelandic}. For your reference, here is a link to guidelines
defining plagiarism and discussing how to avoid it (from the
University of Oxford):

\url{http://www.ox.ac.uk/students/academic/goodpractice/about/}

This template uses BibLaTeX and Biber to produce the Bibliography.
This is because they have better support for accented and Icelandic
characters in UTF-8.  They also do more checking on the bibliography
to look for problems.  Unfortunately, this means that if you have a
semi-broken references.bib, Biber will refuse to produce a
Bibliography.  This is a feature because it is telling you that you
need to fix the problems.  BibTeX will also give errors, but will
continue to generate a Bibliography.  This can cause problems later.
For more information read
\url{http://afs.rnd.ru.is/project/htgaru/trunk/how-to-get-around-projects.pdf}.


\section{Document Structure}
The default document structure is the style called IMRAD:\@ Introduction, Methods, Resarch, And, Discussion.  
This is also known as the American Psychological Association (APA) format and is commonly used in the sciences.\cite{web:uta.fi:imrad}.
The template's current editor believes that this format is a good starting point, but may not be appropriate for some engineering works.  
In those cases, following a more technical report style may be easier: Introduction, Background, Design, Implementation, Analysis,  Conclusion.
For some works, this is still too confining.
In those cases, the editor recommends the Introduction, Body, Conclusion (IBC) format where the body is a series of IMRAD or other structures.
Each section of the body should ask a question and answer it relating to the topic.
If you are writing a thesis that is primarily a collection of papers, then the body section would consist of these papers.

\subsection{Reports for classes}
If you are writing a report for a class using this template, you probably want to get rid of a lot of the extra pages.
The set used in these instructions are a good place to start.
See Listing~\ref{src:minimal-document} for a quick example.
\lstinputlisting[language={[LaTeX]TeX}, caption={Minimal document example},label={src:minimal-document}]{minimal-document.tex}

\chapter{\LaTeX{} techniques}
\section{Fitting things on the page\label{sec:fitting}}
If you see little black boxes on the right in draft mode, this is
\LaTeX{} trying to tell you that it couldn't make it fit nicely.

\section{FixMe\label{sec:fixme} reminders}
\fxnote{This is a note.  It will disappear in final mode, which makes it a safe place to put reminders.}
\fxwarning{This is someting more important that needs to be corrected, but will also disappear in final mode.}
\fxerror{This is something very important to be fixed.  This note will disappear in final mode.}
%\fxfatal{This is a fatal issue that must be fixed before submission (final).  
If any of these still exist except for note when we are not in draft mode, \LaTeX{} will have errors during processing.
If the Fatal error is still there, \LaTeX{} will refuse to produce the document\footnote{This is a {\em feature\/} so that you don't submit the document without fixing really important issues.}.
The Fatal error has been commented out in this document so that we can successfully generate documents but it can be seen in the document's source.

\subsection{Indexes, glossaries, and abbreviations}
This section demonstrates indexes and glossaries.\index{glossary}
Note that you may get errors in this section if you turn the index and glossary options off.
In order to generate the glossary, you will need to run \verb|makeglossaries|, which requires the Perl language to be installed.  On windows, you can get a perl interpreter at \url{http://www.activestate.com/activeperl}
See the \verb|README.TXT| for details.

You might expect to find the term \gls{led} in a \gls{glossary}.
For maximum safety, the glossary and acronyms should be declared using newglossaryentry before \verb|\begin{document}| which is done automatically if you put them into \verb|glossary.tex|.
This is {\em highly recommended\/} in case you reorder your text.
The author did his \gls{dissertation} on \gls{rfid} security.  
You can even have symbols, such as the one for \gls{ohm} which is \glssymbol{ohm}


\subsection{Examples}
Some examples of how to use features of \LaTeX{} follow.

\subsubsection{Equations and references}
\begin{figure}
\centering
\includegraphics[height=20mm]{ru-logo}
\caption{The logo of Reykjavík University}\label{fig:ru-logo}
\end{figure}
Figures and Tables need to be properly formatted and referenced in the
text.  Number figures/tables consecutively, include captions, and
refer to the figures/tables in the text
e.g. Figure~\ref{fig:ru-logo}. Equations need to be numbered
consecutively as well.  Equations need to have each of their variables
defined when they are first used or redefined.  If you need to refer
to particular places in the document, use numbered references.  Note
that the \verb|\ref{BLAH}| and the \verb|\label{BLAH}| must match.
When using \verb|\label{}|, it has to be AFTER the item you are trying
to point to without spaces or inside of its declaration.  It is often
safest to put inside e.g.:
\begin{verbatim}\caption{This is a caption\label{fig:mycaption}}\end{verbatim}

As an example of an equation, in Equation~\ref{eq:freq} is the
relationship between angular frequency and hertz:
\begin{equation}
f = \frac{1}{T} = \frac{1}{2\pi\omega}\label{eq:freq}
\end{equation}
where $f$ is frequency in \si{\hertz}, $T$ is period in \si{\second},
and $\omega$ is angular frequency in \si{\radian\per\second}.
Resistance is in units of \si{\ohm}.  This example uses \verb|\si|
from the siunitx package to format the units correctly.  If you want
to format it with a value, it uses \verb|\SI{value}{unit}| e.g.:
\verb|\SI{2.99E8}{\meter\per\second}| becomes
\SI{2.99E8}{\meter\per\second}

As an example of referring to a particular part of the document, see
Section~\ref{sec:fitting} on page~\pageref{sec:fitting} for information
on how to fit text on a page.

\subsubsection{Chemical Formulas}
Chemical formulas are best formatted by the mhchem package (included
in the custom.tex by default).
\href{http://mirrors.ctan.org/macros/latex/contrib/mhchem/mhchem.pdf}{Click
  here for the documentation}.  This example is disabled by default
because it breaks due to a change in \LaTeX3 that hasn't propogated to
all distributions.


Here is an example in reaction~\ref{react:so4}:

\reaction[react:so4]{SO4^2- + Ba^2+ -> BaSO4 v}


\subsubsection{Tables}
When making tables, there are some very important rules for making
them look professional.  See
\url{http://ctan.uib.no/macros/latex/contrib/booktabs/booktabs.pdf}
For an example of a simple table, see Table~\ref{tab:example}.  Some general
rules are to always avoid \verb|\hrule| and vertical separator lines.

\begin{table}\centering
\caption{Listing of real name to username\label{tab:example}}
\begin{tabular}{lll}\toprule
Firstname &Lastname &username\\\midrule
Joseph &Foley &foley\\\bottomrule
\end{tabular}
\end{table}

\subsection{Mutilingual support}
In the main file, you have already seen the ML macros used to create multi-lingual documents.
You may want to read more about the babel and polyglossia packages to do even more.
Quotation marks can be made smart by using \enquote{encite}.
This will change the quotation marks to follow whatever the current language is set to.
The Context Sensite Quotes (csquote) package is full of useful macros of this sort. 
One caveat: don't use the shorcut "` and '" in the macros nor in the text\footnote{This is the shortcute for German style quotations, which is the same as Icelandic.}.
It doesn't work due to interactions with the multilanguage macros.
If you need to use those quotes, you can switch languages temporarily,  \selectlanguage{icelandic}
"`Like this'" \resetlanguage{}


\subsection{URLs and paths}
URLs are easy to use in both the bibliography and text.  In your
\path{.bib} file, make sure that you only put URLs into the \path{url}
field.  Putting it anywhere else is likely to cause it to go off the
page.  If you are using the \path{hyperref} package, then these will
be clickable.  You can put links in the text with
\verb|\url| like this: \url{http://www.ru.is}.  If you need to
specify a location of a file, use the
\verb|\path| similarly: \path{C:\Documents and Settings\ }.  The only
difference is that paths are not clickable.

\subsection{Hints}
When placing tables and figures, {\em do not\/} force \LaTeX{} to
place them in a given place.  Doing so almost always causes problems
later.  It is better to give hints \verb|[htb]| than to
force it \verb|!h]|.

Add linebreaks at least at the end of every sentence.  This is
especially important if you are using SVN or collaborating with
others.  Everything is done by line number, so if an error or
issue is on line 5 which happens to be a page of text, good luck
finding the problem.

Along those same lines, if you are not storing your thesis in version
control (SVN, git, etc) or on a backed up drive, you are asking for
trouble.
\subsection{Debugging}
When things go wrong, it is sometimes hard to figure out what to do.
The first step is always to go look at the log file.  If you are
making changes and they are not showing up, it is likely that you are
viewing an old file.  Delete the \path{.pdf} and \path{.aux} and try
again.  On some platforms, the package manager can get out of date.
Make sure you have updated/synchronized before running \path{pdflatex} again.

%% ---------------------------------------------------------------
\printbibliography{} %%RUM: "References"

%% If appendices are needed, uncomment the following line
%% and include the appendices in separate files
\appendix{}%%RUM: "Appendicies (as appropriate)
\chapter{Code}\label{cha:code}
% \begin{minted}[linenos, breaklines]{python}
% def boring(args=None):
%     pass
% \end{minted}
\inputminted[linenos, breaklines, frame=lines, stripall]{python}{src/main.py}
% \lstinputlisting[language=Python]{src/main.py}
% You can put code in your document using the listings package, which is
% loaded by default in \path{custom.tex}.  Be aware that the listings
% package does not put code in your document if you are in draft mode
% unless you set the \texttt{forcegraphics} option.
%
% There is an example java (Listing~\ref{src:Data_Bus.java}) and XML
% file (Listing~\ref{src:AndroidManifest.xml}).  Thanks to the
% \texttt{url} package, you can typeset OSX and unix paths like this:
% \path{/afs/rnd.ru.is/project/thesis-template}.  Windows paths:
% \path{C:\windows\temp\ }.  You can also typeset them using the menukey
% package, but it tends to delete the last separator and has other
% complications.\footnote{The menukey package has issues with biblatex,
%   read \path{custom.tex} for more information.}
%
% If you are trying to include multiple different languages, you should
% go read the documentation and set these up in \path{custom.tex}.  You
% will save yourself a lot of effort, especially if you have to fix
% anything.
%
% %I have put the source code in the \directory{src/} folder.
% \lstinputlisting[language=Java, firstline=1,
% lastline=40, caption={Data\_Bus.java: Setting up the class.},
% label={src:Data_Bus.java}]{src/Data_Bus.java}
%
% \lstinputlisting[language={[android]XML}, firstline=1, lastline=20,
% caption={AndroidManifest.xml: Configuration for the Android UI.},
% label={src:AndroidManifest.xml}]{src/AndroidManifest.xml}

%%% Local Variables:
%%% mode: latex
%%% TeX-master: "msc-tannock-2016"
%%% End:
 % as an example, perhaps some of your code

%\backmatter{} % Sections after this don't get numbers
%% We prefer that all elements be numbered

%%%%%%%%%%%%% SHOW GLOSSARY %%%%%%%%%%%%%%%%%
%% Glossary, optional.  A good idea on longer documents
%% Remember to run "makeglossaries <DEGREE-NAME-YEAR>"
%%   to update the glossary
\printglossary{}%%RUM: Not mentioned
\printacronyms{}%%RUM: Not mentioned

%% Glossary debugging code
%% Uncomment this to figure out of one of the glossary entries is broken
%% by putting all of the entries in the glossary.
%\ifthenelse{\boolean{debug}}{%
%  \glsaddall{} %put all entries in
%  \printglossaries{}}{}% dump them all
%% This command dumps all of the glossaries

%%%%%%%%%%%%% SHOW INDEX %%%%%%%%%%%%%%%%%%
%% Index, optional.  A good idea on longer documents

% You can put instructions at the beginning of the index:
%\renewcommand{\preindexhook}{%
%  The first page number is usually, but not always,
%  the primary reference to the indexed topic.\vskip\onelineskip}

%% You may have to run "makeindex" to have it be generated
%% Depending upon which package you chose.
%% 
\clearforchapter{}
\printindex{}%%RUM: Not mentioned

\backcover{}%%RUM: "Back cover (standard)
\end{document}

%%% Local Variables: 
%%% mode: latex
%%% TeX-master: t
%%% End: 
}{}

\section{Introduction}
\label{sec:Introduction}
\subsection{Permutations}
\label{sub:Permutations}

\begin{frame}{Permutations}
  A \emph{permutation} is a \emph{bijection}, \(\pi\), from the set
  \(\nrange{n} = \setrange{1}{n}\) to itself.

  More intuitively \textquote[Knuth~\cite{Knuth}]{A \emph{permutation of \(n\) objects} is
  an arrangement of \(n\) distinct objects in a row}.

  We write permutations in \emph{one-line notation}, writing
  the entries of the permutation in order
  \begin{equation*}
    \pi = \perm{\pi(1),\pi(2),\dotso,\pi(n)}
  \end{equation*}
  \begin{example} The \(6\) permutations on \(\nrange{3}\) are
    \begin{equation*}
      123, 132, 213, 231, 312, 321
    \end{equation*}
  \end{example}
\end{frame}

\begin{frame}
  We can display a permutation in a \emph{plot} to give a graphical
  represention. We plot the points \((i,\pi(i))\) in a Cartesian coordinate
  system.
  \begin{figure}[htb]
    \centering
    \pattern{}{2,3,1}{}
    \caption{Plot of the permutation \(\perm{2,3,1}\)}
  \end{figure}
  In this setting we call the elements of the permutations \emph{points}.

  The set of all permutations of length \(n\) is \(\mathfrak{S}_n\) and
  has size \(n!\). The set of all permutations is
  \(\mathfrak{S}=\bigcup_{i=0}^{\infty}\mathfrak{S}_i\).
\end{frame}
\subsection{Classical Permutation Patterns}
\label{sub:Classical Permutation Patterns}

\begin{frame}{Classical Permutation Patterns}
  \emph{Classical permutation patterns} capture many interesting combinatorial
  objects and properties.

  \begin{definition}[Order Isomorphism]
    Two substrings \(\alpha_1\alpha_2\dotsm\alpha_n\) and
    \(\beta_1\beta_2\dotsm\beta_n\) are said to be \emph{order isomorphic}
    if \(\alpha_r<\alpha_s\) if and only if \(\beta_r<\beta_s\).
  \end{definition}
\end{frame}

\begin{frame}
  \begin{definition}
    A permutation \(\pi\) is said to \emph{contain} the \emph{classical
    permutation pattern} \(\sigma\)
    (denoted \(\sigma \preceq \pi\)) if there is some
    subsequence \(i_1i_2\dotsm{}i_k\) such that the sequence
    \(\pi(i_1)\pi(i_2)\dotsm\pi(i_k)\) is order isomorphic to
    \(\sigma(1)\sigma(2)\dotsm\sigma(k)\).
  \end{definition}

  If \(\pi\) does not contain \(\sigma\) we say that \(\pi\) \emph{avoids}
  \(\sigma\).

  The set of permutations of length \(n\) avoiding a pattern \(\sigma\) is
  denoted as \(\Av_n(\sigma)\) and
  \begin{equation*}
    \av{\sigma} = \bigcup_{i=0}^{\infty}\Av_i(\sigma)
  \end{equation*}
\end{frame}

\begin{frame}
  \begin{example}
    The permutation \(\pi = \perm{2,4,1,5,3}\) contains the pattern
    \(\sigma = \perm{2,3,1}\)
    \begin{figure}[htb]
      \centering
      \pattern{2,4,5}{2,4,1,5,3}{}
      \caption{Plot of the permutation \(\perm{2,4,1,5,3}\) with an occurrence
      of \(\perm{2,3,1}\) indicated}
    \end{figure}
  \end{example}
\end{frame}

\subsection{Mesh Patterns}
\label{sub:Mesh Patterns}
\begin{frame}{Mesh Patterns}
  \emph{Mesh patterns} are a natural extension of classical permutation patterns.
  \begin{definition}
    A \emph{mesh pattern} is a pair
    \begin{equation*}
      p = (\tau,R)\text{ with } \tau \in \mathfrak{S}_k \text{ and } R \subseteq
      [0,k]\times [0,k].
    \end{equation*}
  \end{definition}
  We say that \(\tau\) is the \emph{underlying classical pattern} of \(p\).
\end{frame}

\begin{frame}
  \begin{example}
      The pattern \(p=\mperm{2,1,3}{\{(0,1),(0,2),(0,3),(1,0),(1,1),(2,1),(2,2)\}}=
      \textpattern{}{ 2, 1, 3 }{ 0/1, 0/2,0/3, 1/0, 1/1, 2/1, 2/2 }\) is contained in
      \(\pi = \perm{3,4,2,1,5}\).
      \begin{figure}[htb]
        \centering
      \pattern{1,3,5}{3,4,2,1,5}{0/2,0/3,0/4,0/5,
                                 1/0,1/1,1/2,
                                 2/0,2/1,2/2,
                                 3/2,3/3,3/4,
                                 4/2,4/3,4/4}
      \caption{ An occurrence of \(p\) in \(\pi\) }
      \end{figure}
  \end{example}
\end{frame}

\begin{frame}
  \begin{example}
      The pattern \(q=\mperm{2,1}{\{(0,1),(0,2),(1,0),(1,1)\}}=
      \textpattern{}{ 2, 1 }{ 0/1, 0/2, 1/0, 1/1}\) is contained in
      \(p=\mperm{2,1,3}{\{(0,1),(0,2),(0,3),(1,0),(1,1),(2,1),(2,2)\}}=
      \textpattern{}{ 2, 1, 3 }{ 0/1, 0/2,0/3, 1/0, 1/1, 2/1, 2/2 }\) as a
      subpattern.
      \begin{figure}[htb]
        \centering
      \pattern{1,2}{ 2, 1, 3 }{ 0/1, 0/2,0/3, 1/0, 1/1, 2/1, 2/2 }
      \caption{ An occurrence of \(q\) in \(p\) }
      \end{figure}
  \end{example}
\end{frame}
%%RUM: Introduction
\chapter[Coincidence Classification]{Coincidences amongst families of mesh patterns and classical patterns}\label{chap:coincs}
One interesting question to ask about permutation patterns considers when a
pattern may be avoided by, or contained in, arbitrary permutations. Two
patterns \(\pi\) and \(\sigma\) are said to be \emph{coincident} if the set
of permutations that avoid \(\pi\) is the same as the set of permutations that
avoid \(\sigma\), \ie \(\av{\pi} = \av{\sigma}\). This extends to sets of patterns as well as single patterns.

We consider the avoidance sets, \(\av{\{\pi,q\}}\) where \(\pi\) is a classical pattern
of length \(3\) and \(q\) is a mesh pattern of length \(2\) in order to
establish some rules about when these two sets give the same avoidance set.
We fix \(\pi\) in order to define the  coincidence and say that \(\pi\) is the
\emph{dominating pattern}.

It is useful to be able to modify a mesh pattern by adding points to an
already existing mesh pattern. First adding a single point into a pattern.
\begin{definition}
Given a mesh pattern \(p = (\pi,R)\) \(\addp[D]{p}{(a,b)}\) gives a mesh
pattern \(p^\prime = (\pi^\prime,R^\prime)\) with length \(\setsize{\pi}+1\) defined by
\begin{equation*}
    \pi^\prime(i) = \begin{cases}
    \pi(i) & \text{if } i\neq a+1 \text{ and } \pi(i)<b\\
    \pi(i)+1 & \text{if } i\neq a+1 \text{ and } \pi(i)>b\\
    b+1 & \text{if } i = a+1
\end{cases}
\end{equation*}
and
\begin{equation*}
        R^\prime = \bigcup_{(i,j)\in R}r((i,j)),
\end{equation*}
where \(r((i,j))\) is defined by
\begin{equation*}
    r((i,j)) =
    \begin{cases}
        \{(i,j)\} & \text{if } i<a, j<b\\
        \{(i,j),(i,j+1)\} & \text{if } i<a, j=b\\
        \{(i,j+1)\} & \text{if } i<a, j>b\\
        \{(i,j),(i+1,j)\} & \text{if } i=a, j<b\\
        \{(i,j+1),(i+1,j+1)\} & \text{if } i=a, j>b\\
        \{(i+1,j)\} & \text{if } i>a, j<b\\
        \{(i+1,j),(i+1,j+1)\} & \text{if } i>a, j=b\\
        \{(i+1,j+1)\} &\text{if } i>a, j>b\\
    \end{cases}
\end{equation*}

If the shading set \(D\) is non-empty we can modify definition of the directions
slightly
\begin{equation*}
    \begin{aligned}
        N &= \{(a,b+1),(a+1,b+1)\} \\
        E &= \{(a+1,b),(a+1,b+1)\} \\
        S &= \{(a,b),(a+1,b)\} \\
        W &= \{(a,b),(a,b+1)\}
    \end{aligned}
\end{equation*}
And we add the union of the sets in \(D\) into the mesh set \(R^\prime\).
\end{definition}

Given a mesh pattern \(p\), \(\addp[D]{p}{(a,b)}\) is the operation that returns
a mesh pattern equivalent to placing a point in the center of box \((a,b)\),
where \((a,b)\) is not in the mesh set of \(p\) , with shading defined by
\(D\subseteq\{N,E,S,W\}\).
The set \(D\) defines the shading by indicating that the boxes in the cardinal
directions in \(D\) next to the point are shaded in the resulting pattern.
Since there is no ambiguity we let \(\addpe[D]\) be equivalent to \(\addp[D]{\varepsilon}{(0,0)}\).

\begin{example}
    The result of adding a single point to the empty permutation for each cardinal direction.
 \begin{align*}
     \addpe[\{N\}] = \textpattern{}{1}{0/1,1/1} &\qquad\addpe[\{E\}] =\textpattern{}{1}{1/1,1/0} \qquad \\
     \addpe[\{S\}] = \textpattern{}{1}{0/0,1/0} &\qquad\addpe[\{W\}] =\textpattern{}{1}{0/0,0/1} \qquad
\end{align*}
    A more complex example for \texttt{add\_point} could be
\begin{equation*}
\addp[\{N,E\}]{\pattern{}{2,1,3}{0/3}}{(2,3)} = \pattern{}{2,1,4,3}{0/3,0/4,2/4,3/4,3/3}
\end{equation*}
\end{example}

It is also useful to think about adding an ascent, or descent, into a pattern

\begin{definition}
Considering
only adding the ascent, as adding a descent is very similar.
Given a mesh pattern \(p = (\pi, R)\) \(\adda[D]{p}{(a,b)}\) gives a mesh
pattern \(p^\prime = (\pi^\prime,R^\prime)\) with length \(\setsize{\pi}+2\) defined by
\begin{equation*}
    \pi^\prime(i) = \begin{cases}
    \pi(i) & \text{if } i\neq a+1,a+2 \text{ and } \pi(i)<b\\
    \pi(i)+2 & \text{if } i\neq a+1,a+2 \text{ and } \pi(i)>b\\
    b+j, b+j & \text{if } i = a+j, j\in\{1,2\}
\end{cases}
\end{equation*}
and
\begin{equation*}
        R^\prime = \{(a+1,b),(a,b),(a+1,b),(a+2,b),(a+1,b+2)\} \cup \bigcup_{(i,j)\in R}r((i,j)),
\end{equation*}
where \(r((i,j))\) is defined by
\begin{equation*}
    r((i,j)) =
    \begin{cases}
        \{(i,j)\} & \text{if } i<a, j<b\\
        \{(i,j),(i,j+1),(i,j+2)\} & \text{if } i<a, j=b\\
        \{(i,j+2)\} & \text{if } i<a, j>b\\
        \{(i,j),(i+1,j),(i+2,j)\} & \text{if } i=a, j<b\\
        \{(i,j+2),(i+1,j+2),(i+2,j+2)\} & \text{if } i=a, j>b\\
        \{(i+2,j)\} & \text{if } i>a, j<b\\
        \{(i+2,j),(i+2,j+1),(i+2,j+2)\} & \text{if } i>a, j=b\\
        \{(i+2,j+2)\} &\text{if } i>a, j>b\\
    \end{cases}
\end{equation*}
\end{definition}

Given a pattern \(p\), \(\addd{p}{(a,b)}\), and \\ \(\adda{p}{(a,b)}\),
are the operations that return a mesh pattern equivalent to placing a
decrease, or increase, in the center of box \((a,b)\), where \((a,b)\) is not
in the mesh set of \(p\), in \(p\).

\begin{example}
    \begin{align*}
        \addae = \pattern{}{1,2}{0/1,1/0,1/1,1/2,2/1}\\
        \addde = \pattern{}{2,1}{0/1,1/0,1/1,1/2,2/1}
    \end{align*}

    \noindent A more complex example is
\begin{equation*}
\adda{\pattern{}{2,3,1}{0/0,1/0,1/3,
                        3/0,3/1,3/3}}{(1,1)}
= \pattern{}{4,2,3,5,1}{0/0,
                        1/0,1/2,1/5,
                        2/0,2/1,2/2,2/3,2/5,
                        3/0,3/2,3/5,
                        5/0,5/1,5/2,5/3,5/5}
\end{equation*}
\end{example}

We now attempt to fully classify coincidences in families characterised by avoidance
of a classical pattern of length \(3\) and a mesh pattern of length \(2\), that is
finding and explaining all coincidences where \(\av{\{p,m\}} = \av{\{p,m^\prime\}}\).

It can be easily seen that in order to classify coincidences one need only
consider coincidences within the family of mesh patterns with the same underlying
classical pattern, this is due to the fact that \(\perm{2,1}\in\av{\mperm{1,2}{R}}\)
and \(\perm{1,2}\in\av{\mperm{2,1}{R}}\) for all mesh-sets \(R\).

We know that there are a total of \(512\) mesh-sets for each underlying classical
pattern. By use of the previous results of \textcite{DBLP:journals/corr/ClaessonTU14}\footnote{
The authors use the Simultaneous Shading Lemma, a closure result and one worked out special case.
} the number of coincidence classes can be reduced to \(220\).
% the Shading
% Lemma\cite[Lemma~3.11]{DBLP:journals/combinatorics/HilmarssonJSVU15},
% Simultaneous Shading Lemma\textcite[Lemma~7.6]{DBLP:journals/corr/ClaessonTU14}, and
% one special case, we can reduce the number of equivalence classes to \(220\).

By discussion of a number of rules we will show that the number of coincidence
classes follows the values shown in  \TableRef{tab:domclasses}. The experimental
data in the last row of the table is calculated on permutations up to length 11.

\begin{table}[htb]
\begin{center}
\begin{tabularx}{\textwidth}{c|Y|Y|Y|Y|}
\cline{2-5}
& \multicolumn{4}{|c|}{Dominating Pattern}\\
\cline{2-5}
& \multicolumn{2}{|c|}{\(\perm{2,3,1}\)}&\multicolumn{2}{|c|}{\(\perm{3,2,1}\)}\\
\cline{2-5}
& \(\perm{1,2}\)&\(\perm{2,1}\)&\(\perm{1,2}\)&\(\perm{2,1}\)\\
\hline
\multicolumn{1}{|r|}{No Dominating rule}&\(220\)&\(220\)&\(220\)&\(220\)\\
\hline
\multicolumn{1}{|r|}{First Dominating rule}&\(85\)&\(43\)&\(220\)&\(29\)\\
\hline
\multicolumn{1}{|r|}{Second Dominating rule}&\(59\)&\(39\)&\(220\)&\(29\)\\
\hline
\multicolumn{1}{|r|}{Third Dominating rule}&\(56\)&\(39\)&\(220\)&\(29\)\\
\hline
\multicolumn{1}{|r|}{Experimental class size}&\(56\)&\(39\)&\(213\)&\(29\)\\
\hline
\end{tabularx}
\end{center}
    \caption{Coincidence class number reduction by application of Dominating Rules}
    \label{tab:domclasses}
\end{table}

From the table it can be seen that the rules established capture almost all
coincidences. However, there are still some coincidences that are not able to
be explained by the rules. This shows that complete coincidence classification
of mesh patterns is a very difficult task, even when we have additional tools
available.
\section{Coincidence classes of Av(\{321, (21, \textit{R})\}).}
Through experimentation, up to permutations of length 11, we discover that there
are at most 29 coincidence classes of mesh patterns with underlying classical
pattern \(\perm{2,1}\).

\begin{proposition}[First Dominating Pattern Rule]
    \label{prop:dom1}
    Given two mesh patterns \(m_1 =(\sigma, R_1)\) and \(m_2 = (\sigma, R_2)\),
    and a dominating classical pattern \(\pi = (\pi,\emptyset)\) such that
    \(\setsize{\pi} \le \setsize{\sigma} + 1\), the sets \(\av{\{\pi,m_1\}}\) and
    \(\av{\{\pi,m_2\}}\) are coincident if

    \begin{enumerate}
        \item \(R_1 \triangle R_2 = \{(a,b)\}\)
        \item \(\pi \preceq \addp{\sigma}{(a,b)}\)\label{prop:dom1:cont}
    \end{enumerate}
\end{proposition}
In order to prove this proposition we must first make the following note.

\begin{note}
    \label{not:downcmesh}
    Let \(R^\prime \subseteq R\). Then any occurence of \((\tau, R)\) in a permutation
    is an occurence of \((\tau,R^\prime)\).
\end{note}
% \begin{proof}
%     Assume without meaningful loss of generality that \(R^\prime\) is a proper
%     subset of \(R\).
%
%     Consider an occurence of \((\tau, R)\) in a permutation \(\sigma\), obviously
%     this corresponds to an occurrence of \(\tau\) in \(\sigma\). Now consider
%     the mesh sets \(R\) and \(R^\prime\), since \(R^\prime \subseteq R\) then
%     there are more restrictions on where points are in an occurrence of \((\tau,R)\).
%     Namely, for every shaded box in \(R\) the corresponding region in \(\sigma\)
%     must contain no points, since \(R^\prime\) has less shading than \(R\) there
%     exists a region in the occurence of \((\tau,R)\) in \(\sigma\) that is
%     now devoid of restrictions. However, by removing restrictions we cannot
%     make an occurence become not an occurence, and therefore the same occurrence
%     of \(\tau \) in \(\sigma\) is now an occurence of \((\tau, R^\prime)\).
% \end{proof}
\begin{proof}[Proof of \PropositionRef{prop:dom1}]
    We need to prove that \(\av{\{\pi,m_1\}} = \av{\{\pi,m_2\}}\).

    \noindent Assume without meaningful loss of generality that \(R_2 = R_1 \cup \{(a,b)\}\).
    % Consider a permutation \(\omega\) that contains an occurrence of \(m_2\)
    % by \NoteRef{not:downcmesh} any of these occurrences is also an occurence
    % of \(m_1\). This proves that every occurence of \(m_2\) is also an
    % occurence of \(m_1\) and therefore \(\av{\{\pi,m_1\}} \subseteq \av{\{\pi,m_2\}}\).
    Since \(R_1\) is a subset of \(R_2\), \NoteRef{not:downcmesh} states that
    \(\av{\{\pi,m_1\}} \subseteq \av{\{\pi,m_2\}}\)

    Now we consider a permutation \(\omega^\prime\in \av{\pi}\),
    containing an occurrence of \(m_1\). Consider placing a point in the
    region corresponding to the box \((a,b)\), regardless of where in this
    region we place the point by condition~\ref{prop:dom1:cont} of the
    Proposition we create an occurrence of \(\pi\), therefore there can be no
    points in this region, which could have been represented in the mesh set
    \(R_1\) by adding the box \((a,b)\). Hence every occurrence of
    \(m_1\) is in fact an occurrence of \(m_2\), and we have that
    \(\av{\{\pi,m_2\}} \subseteq \av{\{\pi,m_1\}}\).

    Taking both directions of the containment we can therefore draw the
    conclusion that \(\av{\{\pi,m_1\}} = \av{\{\pi,m_2\}}\).
\end{proof}

All coincidence classes of \(\av{\{\perm{3,2,1}, \mperm{2,1}{R}\}}\) can be
explained by application of \PropositionRef{prop:dom1}. % repeated?

This rule can be understood very in graphical form. In the
pattern in \FigureRef{fig:rule1} we can gain shading in three boxes since if
there is a point in any of these boxes we would get an occurrence of the
dominating pattern \(\perm{3,2,1}\).

\begin{figure}[htb]
    \begin{center}
    \begin{tikzpicture}[scale=\picscale]
        \modpattern{}{2,1}{1/0,1/1,1/2,2/2}
    \end{tikzpicture}
        \raisebox{2ex}{\(\mapsto\)}
    \begin{tikzpicture}[scale=\picscale]
        \filldraw (2.5,0.5) circle (3pt);
        \modpattern{}{2,1}{1/0,1/1,1/2,2/2,2/0}
    \end{tikzpicture}

        \caption{Visual depiction of first dominating pattern rule.}
        \label{fig:rule1}
    \end{center}
\end{figure}

Note that there are two natural extensions of this rule. We can replace \(\pi\)
with a set of classical patterns, or we can consider \(\pi\) to be a mesh pattern.

\section{Equivalence classes of Av(\{231, (21, \textit{R})\}).}
By application of \PropositionRef{prop:dom1} we obtain \(43\) coincidence
classes. Experimentation shows that there are in fact at most \(39\) coincidence
classes, for example the following two patterns are coincident in
\(\av{\perm{2,3,1}}\) but this is not explained by \PropositionRef{prop:dom1}.
\begin{equation*}
    m_1 = \pattern{}{2,1}{0/0,0/2,1/0,2/1} \text{ and } m_2 = \pattern{}{2,1}{0/0,0/2,1/0,1/1,2/1}
\end{equation*}

Consider an occurrence of \(m_1\) in a permutation in \(\av{\perm{2,3,1}}\), consisting
of elements \(x\) and \(y\). If the region corresponding
to the box \((1,1)\) is empty we have an occurrence of \(m_2\). Otherwise, if there is any
increase in this box then we would have an occurrence of \(\perm{2,3,1}\),
however, since we are in \(\av{\perm{2,3,1}}\) this is not possible. This box
must therefore contain a (non-empty) decreasing subsequence.
 This gives rise to the following lemma:

\begin{lemma}
    \label{lem:incdec}
    Given a mesh pattern \(m =(\sigma, R)\), where the box \((a,b)\) is not
    in \(R\), and a dominating classical pattern \(\pi = (\pi,\emptyset)\) if
    \(\pi \preceq \adda{\sigma}{(a,b)}\)\\(\(\pi \preceq \addd{\sigma}{(a,b)}\))
    then in any occurrence of \(m\) in a permutation \(\varrho\) the region
    corresponding to the box \((a,b)\) can only contain an increasing
    (decreasing) subsequence of \(\varrho\).
\end{lemma}
The proof is analogous to the proof of \PropositionRef{prop:dom1}.

Going back to our example mesh patterns
\begin{equation*}
    \begin{tikzpicture}[scale=\picscale]
        \modpattern{}{2,1}{0/0,0/2,1/0,2/1}
        \filldraw (1.66,1.33) circle (3pt);
        \filldraw (1.33,1.66) circle (3pt);
    \end{tikzpicture}
\end{equation*}
We know that the region corresponding to the box \((1,1)\) contains a decreasing subsequence.
If we let \(z\) be the topmost point in this decereasing subsequence, then \(xz\)
is an occurrence of \(m_2\). This shows that our two example patterns are coincident.

This result generalises into the following rule for categorising coincidences
of mesh patterns in cases where there is a dominating classical pattern.

\begin{proposition}[Second Dominating Pattern Rule]
    \label{prop:dom2}
    Given two mesh patterns \(m_1 =(\sigma, R_1)\) and \(m_2 = (\sigma, R_2)\),
    and a dominating classical pattern \(\pi = (\pi,\emptyset)\) such that
    \(\setsize{\pi} \le \setsize{\sigma} + 2\), the sets \(\av{\{\pi,m_1\}}\) and
    \(\av{\{\pi,m_2\}}\) are coincident if

    \begin{enumerate}
        \item \(R_1 \triangle R_2 = \{(a,b)\}\)
        \item   \begin{enumerate}
                \item\label{prop:dom2:condc} \(\pi \preceq \adda{\sigma}{(a,b)}\) and
                        \begin{enumerate}
                            \item \((a+1,b) \in \sigma\) and \((a+1,b-1)\notin R\) and \\
                                \((x,b-1)\in R \implies (x,b) \in R \) (where \(x\neq a,a+1\)) and\\
                                  \((a+1,y)\in R \implies (a,y) \in R\) (where \(y\neq b-1,b\)).
                            \item \((a,b+1) \in \sigma\) and \((a-1,b+1)\notin R\) and \\
                                  \((x,b+1)\in R \implies (x,b) \in R\) (where \(x\neq a-1,a\)) and\\
                                  \((a-1,y)\in R \implies (a,y) \in R\) (where \(y\neq b,b+1\)).
                        \end{enumerate}
                    \item \(\pi \preceq \addd{\sigma}{(a,b)}\) and
                        \begin{enumerate}
                            \item \((a+1,b+1) \in \sigma\) and \((a+1,b+1)\notin R\) and \\
                                  \((x,b+1)\in R \implies (x,b) \in R\) (where \(x\neq a,a+1\)) and\\
                                  \((a+1,y)\in R \implies (a,y) \in R\) (where \(y\neq b,b+1\)).
                            \item \((a,b) \in \sigma\) and \((a-1,b-1)\notin R\) and \\
                                  \((x,b+1)\in R \implies (x,b) \in R\) (where \(x\neq a-1,a\)) and\\
                                  \((a-1,y)\in R \implies (a,y) \in R\)  (where \(y\neq b-1,b\)).
                        \end{enumerate}
                \end{enumerate}
    \end{enumerate}
\end{proposition}
\begin{proof}
    We need to prove that \(\av{\{\pi,m_1\}} = \av{\{\pi,m_2\}}\).

    \noindent Assume without meaningful loss of generality that \(R_2 = R_1 \cup \{(a,b)\}\).

    Consider a permutation \(\omega\) that contains an occurrence of \(m_2\).
    By \NoteRef{not:downcmesh} any of these occurrences is also an occurence
    of \(m_1\). This proves that every occurence of \(m_2\) is also an
    occurence of \(m_1\) and therefore \(\av{\{\pi,m_1\}} \subseteq \av{\{\pi,m_2\}}\).

    We will consider taking the first branch of every choice.
    Now consider a permutation \(\omega^\prime \in \av{\pi}\). Suppose
    \(\omega^\prime\) contains \(m_1\) and consider the region corresponding to
    \((a,b)\) in \(R_1\).

    If the region is empty, the occurrence of \(m_1\) is trivially an occurence
    of \(m_2\).

    Now consider if the region is non-empty, by \LemmaRef{lem:incdec} and
    condition~\ref{prop:dom2:condc} of the proposition this region must contain
    a decreasing subsequence.
    We can choose the topmost point in the region to replace the corresponding
    point in the mesh pattern and the points from the subsequence are now in
    the box southeast of the point. The other conditions allow this to be done
    without points being present in regions that were shaded. Hence there are no
    points in the region corresponding to the box \((a,b)\) in the mesh
    pattern, and therefore we can shade this region. This implies that every
    occurrence of \(m_1\) in \(\av{\pi}\) is in fact an occurrence of \(m_2\)
    so \(\av{\{\pi,m_2\}} \subseteq \av{\{\pi,m_1\}}\).

    Similar arguments satisfy the remainder of the branches.
\end{proof}
\begin{figure}[ht]
\centering
\begin{tikzpicture}[scale = 0.75]

     % Make the grid
     \draw[step=1cm,gray!50] (2/3,2/3) grid (5+1/3,5+1/3);

     % Make black cross in the middle
     \draw[-,very thick] (3,2/3) -- (3,5+1/3);
     \draw[-,very thick] (2/3,3) -- (5+1/3,3);
     \draw[-,very thick,gray] (4,2) -- (2,2) -- (2,4) -- (4,4) -- (4,2);

     \fill[black] (3,3) circle (0.1);
     \fill[red] (2,3) circle (0.1);

     % Draw the dots on the side
     \foreach \x in {1,2,4,5}{
         \draw[dotted, gray!50] (\x, 0) -- (\x, 2/3);
         \draw[dotted, gray!50] (\x, 5+1/3) -- (\x, 6);
         \draw[dotted, gray!50] (0,\x) -- (2/3,\x);
         \draw[dotted, gray!50] (5+1/3,\x) -- (6,\x);
        }
     % Draw black dots next to the cross
     \draw[dotted, black, very thick] (3, 0) -- (3, 2/3);
     \draw[dotted, black, very thick] (3, 5+1/3) -- (3, 6);
     \draw[dotted, black, very thick] (0, 3) -- (2/3, 3);
     \draw[dotted, black, very thick] (5+1/3, 3) -- (6, 3);

     \draw[->-=.5,thick] (3,3) to[bend right=15](2,4);

     \draw[->, thick] (5,3) -- (5,4);
     \draw[->, thick] (1,3) -- (1,4);
     \draw[->, thick] (3,5) -- (2,5);
     \draw[->, thick] (3,1) -- (2,1);

     \node [black,right] at (6,3) {$b$};
     \node [black,below] at (2,0) {$a$};

\end{tikzpicture}
    \caption{If the conditions of \PropositionRef{prop:dom2} are satisfied the box
    \((a-1,b)\) can be shaded.}
\end{figure}

This proposition essentially states that we slide all of the points in the box we desire
to shade diagonally, and chose the topmost/bottommost point to replace the original
point in the mesh pattern.

By taking the First Dominating Pattern Rule and the Second Dominating Pattern
Rule together coincidences of classes of the form \(\av{\{\perm{2,3,1},\mperm{2,1}{R}\}}\)
are fully explained, obtaining 39 coincidence classes of mesh patterns.

\section{Equivalence classes of Av(\{231, (12, \textit{R})\}).}
When considering the coincidence classes of \(\av{\perm{2,3,1},\mperm{1,2}{R}}\)
we first apply the two Dominating Pattern rules previously established.
Starting from \(220\) classes, application of the first Dominating Pattern rule
gives \(85\) classes. Following this with the second Dominating Pattern rule
reduces the number of classes to \(59\). However we know that there are patterns
where the coincidences are not explained by the rules given above.

For example the patterns
\begin{equation*}
    m_1 = \pattern{}{1,2}{0/2,2/0,2/1} \text{ and } m_2 = \pattern{}{1,2}{0/2,1/0,2/0,2/1}
\end{equation*}
are experimentally coincident. This coincidence is not explained by our rules,
it is necessary to attempt to capture these coincidences by establishing more
rules.

Consider an occurrence of \(m_1\) in a permutation, if the region corresponding
to the box \((1,0)\) is empty then we have an occurrence of \(m_2\). Now look at
the case when this region is not empty. Consider choosing the rightmost point in
region.
% In order to rigorously establish this coincidence we need to consider what would
% happen if we were to choose a point in the box \((1,0)\) that we would like to shade.
% In order to get a chance of getting the shading we want consider choosing the
% rightmost point in the box.
This gives us an occurrence of the following mesh pattern.
\begin{equation*}
    \pattern{}{2,1,3}{0/3,2/0,2/1,3/0,3/1,3/2}
\end{equation*}
By application of \PropositionRef{prop:dom1} we then achieve the following
mesh pattern
\begin{equation*}
    \pattern{2,3}{2,1,3}{0/3,1/3,2/0,2/1,3/0,3/1,3/2}
\end{equation*}
If we look at the highlighted points we see that the subpattern is an occurrence
of the mesh pattern that we originally desired. This gives rise to the following
rule:

\begin{proposition}[Third Dominating Pattern Rule]
    \label{prop:dom3}
    Given two mesh patterns \(m_1 =(\sigma, R_1)\) and \(m_2 = (\sigma, R_2)\),
    and a dominating classical pattern \(\pi = (\pi,\emptyset)\), the sets
    \(\av{\{\pi,m_1\}}\) and \(\av{\{\pi,m_2\}}\) are coincident if
    \begin{enumerate}
        \item \(R_1 \triangle R_2 = \{(a,b)\}\)
        \item\label{prop:dom3:condocc} \(\addp[D]{\mperm{\sigma}{R_1}}{(a,b)}\)
            where \(D\in\{N,E,S,W\}\)
            is coincident with a mesh pattern containing an occurrence of
            \(\mperm{\sigma}{R_2}\).
    \end{enumerate}
\end{proposition}
\begin{proof}
    We need to prove that \(\av{\{\pi,m_1\}} = \av{\{\pi,m_2\}}\).

    \noindent Assume without meaningful loss of generality that \(R_2 = R_1 \cup \{(a,b)\}\).

    Consider a permutation \(\omega\) that contains an occurrence of \(m_2\).
    By \NoteRef{not:downcmesh} any of these occurrences is also an occurence
    of \(m_1\). This proves that every occurence of \(m_2\) is also an
    occurence of \(m_1\) and therefore \(\av{\{\pi,m_1\}} \subseteq \av{\{\pi,m_2\}}\).

    Now consider a permutation \(\varrho\) in \(\av{\pi}\) that contains an
    occurrence of \(m_1\). If the region corresponding to the box \((a,b)\)
    is empty then we have an occurrence of \(m_2\). If the region is non-empty
    then by condition~\ref{prop:dom3:condocc}  of the proposition there exists
    a direction such that there exists an occurrence of a mesh pattern of
    length one longer than \(m_1\) in this position. This mesh pattern is
    coincident with another mesh pattern that contains an
    occurrence of \(m_2\). Hence, every occurrence of \(m_1\) leads to an occurrence
    of \(m_2\).  Thus \(\av{\{\pi,m_2\}} \subseteq \av{\{\pi,m_1\}}\) and the
    two patterns are coincident.
\end{proof}

By application of this rule we can reduce the number of classes in
\(\av{\{\perm{2,3,1},(\perm{1,2}, R)\}}\) to \(56\).

\section{Equivalence classes of Av(\{321, (12, \textit{R})\}).}
When considering coincidences of mesh patterns with underlying classical
pattern \(\perm{1,2}\) in \(\av{\perm{3,2,1}}\) application of the previously
established rules give no coincidences. Through experimentation we discover
that there are \(7\) non-trivial coincidence classes (all others are singletons)
which can be explained through the
use of two different lines of reasoning. Since the number of coincidences is
so small we will reason for these coincidences without attempting to generalise
into concrete rules.

Intuitively it is easy to see why our previous rules have no
power here. There is nowhere that it is possible to add a single point to
gain an occurrence of \(\pi = \perm{3,2,1}\). It is also impossible to have
a position where addition of an increase, or decrease, provides extra
shading power.

The patterns
\begin{equation*}
    m_1 = \pattern{}{1,2}{0/0, 0/2, 1/1, 1/2, 2/1} \text { and } m_2 \pattern{}{1,2}{0/0, 0/1, 0/2, 1/1, 1/2, 2/1}
\end{equation*}
are equivalent in \(\av{\perm{3,2,1}}\). (There are 3 symmetries of these
patterns that are also equivalent to each other by the same reasoning.)

Consider the region corresponding to the box \((0,1)\) in any occurrence of \(m_1\),
in a permutation. By \LemmaRef{lem:incdec} it must contain an increasing subsequence.
If the region is empty then we have an occurrence of \(m_2\).
If there is only one point in the region we can choose this to replace the \(1\) in
the mesh pattern to get the required shading. If there is more than one point then
choosing the two leftmost points gives us the following mesh pattern.
\begin{equation*}
    \pattern{3,4}{2,3,1,4}{0/0, 0/1, 0/2, 0/3, 0/4,
                            1/0, 1/1, 1/2, 1/3, 1/4,
                            2/0, 2/1,2/2, 2/4,
                            3/1, 3/2, 3/3, 3/4,
                            4/1, 4/2, 4/3}
\end{equation*}
Where the two highlighted points are the original two points.
Now if we take the other two points as the points in our
mesh pattern then we get an occurrence of the pattern
we originally desired, and hence the two patterns are coincident.

The other reasoning applies to the patterns
\begin{equation*}
    m_1 = \pattern{}{1,2}{0/1,0/2,2/0,2/1} \text{ and } m_2 = \pattern{}{1,2}{0/1,0/2,1/1,2/0,2/1}
\end{equation*}
which are coincident by experimentation.

In order to prove this coincidence we will proceed by mathematical
induction on the number of points in region corresponding to the middle box.
We call this number \(n\).
\begin{description}
    \item [Base Case \((n=0)\):] The base case holds since we can freely shade the
        box if it contains no points.
    \item [Inductive Hypothesis \((n=k)\):] Suppose that we can find an occurrence
        of the second pattern if we have an occurrence of the first with \(k\) points
        in the middle box.
    \item [Inductive Step \((n=k+1)\)] Suppose that we have \((k+1)\) points in
        the middle box. Choose the bottom most point in the middle box, this gives the
        mesh pattern
     \begin{equation*}
        \begin{tikzpicture}[scale=0.5]
            \modpattern[5]{}{1,2,3}{0/1,0/2,0/3,1/1,2/1,3/0,3/1,3/2}
            \draw (2.5,0.5) node {\(X\)};
        \end{tikzpicture}
    \end{equation*}
        Now we need to consider the box labelled \(X\). If this box is empty then
        we have an occurrence of \(m_2\) and are done. If this box contains
        any points then we gain some extra shading on the mesh pattern
        due to the dominating pattern
     \begin{equation*}
        \begin{tikzpicture}[scale=0.5]
            \modpattern[5]{2,3}{1,2,3}{0/1,0/2,0/3,1/1,1/2,1/3,2/1,3/0,3/1,3/2}
            \filldraw (2.5,0.5) circle (4 pt);
        \end{tikzpicture}
    \end{equation*}
        The two highlighted points form an occurrence of \(m_1\) with \(k\)
        points in the middle box, and thus by the Inductive Hypothesis we
        are done.
\end{description}
By induction we have that every occurrence of \(m_1\) leads to an occurrence of \(m_2\)
and by \NoteRef{not:downcmesh} every occurrence of \(m_2\) is an occurrence of
\(m_1\) so the two patterns are coincident.
This argument applies to another two pairs of classes. Therefore in
total in \(\av{\perm{3,2,1}, \mperm{1,2}{R}}\) there are \(213\) coincidence
classes.

\chapter{Wilf-equivalences under dominating  patterns}

Another question often asked in thefield of permutation patterns is that
of Wilf-equivalence. Two patterns \(\pi\) and \(\sigma\) are said to be
Wilf-equivalent if their avoidance sets have the same size at each
length. More formally

\begin{definition}[Wilf-equivalence]
    Two patterns \(\pi\) and \(\sigma\) are said to be \emph{Wilf-equivalent}
    if for all \(k_{} \ge 0\), \(\size{\av[k]{\pi}} = \size{\av[k]{\sigma}}\).

    Two sets of permutation patterns \(R\) and \(S\) are are
    \emph{Wilf-equivalent} if for all \(k_{} \ge 0\),
    \(\size{\av[k]{R}} = \size{\av[k]{S}}\).
\end{definition}

Wilf-equivalence is of interest as if two permutation classes are enumerated
in the same way then there should exist a bijection between them, and therefore
any other combinatorial object that they represent.

Coincident pattern classes are also Wilf-equivalent. This is due to the fact that
if \(\av[k]{S} = \av[k]{R}\) then obviously \(\size{\av[k]{R}} = \size{\av[k]{S}}\).
Coincidence is therefore a stronger equivalence condition than Wilf-equivalence.

There are a number of symmetries we can use when examining Wilf-equivalences
to reduce the amount of work, it can be easily seen that the reverse, complement
and inverse operations preserve enumeration, and therefore these classes are trivially
Wilf-equivalent.

\begin{align*}
    \rev{\textpattern{}{2,3,1}{}} &= \textpattern{}{1,3,2}{}\\
    \com{\textpattern{}{2,3,1}{}} &= \textpattern{}{2,1,3}{}\\
    \inv{\textpattern{}{2,3,1}{}} &= \textpattern{}{3,1,2}{}
\end{align*}

The group of symmetries on permutations is isomorphic to the dihedral group
of order \(8\), the group of symmetries of a square. Composition of the above
symmetries gives the remaining 5 symmetries. If we consider generators of the
group the operations \emph{reverse-inverse} and \emph{reverse} correspond to
the generators of the dihedral group.

Since we are always considering Wilf-equivalences in the set \(\av{S}\) we
must only use these symmetries when they preserve the dominating pattern.

Throughout this section we will consider Wilf-equivalences of patterns
whilst avoiding the \emph{Dominating Pattern} \(\perm{2,3,1}\). We will
use \(\mathcal{C}\) to denote the set of these avoiders and \(C(x)\) will
be the usual Catalan generating function satisfying \(C(x) = 1 + C(x)^2\).
This is easy to see by structural decomposition around the maximum, as
shown in \FigureRef{fig:decompmax}

\begin{figure}[!ht]
    \centering
    \decompmax{0/1,1/0}{A}{}{}{B}
    \caption{Structural decomposition of a typical avoider of 231}
    \label{fig:decompmax}
\end{figure}

The elements to the left of the maximum, \(A\), have the structure of a \(231\)
avoiding permutation, and the elements to the right of the maximum, \(B\), have the
structure of a \(\perm{2,3,1}\) avoiding permutation. Furthermore, all the
elements in \(A\) lie below all of the elements in \(B\).

We can also decompose a permutation avoiding \(\perm{2,3,1}\) around the leftmost
point, giving a similar figure.

\section{Wilf-classes with patterns of length 1.}
When considering the mesh patterns of length \(2\) it will be useful to know
the Wilf-equivalence classes of the mesh patterns of length \(1\) inside
\(\av{\perm{2,3,1}}\).

It can be seen that the patterns in the following set are set equivalent.
\begin{equation*}
    \left\{
        \begin{array}{c}
        \textpattern{}{1}{},
        \textpattern{}{1}{0/0},
        \textpattern{}{1}{1/0},
        \textpattern{}{1}{0/0,1/0},
        \textpattern{}{1}{0/1},\\
        \textpattern{}{1}{0/0,0/1},
        \textpattern{}{1}{1/1},
        \textpattern{}{1}{1/0,1/1},
        \textpattern{}{1}{0/1,1/1}
    \end{array}
    \right\}
\end{equation*}
This is due to the fact that every permutation except the empty
permutation must contain an occurrence of all of these patterns.

The pattern \(\textpattern{}{1}{0/0,1/0,0/1,1/1}\) is in its own
Wilf-class since the only permutation containing this pattern is
the permutation \(\perm{1}\). The avoiders of this pattern therefore
have generating function \(E(x) = C(x)-x\).

The avoiders of the pattern \(p_1 = \textpattern{}{1}{0/1,1/0}\) can be
decomposed around the maximum element in order to give the following
structural decomposition.
\begin{equation*}
    \mathcal{F} = \varepsilon \sqcup \decompmax{0/1,1/0}{\mathcal{F}}{}{}{\mathcal{C}\setminus\varepsilon}
\end{equation*}
Since if the upper right section was empty the maximum would create an
occurrence of the pattern, however no points in this section can create
an occurrence since the maximum lies in a region corresponding to the
shading in \(p_1\). The lower right region however can create occurrences
of \(p_1\) and therefore must also avoid \(p_1\) as well as \(\perm{2,3,1}\).
This gives the generating function of avoiders to be the function \(F(x)\)
satisfying.
\begin{align*}
    F(x) &= 1 + xF(x)(C(x)-1) \\
    \shortintertext{Solving for \(F\) gives}
    F(x) &= \frac{1}{1+x-xC(x)}
\end{align*}

It can be shown by use of \PropositionRef{prop:dom2} that the patterns
\textpattern{}{1}{0/0,1/1} and \(p_2 = \textpattern{}{1}{0/0,0/1,1/1}\) are
set-equivalent. Consider the decomposition of an avoider of \(p_2\) in
\(\av{\perm{2,3,1}}\) around the maximum
\begin{equation*}
    \mathcal{G}_1 =\varepsilon \sqcup \decompmax{0/1,1/0}{\mathcal{C}\setminus\varepsilon}{}{}{\mathcal{C}}
\end{equation*}
This can be explained succinctly by the fact that a permutation containing
\(p_2\) starts with it's maximum, by not allowing the left part of the
\(\perm{2,3,1}\) avoider to be empty we prevent an occurrence from ever happening.

Consider \(p_3 = \textpattern{}{1}{0/1,1/0,1/1}\), avoiding this pattern
means that a permutation does not end with it's maximum. We can perform a similar
decomposition as before to get
\begin{equation*}
    \mathcal{G}_2 =\varepsilon \sqcup \decompmax{0/1,1/0}{\mathcal{C}}{}{}{\mathcal{C}\setminus\varepsilon}
\end{equation*}

Now consider \(p_4 = \textpattern{}{1}{0/0,0/1,1/0}\), the avoiders of this
pattern are permutations that do not start with their minimum. In this case we
perform the decomposition around the leftmost element
\begin{equation*}
    \mathcal{G}_3 =\varepsilon \sqcup \decompleft{0/1,1/0}{\mathcal{C}\setminus\varepsilon}{}{}{\mathcal{C}}
\end{equation*}
All of these classes have the same generating function namely
\begin{equation*}
    G(x) = 1 + xC(x)(C(x)-1)
\end{equation*}

There is one pattern of length \(1\) still to consider. The pattern
\(p_5 = \textpattern{}{1}{0/0,1/0,1/1}\) is avoided by all permutations
that do not end in their minimum. Considering the standard decomposition
of a \(\perm{2,3,1}\) avoider around the maximum we can see that an avoider
of \(p_5\) must fit into the following form precisely once.
\begin{equation*}
    \mathcal{H} = \varepsilon \sqcup
    \underbrace{\decompmax{0/1,1/0}{\mathcal{C}\setminus\varepsilon}{}{}{\mathcal{C}}}_\text{
        Minimum comes before the maximum.
    } \sqcup
    \underbrace{\decompmax{0/0,0/1,1/0}{}{}{}{\mathcal{H}}}_\text{
        Minimum is after the maximum, cannot be last.
    }
\end{equation*}
Therefore this particular class has generating function \(H(x)\) satisfying
\begin{equation*}
    H(x) = 1 + xC(x)(C(x)-1) + xH(x)
\end{equation*}

\section{Wilf-classes with patterns of length 2}
By use of the set equivalences from \ChapterRef{chap:coincs} we know there are
at most \(95\) Wilf-equivalence classes.

In order to consider symmetries we must only take the symmetries that preserve
the pattern \(\perm{2,3,1}\). If we take any of the symmetries alone the
permutation is different. The only symmetry that preserves the pattern
\(\perm{2,3,1}\) is that of \emph{reverse-complement-inverse}. Using this
set of symmetries to merge classes gives us \(61\) classes of trivial
Wilf-equivalences.

Computing avoiders up to length \(10\) gives us \(23\) Wilf-classes,
of which \(13\) are non-trivial.

When considering explanations of Wilf-equivalences we consider how the
permutations correspond to set-partitions.

\begin{note}
    \label{not:setequiv}
    The avoiders of the pattern \(\mperm{2,3,1}{\{(1,0),(1,1),(1,2),(1,3)\}}\)
    in \(\mathfrak{S}_n\) are in one-to-one correspondence with partitions of
    \(\nrange{n}\). \cite{DBLP:journals/ejc/Claesson01}
\end{note}
\begin{proof}
    Let \(\pi\) be a permutation in \(\av[n]{\perm{2,3,1}}\) take the permutation
    in one-line notation and insert a dash between each ascent in \(\pi\). This
    corresponds to set partitions where the blocks are the elements between the
    dashes, the blocks are listed in increasing order of their least element,
    with the elements written in each block in descending order.
\end{proof}

\begin{example}
    Given the permutation \(\pi = \perm{5,4,2,1,3,9,6,8,7}\) this corresponds to the
    partition \(\{\{1,2,4,5\},\{3\},\{6,9\},\{8,7\}\}\).
\end{example}

We are looking at permutations in \(\av{\perm{2,3,1}}\), all of these permutations
also avoid the vincular pattern in \NoteRef{not:setequiv}.

\subsection{}
The set/symmetry classes containing the following patterns are
Wilf-equivalent in \(\av{\perm{2,3,1}}\) but are not set equivalent or
symmetries of each other
\begin{equation*}
    m_1 = \pattern{}{1,2}{0/0,0/2,
                    1/0,1/1,
                    2/0,2/1,2/2}
    \text{ and }
    m_2 = \pattern{}{1,2}{0/0,
                    1/0,1/1,1/2,
                    2/0,2/1,2/2}
\end{equation*}
In this case it is better to consider the containers of the patterns
instead of the avoiders due to the amount of shadings in the mesh.

We look at the containers of the pattern \(m_1\), there can only ever
be one occurrence of this pattern in a permutation corresponding to
the last point in the permutation and the minimum. The only form that
points in either of the two boxes can take is a decreasing sequence.
For a permutation of length \(k\) if we fix the number of points in
one of the boxes the number of points in the other box is determined.
Therefore we can have any number of points from \(\setrange{0}{k-2}\)
points in the bottom box. Therefore there are \(k-1\) containers of
length \(k\). These permutations correspond to set partitions of \(k\)
points into exactly two non-overlapping parts partitioned by the first
element and the minimum.

Now consider the containers of \(m_2\), we know that the unshaded region
must contain a decreasing subsequence, with the point corresponding to
the \(1\) in the mesh pattern. This decreasing subsequence has \(k-1\)
points, we can put the point corresponding to the \(2\) above any of these
points and therefore there are \(k-1\) containers of length \(k\).

Therefore these two patterns have been shown to have the same number of
avoiders of length \(k\) for all \(k\) and therefore all Wilf-equivalent.

\subsection{}
The classes containing the following patterns are Wilf-equivalent
when avoiders are considered in \(\av{\perm{2,3,1}}\)
\begin{align}
    m_1 = \pattern{}{1,2}{0/0,0/1,1/1,1/2,2/0,2/1} &\text{ and }
    m_2 = \pattern{}{2,1}{0/0,0/1,1/1,1/2,2/0,2/1} \label{eq:singleud}\\
    \text{ and }
    m_3 = \pattern{}{2,1}{0/1,0/2,1/0,1/1,1/2,2/1} &\text{ and }
    m_4 = \pattern{}{2,1}{0/1,0/2,1/0,1/1,1/2,2/0} \label{eq:other}
\end{align}

First we consider the coincidence in two parts then consolidate these parts.
Consider the coincidence shown in \eqref{eq:singleud}. The easiest way to
show that these are equinumerous is to consider the containers in the realm
of set partitions.

Due to the shading we know the following about the points corresponding
to the points in the patterns.
\begin{itemize}
    \item The point corresponding to the first point in both patterns
        must lie in the first block of our set partition (there are
        no points southwest from it in the permutation).
    \item The point corresponding to the second point in both patterns
        is a block bottom (there are no points southeast of it in
        the permutation).
    \item The block containing the point corresponding to the second
        point in both patterns contains only the point (it is a singleton
        block).
\end{itemize}
This tells us that an occurence of the patterns must happen when there
is a singleton block occurring after the first block. The difference
between the patterns is in the underlying classical pattern. This
means that permtuations containing \(m_1\) correspond to set partitions
with a singleton block with value one higher than some element in the
block containing \(1\). The permutations containing \(m_2\) correspond
to the set partitions containing a block with block bottom having value
one lower than some element in the block containing \(1\) and if this
block is not the block containing \(1\) then it is a singleton block.

Consider an avoider of \(\perm{2,3,1}\) and \(m_3\). We can perform
the decomposition around the maximum
\begin{equation*}
    \mathcal{I}_1 = \varepsilon \sqcup
    \decompmax{1/0,0/1}{\mathcal{I}_1}{}{}{\mathcal{G}_1}
\end{equation*}
Since only the first point in the top right region can
create an occurrence of \(m_3\) if and only if it is the
maximum in this region we must avoid starting with the maximum.

Looking at avoiders of \(\perm{2,3,1}\) and \(m_4\) we
can perform the same decomposition around the maximum to get
\begin{equation*}
    \mathcal{I}_2 = \varepsilon \sqcup
    \decompmax{1/0,0/1}{\mathcal{I}_2}{}{}{\mathcal{G}_3}
\end{equation*}
Since an occurrence of \(m_4\) can never occur in the
top right region, and could only occur between the first point
in the region and the maximum, if and only if this first point is the minimum.
Since both \(\mathcal{G}_1\) and \(\mathcal{G}_3\) have the same
enumeration, \(\mathcal{I}_1\) and \(\mathcal{I}_2\) must also have
the same enumeration and are therefore Wilf-equivalent.

Now we must consolidate these two subclasses. In order to do this we must
consider the decomposition around the leftmost point of a permutation in
\(\av{\perm{2,3,1},m_1}\) we gain the following.
\begin{equation*}
    \mathcal{I}_3 = \varepsilon \sqcup
    \decompleft{1/0,0/1}{\mathcal{I}_3}{}{}{\mathcal{G}_3}
\end{equation*}
It is therefore obvious that avoiders of \(m_1\) and avoiders of \(m_4\)
have the same enumeration, and therefore all four patterns are Wilf-equivalent
in \(\av{\perm{2,3,1}}\) with generating function satisfying
\begin{equation*}
    I(x) = 1 + xI(x)G(x)
\end{equation*}

\subsection{}
The classes containing the following patterns are Wilf-equivalent
in \(\av{\perm{2,3,1}}\)
\begin{equation*}
    m_1 = \pattern{}{1,2}{0/0,0/1,0/2,
                    1/0,1/1,1/2,
                    2/0,2/1,2/2}
    \text{ and }
    m_2 = \pattern{}{2,1}{0/0,0/1,0/2,
                    1/0,1/1,1/2,
                    2/0,2/1,2/2}
\end{equation*}
It is obvious that these two are Wilf-equivalent since the only
permutations that contain these patterns are \(\perm{1,2}\) and
\(\perm{2,1}\) respectively, therefore the avoiders of these patterns
are counted by the Catalan numbers at all lengths except for length
\(2\) where there is precisely \(1\) avoider.

\subsection{}

The classes containing the following patterns are Wilf-equivalent
in \(\av{\perm{2,3,1}}\)
\begin{equation*}
    m_1 = \pattern{}{2,1}{0/0,0/1,0/2,
                    1/1,1/2,
                    2/0,2/2}
    \text{ and }
    m_2 = \pattern{}{1,2}{0/1,0/2,
                    1/0,1/2,
                    2/0,2/1,2/2}
\end{equation*}

First consider the structure of an avoider of \(m_1\) and \(\perm{2,3,1}\)
we can perform the usual structural decomposition of an avoider
of \(\perm{2,3,1}\) where we consider decomposition around the maximum.
Any permutation in \(\av{\perm{2,3,1},m_1}\) is in the following set
precisely once.
\begin{equation*}
    \mathcal{J}_1 = \varepsilon \sqcup
    \decompmax{0/1,1/0}{\mathcal{C}\setminus\varepsilon}{}{}{\mathcal{C}} \sqcup
    \decompmax{0/1,1/0,0/0}{}{}{}{\mathcal{F}}
\end{equation*}

Now consider the decomposition around the maximum of a permutation in
\(\av{\perm{2,3,1},m_2}\) This fits into the following set.
\begin{equation*}
    \mathcal{J}_2 = \varepsilon \sqcup
    \decompmax{0/1,1/0}{\mathcal{C}}{}{}{\mathcal{C}\setminus\varepsilon} \sqcup
    \decompmax{0/1,1/0,1/1}{\mathcal{F}}{}{}{}
\end{equation*}
Therefore both of these sets of avoiders are enumerated in the same manner
having generating function satisfying
\begin{equation*}
    J(x) = 1 + xC(x)(C(x)-1) + xF(x)
\end{equation*}

\subsection{}
Consider the containers of the patterns
\begin{equation*}
    \pattern{}{1,2}{0/0,0/1,0/2,1/0,1/2,2/0,2/1,2/2},
    \pattern{}{1,2}{1/1,0/1,0/2,1/0,1/2,2/0,2/1,2/2},
    \pattern{}{2,1}{0/0,0/1,0/2,1/0,1/2,1/1,2/1,2/2},
    \pattern{}{2,1}{0/0,0/1,0/2,1/0,1/2,2/0,1/1,2/2},
    \pattern{}{2,1}{0/0,0/1,0/2,1/0,1/2,2/0,2/1,1/1},
    \pattern{}{2,1}{0/0,0/1,0/2,1/1,1/2,2/0,2/1,2/2}
\end{equation*}

For each of these patterns there is precisely one occurrence
in any permutation containing the pattern. Now consider
the points in the free box in each case. Each of these regions
must contain an avoider of \(\perm{2,3,1}\) that is of length
\(n-2\). Therefore these classes are all Wilf-equivalent
and the number of length \(n\) avoiders is
\begin{equation*}
    K_n = C_n - C_{n-2}
\end{equation*}
for \(n \ge 2\) where \(C_n\) is the \(n\)th Catalan number, the number of
\(\perm{2,3,1}\) avoiders of length \(n\).

\subsection{}
Now consider the containers of the patterns
\begin{equation*}
    \pattern{}{2,1}{0/0,0/1,0/2,1/0,1/2,2/0,2/1,2/2},
    \pattern{}{1,2}{0/0,1/1,0/2,1/0,1/2,2/0,2/1,2/2},
    \pattern{}{1,2}{0/0,0/1,1/1,1/0,1/2,2/0,2/1,2/2}
\end{equation*}
Each of these patterns again occurs precisely once in
any containing permutation. However this time when considering
the free box we need to take into consideration \LemmaRef{lem:incdec}
and so the empty box can only contain a decreasing subsequence. There
is precisely one decreasing subsequence at every length, and so there
is exactly one container of each pattern at each length. The three
patterns are Wilf-equivalent and have \(C_n - 1\) avoiders of
length \(n\) for all \(n \ge 2\).

\subsection{}
The containers of the following patterns can only
have exactly one occurrence.
\begin{equation*}
    \pattern{}{1,2}{0/0,0/1,0/2,1/1,2/0,2/1,2/2}
    \pattern{}{2,1}{0/0,0/1,1/0,1/1,1/2,2/0,2/2}
    \pattern{}{1,2}{0/0,0/1,1/1,1/2,2/0,2/1,2/2}
    \pattern{}{1,2}{0/1,0/2,1/0,1/1,2/0,2/1,2/2}
    \pattern{}{2,1}{0/0,0/1,0/2,1/0,1/2,2/0,2/1}
    \pattern{}{1,2}{0/0,0/1,1/0,1/2,2/0,2/1,2/2}
\end{equation*}
Once again we consider the free boxes, obviously for every pattern
except the first the two boxes are independent, and one contains
any avoider of \(\perm{2,3,1}\) and the other must contain a
decreasing sequence by \LemmaRef{lem:incdec}.
Let us consider the first pattern on it's own. In order to
avoid \(\perm{2,3,1}\) across the free boxes we can add some additional
restrictions
\begin{equation}
    \raisebox{0.6ex}{
    \begin{tikzpicture}[scale=1, baseline=(current bounding box.center)]
        \fill[pattern=north east lines, pattern color=black!75] (1.5, 0) rectangle +(0.5,1);
        \fill[pattern=north east lines, pattern color=black!75] (1.0, 2) rectangle +(0.5,1);
        \draw (1.5,0) -- (1.5,1)
              (1.5,2) -- (1.5,3);
        \modpattern[2]{}{1,2}{0/0,0/1,0/2,1/1,2/0,2/1,2/2}
    \end{tikzpicture}
    }
    \label{eqn:splitboxes}
\end{equation}
Now we can see that the top free box must contain a decreasing sequence,
and the bottom must contain an avoider of \(\perm{2,3,1}\) and these
two do not interact in any manner. The containers of this pattern are counted
the same as the other patterns, and due to this they are Wilf-equivalent.

\subsection{}
The classes containing the following patterns are Wilf-equivalent
\begin{equation*}
    \begin{aligned}
        m_1 = \pattern{}{1,2}{0/0,0/1,1/1,1/2,2/0,2/2} &\text{ and }
        m_2 = \pattern{}{1,2}{0/0,0/1,1/0,2/0,2/1,2/2} \\
        \text{ and }
        m_3 = \pattern{}{1,2}{0/2,1/0,1/1,2/0,2/1,2/2} &\text{ and }
        m_4 = \pattern{}{1,2}{0/1,1/0,1/1,2/0,2/1,2/2}
    \end{aligned}
\end{equation*}

First consider the decomposition of avoiders of \(m_1\) in \(\av{\perm{2,3,1}}\)
around the maximum.
\begin{equation*}
    \mathcal{L}_1 =
    \varepsilon \sqcup
    \decompmax{0/1,1/0}{\scriptstyle \mathcal{G}_1\setminus\varepsilon}{}{}{\mathcal{C}}
    \sqcup
    \decompmax{0/0,0/1,1/0}{}{}{}{\mathcal{L}_1}
\end{equation*}

Now we decompose the avoiders of \(m_2\) around the leftmost point
\begin{equation*}
    \mathcal{L}_2 =
    \varepsilon \sqcup
    \decompleft{0/1,1/0}{\mathcal{C}\setminus\varepsilon}{}{}{\mathcal{C}\setminus\varepsilon}
    \sqcup
    \decompleft{0/1,1/0,1/1}{\mathcal{L}_2}{}{}{}
\end{equation*}
This gives us two generating functions satisfying the following pair of equations
\begin{align}
    L_1(x) &= 1 + xC(x)(G(x)-1) +xL_1(x) \label{eqn:genL1}\\
    \text{ and }L_2(x) &= 1 + x(C(x)-1)^2 + xL_2(x) \label{eqn:genL2}
\end{align}
In order for these two functions to give the same value it is necessary to show that
\eqref{eqn:genL1} and \eqref{eqn:genL2} are equal, this occurs if
\(C(x)(G(x)-1) = (C(x)-1)^2\).
\begin{align*}
    &&(C(x)-1)^2 &= C(x)(G(x)-1)\\
        &\Leftrightarrow& x^2C^4(x) &=  xC(x)(C(x)-1)C(x) &&\text{By definition of \(G\) and \(C\)}\\
        &\Leftrightarrow& x^2C^4(x)&=xC^3(x)-xC^2(x)\\
        &\Leftrightarrow& xC^2(x) &= C(x) -1 &&\text{Divide by \(xC^2(x)\)}\\
        &\Leftrightarrow& C(x) &= 1 +xC^2(x)
\end{align*}
The final line is always satisfied since it is the form of \(C(x)\), and
therefore the two generating functions are equal.

Now we look at the other patterns. In particular note that any container of these
patterns can contain the pattern precisely once, \(m_2\) specifies the
minimum and last point, \(m_3\) and \(m_4\) both use the last point and the
previous block bottom (in the set partition context).

Looking at the structure of a container of \(m_3\) in \(\av{\perm{2,3,1}}\)
\begin{equation*}
    \raisebox{0.6ex}{
        \begin{tikzpicture}[scale=1, baseline=(current bounding box.center)]
            \modpattern[2]{}{1,2}{0/2,1/0,1/1,2/0,2/1,2/2}
            \draw (0.5,0.5) node {\(b_1\)};
            \draw (0.5,1.5) node {\(b_2\)};
            \draw (1.5,2.5) node {\(b_3\)};
        \end{tikzpicture}
    }
\end{equation*}
The boxes \(b_2\) and \(b_3\) must contain a decreasing sequence by
\LemmaRef{lem:incdec}. The box labelled \(b_1\) must contain an avoider
of \(\perm{2,3,1}\). However note that the points in this box can have
interaction with any points in box \(b_2\). If there is just one point
in \(b_2\) then any points in \(b_1\) to the left of this point must
be lower than any points to the right of this point. By extension,
if \(b_2\) contains a decreasing sequence with \(k\) points, there
are \(k+1\) non-interacting avoiders of \(\perm{2,3,1}\) in \(b_3\).

Now in \(m_2\) and \(m_4\) containers we can use the same method as in
\eqref{eqn:splitboxes} to separate the two decreasing sequences in the
free boxes in the top row, and the mixing happens in the same manner
as in a container of \(m_3\).
We now have that \(m_2,m_3\) and \(m_4\) have the same number of containers
so are Wilf-equivalent, and that \(m_1\) and \(m_2\) have the same generating
function so all four classes are Wilf-equivalent.
\subsection{}
The classes containing the following patterns are Wilf-equivalent
\begin{equation*}
    \pattern{}{1,2}{} \text{ and }
    \pattern{}{2,1}{}
\end{equation*}
This is true since the only avoiders of these patterns are the decreasing
sequence and the increasing sequence respectively, and both of these
avoid \(\perm{2,3,1}\) in all cases. There is therefore \(1\) avoider at
every length.

\subsection{}
The patterns
\begin{equation*}
    \pattern{}{1,2}{0/0,0/2,1/0,1/2,2/0,2/2},
    \pattern{}{1,2}{1/1,2/1,1/0,1/2,2/0,2/2},
    \pattern{}{1,2}{0/1,1/0,1/2,2/0,2/2},
    \pattern{}{1,2}{0/1,0/2,1/1,1/2,2/0}, \text{ and }
    \pattern{}{1,2}{0/0,0/1,0/2,1/0,1/1,1/2}
\end{equation*}
are Wilf-equivalent.
If \(\mathcal{M}_1\) is the set of avoiders of
\textpattern{}{1,2}{0/0,0/2,1/0,1/2,2/0,2/2}, then by the structural
decomposition around the maximum we have
\begin{equation*}
    \mathcal{M}_1 = \varepsilon \sqcup
    \decompmax{0/0,0/1,1/0}{}{}{}{\mathcal{C}}
\end{equation*}
This is because \(\mathcal{M}_1\) is the set of permutations
who have their minimum occur after their maximum.

% If \(\mathcal{M}_2\) is the set of avoiders of
% \textpattern{}{1,2}{1/1,2/1,1/0,1/2,2/0,2/2}, then by structural
% decomposition around the maximum we have
% \begin{equation*}
%     \mathcal{M}_2 = \varepsilon \sqcup
%     \decompmax{0/1,1/0}{\mathcal{C}}{}{}{\scriptstyle\mathcal{M}_1\setminus\varepsilon}
%     \sqcup \tikz{\filldraw (0,0) circle (3pt);}
% \end{equation*}
% Alternatively,
The pattern \textpattern{}{1,2}{1/1,2/1,1/0,1/2,2/0,2/2} occurs if
the last element is higher than the penultimate element. This can only occur
if the last element is in a single block in the set partition context, In order
to construct a avoider of length \(n\) we can take any avoider of
\(\perm{2,3,1}\) of length \(n-1\) and insert the new maximum into the last
block.  This ensures that the last block is never a singleton. This means
that these permutations are also counted by \(M(x) = 1+xC(x)\).

The avoiders of the third pattern can be decomposed by the maximum
to give
\begin{equation*}
    \mathcal{M}_3 = \varepsilon \sqcup
    \decompmax{0/1,1/0}{\mathcal{F}}{}{}{\mathcal{M}_3}
\end{equation*}
The generating function derived satisfies \(M_3(x) = 1 + xF(x)M_3(x)\).
The fourth pattern can be decomposed around the maximum in a similar
manner.
\begin{equation*}
    \mathcal{M}_4 = \varepsilon \sqcup
    \decompmax{0/1,1/0}{\mathcal{M}_4}{}{}{\mathcal{F}}
\end{equation*}

Finally considering the last pattern, the only way we can construct an
avoider is to take any \(\perm{2,3,1}\) avoider and add a new minimum
at the start of the permutation. Adding a new leftmost point in any other
position would either create an occurrence of \(\perm{2,3,1}\) or the pattern.
Therefore this is also counted by \(M(x) = 1+xC(x)\).

We need to show that the generating function \(M_3(x)\) is the same as \(M(x)\)
\begin{equation*}
    \begin{aligned}
        M_3(x) &= 1 + xF(x)M_3(x) \\
        &= \frac{1}{1-xF(x)} && \text{Solving for \(M_3(x)\)}\\
        &= \frac{1}{1 - \frac{x}{1+x-xC(x)}} &&\text{Substituting for \(F(x)\)}\\
        &= \frac{1-xC(x) +x}{1-xC(x)}\\
        &= 1+xC(x)
    \end{aligned}
\end{equation*}
We have that \(M_3(x) = 1+xC(x)= M(x) \) so all four patterns are Wilf-equivalent.

\subsection{}
The patterns
\begin{equation*}
    \begin{aligned}
        m_1 = \pattern{}{1,2}{0/0,0/1,1/0,1/2,2/1} \text{ and }
        m_2 = \pattern{}{2,1}{0/0,0/1,0/2,1/1,1/2,2/0}\\
        \text{ and }
        m_3 = \pattern{}{1,2}{0/1,0/2,1/0,1/2,2/0,2/2} \text{ and }
        m_4 = \pattern{}{2,1}{1/0,1/1,2/1,2/2}\\
    \end{aligned}
\end{equation*}
Can be shown to be Wilf-equivalent. First we consider a container of \(m_2\) in \(\av{\perm{2,3,1}}\)
\begin{equation*}
    \raisebox{0.6ex}{
    \begin{tikzpicture}[scale=1, baseline=(current bounding box.center)]
        \fill[pattern=north east lines, pattern color=black!75] (2.5, 1) rectangle +(0.5,1);
        \fill[pattern=north east lines, pattern color=black!75] (2.0, 2) rectangle +(0.5,1);
        \draw (2.5,1) -- (2.5,3);
        \draw (1.5,0.5) node {\(b_1\)};
        \draw (2.25,1.5) node {\(\scriptstyle b_2\)};
        \draw (2.75,2.5) node {\(\scriptstyle b_3\)};
        \modpattern[2]{}{2,1}{0/0,0/1,0/2,1/1,1/2,2/0}
    \end{tikzpicture}
    }
\end{equation*}
We can choose the lowest occurrence of \(m_2\) without loss of generality.
The region corresponding to \(b_1\) must avoid the pattern \(\textpattern{}{1}{0/1,1/0}\)
as well as \(\perm{2,1,3}\). The regions corresponding to \(b_2\) and \(b_3\) must
now contain avoiders of \(\perm{2,3,1}\), these regions cannot mix in order to avoid
\(\perm{2,3,1}\). Since we already have an occurrence of
\(m_2\) we do not need to care about creating more occurrences so there are no
other conditions on these boxes.

Now looking at a container of \(m_3\) in \(\av{\perm{2,3,1}}\)
\begin{equation*}
    \raisebox{0.6ex}{
    \begin{tikzpicture}[scale=1, baseline=(current bounding box.center)]
        \fill[pattern=north east lines, pattern color=black!75] (1, 1.5) rectangle +(1,0.5);
        \fill[pattern=north east lines, pattern color=black!75] (2, 1) rectangle +(1,0.5);
        \draw (1,1.5) -- (3,1.5);
        \draw (0.5,0.5) node {\(b_1\)};
        \draw (1.5,1.25) node {\(\scriptstyle b_2\)};
        \draw (2.5,1.75) node {\(\scriptstyle b_3\)};
        \modpattern[2]{}{1,2}{0/1,0/2,1/0,1/2,2/0,2/2}
    \end{tikzpicture}
    }
\end{equation*}
We consider the leftmost occurrence of \(m_3\), the region corresponding
to \(b_1\) must avoid the pattern \(\textpattern{}{1}{0/1,1/0}\) as well as
\(\perm{2,3,1}\) once more. The regions corresponding to \(b_2\) and \(b_3\)
must avoid \(\perm{2,3,1}\) and as in a container of \(m_2\) these regions cannot
mix, as doing so would lead to an occurrence of \(\perm{2,3,1}\).
Therefore both of these sets of containers are enumerated in the same way.

Now we find a structural decomposition for an avoider of \(m_2\).
Decomposing around the maximum we see the set of avoiders of \(m_2\)
have the form
\begin{equation*}
    \mathcal{N}_2 = \varepsilon \sqcup
    \decompmax{0/1,1/0}{\scriptstyle\mathcal{N}_2\setminus\varepsilon}{}{}{\mathcal{C}}
    \sqcup
    \decompmax{0/0,0/1,1/0}{}{}{}{\mathcal{F}}
\end{equation*}
We can decompose an avoider of \(m_1\) in \(\av{\perm{2,3,1}}\) around the leftmost point
in a similar manner
\begin{equation*}
    \mathcal{N}_1 = \varepsilon \sqcup
    \decompleft{0/1,1/0}{\scriptstyle\mathcal{N}_1\setminus\varepsilon}{}{}{\mathcal{C}}
    \sqcup
    \decompleft{0/0,0/1,1/0}{}{}{}{\mathcal{F}}
\end{equation*}
These two decompositions tell us that these two patterns are Wilf-equivalent
and have generating function
\begin{equation}
    N_1 = 1 + x(N_1(x)-1)C(x) +xF(x)
    \label{eqn:ngen}
\end{equation}
Now consider an avoider of \(m_4\) decomposed around the maximum
\begin{equation*}
    \mathcal{N}_4 = \varepsilon \sqcup
    \decompmax{0/1,1/0}{\mathcal{C}}{}{}{\mathcal{O}}
\end{equation*}
Where \(\mathcal{O}\) is the permutations avoiding \(\perm{2,3,1},m_4\) and
\(p = \textpattern{}{1}{0/0,0/1,1/1}\) since if the subsequence in this box
were to start with the maximum then this point and the maximum would create
an occurrence of \(m_4\).
Now consider decomposition of a permutation in \(\mathcal{O}\) around its
leftmost point.
\begin{equation*}
    \mathcal{O} = \varepsilon \sqcup
    \decompleft{0/1,1/0}{\mathcal{C}}{}{}{\scriptstyle\mathcal{N}_4\setminus\varepsilon}
\end{equation*}
This gives us the generating function
\begin{equation*}
    O(x) = 1 + xC(x)(N_4(x)-1)
\end{equation*}
Now we can construct the following for \(N_4\)
\begin{equation}
    N_4(x) = 1 +xC(x)(xC(x)(N_4(x)-1)+1)
    \label{eqn:ngen2}
\end{equation}

All that remains to show Wilf-equivalence is to show that \EquationRef{eqn:ngen} and
\EquationRef{eqn:ngen2} are the same generating function.
First solve \EquationRef{eqn:ngen2} for \(N_4(x)\)
\begin{equation}
    \begin{aligned}
        N_4(x) &= 1+xC(x)(xC(x)(N_4(x)-1)+1)\\
        &= 1 + x^2C^2(x)N_4(x) -x^2C^2(x) +xC(x)\\
        &= 1 + \frac{xC(x)}{1 - x^2C^2(x)}\\
        &= 1 + \frac{xC(x)}{(1-xC(x))(1+xC(x))} && \text{Difference of squares}\\
        N_4(x)&= 1 + \frac{xC^2(x)}{1+xC(x)} && C(x) = \frac{1}{1-xC(x)}
    \end{aligned}
\end{equation}

Now we solve \EquationRef{eqn:ngen} for \(N_1(x)\)
\begin{equation}
    \begin{aligned}
        N_1(x)&= 1+x(N_1(x)-1)C(x)+xF(x)\\
        &= 1+xN_1(x)C(x)-xC(x)+\frac{x}{1+x-xC(x)} &&\text{Substitution of \(F(x)\)}\\
        N_1(x)(1-xC(x))&= \frac{x^2C^2(x) - (x^2+2x)C(x) + 2x+1}{1+x-xC(x)}\\
        N_1(x) &= \frac{x^2C^2(x) - (x^2+2x)C(x) + x+1+x}{x^2C^2(x) - (x^2+2x)C(x) +x+1}\\
        &= 1 +\frac{x}{x^2C^2(x) - (x^2+2x)C(x) +x+1}\\
        &= 1 + \frac{x}{1-x^2C(x)-xC(x)} && xC^2(x) = C(x)-1\\
        &= 1 + \frac{xC^2(x)}{C^2(x) - xC^3(x)(x+1)}\\
        &= 1 + \frac{xC^2(x)}{C(x) - xC^2(x) + xC(x)} && xC^2(x) = C(x)-1\\
        N_1(x)&= 1 + \frac{xC^2(x)}{1+xC(x)} && C(x) = 1+xC^2(x)\\
    \end{aligned}
\end{equation}
We have shown that \(N_1\) and \(N_4\) are indeed the same generating function,
and we have that the classes containing these four patterns are Wilf-equivalent.

\subsection{}
\begin{equation*}
    m_1 = \pattern{}{1,2}{1/0,1/1,1/2,2/0,2/1} \text{ and }
    m_2 = \pattern{}{2,1}{0/1,1/0,1/1,1/2,2/1}
\end{equation*}
Let \(\mathcal{P}_1\) be the set of avoiders of \(m_1\), by structural decomposition
around the leftmost point we have
\begin{equation*}
    \mathcal{P}_1 = \varepsilon \sqcup
    \decompleft{0/1,1/0}{\mathcal{P}_1}{}{}{\mathcal{P}^\prime_1}
\end{equation*}
Where \(\mathcal{P}^\prime_1\) is a permutation avoiding \(\perm{2,3,1}, m_1\)
and \(\textpattern{}{1}{0/0,0/1,1/1}\)
Now consider the decomposition of a permutation in \(\mathcal{P}^\prime_1\) it can
once again be decomposed around the leftmost point
\begin{equation*}
    \mathcal{P}^\prime_1 = \varepsilon \sqcup
    \decompleft{0/1,1/0}{\scriptstyle \mathcal{P}_1\setminus\varepsilon}{}{}{\mathcal{P}^\prime_1}
\end{equation*}
This is a complete decomposition of avoiders of \(m_1\). Pow we
look at an avoider of \(m_2\), this time decomposition is around the maximum
\begin{equation*}
    \mathcal{P}_2 = \varepsilon \sqcup
    \decompleft{0/1,1/0}{\mathcal{P}_2}{}{}{\mathcal{P}^\prime_2}
\end{equation*}

Again we use the same method of decomposition of a permutation in \(\mathcal{P}^\prime_2\)
\begin{equation*}
    \mathcal{P}^\prime_2 = \varepsilon \sqcup
    \decompleft{0/1,1/0}{\scriptstyle \mathcal{P}_2\setminus\varepsilon}{}{}{\mathcal{P}^\prime_2}
\end{equation*}

This gives us a generating function \(P(x)\) satisfying
\begin{align}
    P(x) &= 1 + xP(x)P^\prime(x) \label{eqn:Pgen}\\
    P(x) &= 1 + x(P(x)-1)P^\prime(x)\label{eqn:PprimeGen}
\end{align}

Solving \EquationRef{eqn:PprimeGen} for \(P^\prime(x)\) and substituting into
\EquationRef{eqn:Pgen} gives us the fact that a the generating function for
\(P(x)\) satisfies
\begin{equation}
    P(x) = xP^2(x) - x(P(x) - 1) + 1
\end{equation}

\subsection{}
\begin{gather*}
        \pattern{}{2,1}{0/0,0/1,0/2,1/0,1/1,1/2,2/2},
        \pattern{}{1,2}{0/1,0/2,1/0,1/1,1/2,2/0,2/2},
        \pattern{}{2,1}{0/0,0/1,0/2,1/0,1/1,1/2,2/1},
        \pattern{}{2,1}{0/0,0/1,0/2,1/1,1/2,2/0,2/1},
        \pattern{}{2,1}{0/0,0/1,0/2,1/0,1/1,1/2,2/0},\\
        \pattern{}{1,2}{0/0,0/1,0/2,1/2,2/0,2/1,2/2},
        \pattern{}{1,2}{0/0,0/1,0/2,1/1,1/2,2/0,2/2},
        \pattern{}{1,2}{0/0,0/1,1/0,1/1,1/2,2/0,2/1},
        \pattern{}{1,2}{0/0,0/1,1/0,1/1,1/2,2/0,2/1}
\end{gather*}

In order to gain enumeration, consider decomposition of avoiders
of \(\textpattern{}{2,1}{0/0,0/1,0/2,1/0,1/1,1/2,2/2}\) around
the maximum.
\begin{equation*}
    \mathcal{Q}_1 = \varepsilon \sqcup
    \tikz{\filldraw (0,0) circle (3pt);}
    \sqcup
    \decompmax{0/1,1/0}{\mathcal{C\setminus\varepsilon}}{}{}{\mathcal{C}}
\end{equation*}

This gives us the following generating function
\begin{equation}
    Q(x) = 1 + x + xC(x)(C(x)-1) \label{eq:qgen}
\end{equation}

Now we consider decomposition of an avoider
\textpattern{}{1,2}{0/1,0/2,1/0,1/1,1/2,2/0,2/2} around the maximum,
this avoider must fit into the following form
\begin{equation*}
    \mathcal{Q}_2 = \varepsilon \sqcup
    \decompmax{0/1,1/0}{\mathcal{G}_2}{}{}{\mathcal{C}}
\end{equation*}
This gives us the generating function
\begin{equation*}
    \begin{aligned}
        Q_2(x) &= 1+xC(x)G(x)\\
        &= 1+xC(x)(1+xC(x)(C(x)-1))\\
        &= 1 +xC(x)(C(x)-xC(x)) && C(x) = 1+xC^2(x)\\
        &= 1 +x + xC^2(x) - xC(x) \\
        &= 1 +x + xC(x)(C(x)-1)
    \end{aligned}
\end{equation*}
Therefore this generating function is the same as \EquationRef{eq:qgen}.
We can decompose
\textpattern{}{2,1}{0/0,0/1,0/2,1/0,1/1,1/2,2/1},
\textpattern{}{2,1}{0/0,0/1,0/2,1/1,1/2,2/0,2/1},
\textpattern{}{2,1}{0/0,0/1,0/2,1/0,1/1,1/2,2/0},
\textpattern{}{1,2}{0/0,0/1,0/2,1/2,2/0,2/1,2/2},
and \textpattern{}{1,2}{0/0,0/1,0/2,1/1,1/2,2/0,2/2} around the leftmost point
into an avoider of one of the patterns with generating function \(G(x)\) and
an avoider of \(\perm{2,3,1}\).

Now decompose an avoider of \textpattern{}{1,2}{0/0,0/1,1/0,1/1,1/2,2/0,2/1}
around the leftmost point.
\begin{equation*}
    \mathcal{Q}_3 = \varepsilon \sqcup
    \decompleft{0/1,1/0}{\scriptstyle \mathcal{Q}\setminus\varepsilon}{}{}{\mathcal{C}}
    \sqcup \decompleft{0/0,0/1,1/0}{}{}{}{\mathcal{G}_3}
\end{equation*}
This gives generating function \(Q_3(x)\) satisfying
\begin{equation*}
    \begin{aligned}
        Q(x) &= 1 +xC(x)(P(x)-1) + xG(x)\\
        &= C(x) - xC^2(x) + xC(x) +x^2C^3(x) - x^2C^2 && \text{Solving for \(Q(x)\)}\\
        &= 1 + xC(x) +x^2C^3(x)-xC^2(x) && C = 1+xC^2(x)\\
        &= 1 + x + x^2C^3(x) && xC^2(x) = C(x)-1\\
        &= 1 + x + xC(x)(C(x)-1)
    \end{aligned}
\end{equation*}
This is equivalent to \EquationRef{eq:qgen}, and therefore
the classes containing all of these patterns are Wilf-equivalent.

\chapter{Conclusions and Future work}
From \ChapterRef{chap:coincs} it is can be seen that automatically classifying
coincidences of mesh patterns is a difficult task, establishing rules for longer
dominating patterns requires many more cases to be taken. It would however
be interesting to consider the self application of the  third Dominating Rule to mesh
patterns in order to try to capture some of the coincidences described in
\textcite{DBLP:journals/combinatorics/HilmarssonJSVU15} and \textcite{journals/combinatorics/BrandenC11}.
It is not possible to apply the first and second Dominating Rules to the pattern
itself, since when applying the rules we consider containers of the pattern inside
the avoiders of the dominating pattern. For example if we were to attempt to apply
the first rule to the pattern \(\perm{1,2}\) then we would have to consider
containers of \(\perm{1,2}\) inside \(\av{\perm{1,2}}\), and obviously this can
never occur.

It is also possible to take sets of mesh patterns instead of a single mesh pattern
when considering dominating rules, and expressing coincidence between these sets.
Doing this may give nice enumerative results.

It would be interesting to consider a systematic explanation of Wilf-equivalences
amongst classes where \(\perm{3,2,1}\) is the dominating pattern using the
construction presented in \cite[Sec.~12]{2015arXiv151203226B}, in order to directly
reach enumeration and hopefully establish some of the non-trivial Wilf-equivalences
between classes with different dominating patterns. It is possible to show that
the sets consisting of \textpattern{}{1,2}{1/2,2/0,0/0,1/0,1/1,2/1} and
\(\perm{2,3,1}\), or \(\perm{3,2,1}\), are Wilf-equivalent. 


%\chapter{Methods}

\lipsum[14-20]
%%% Local Variables: 
%%% mode: latex
%%% TeX-master: "msc-tannock-2016"
%%% End: 
%%RUM: "Methods"
%%\part{The Second Part}
%\chapter{Results}

In this section you discuss any issues that came up while developing
the system.  If you found something particularly interesting,
difficult, or an important learning experience, put it here.  This is
also a good place to put additional figures and data.

\lipsum[28-34]

%%% Local Variables: 
%%% mode: latex
%%% TeX-master: "msc-tannock-2016"
%%% End: 
%%RUM: "Results"
%\chapter{Discussion}

\lipsum[35-41]

\section{Summary}
\lipsum[42-43]

\section{Conclusion\label{sec:conclusions}}
\lipsum[44-50]
%%% Local Variables: 
%%% mode: latex
%%% TeX-master: "msc-tannock-2016"
%%% End: 
%%RUM: "Discussion"



%% ---------------------------------------------------------------
\printbibliography{} %%RUM: "References"

%% If appendices are needed, uncomment the following line
%% and include the appendices in separate files
\appendix{}%%RUM: "Appendicies (as appropriate)
\chapter{Equivalence classes of mesh patterns}
\section{Coincidence classes with no dominating pattern}
\begin{center}
\begin{equation}
	\begin{gathered}
		\textpattern{}{1,2}{}
		\textpattern{}{1,2}{0/0}
		\textpattern{}{1,2}{0/1}
		\textpattern{}{1,2}{0/1,0/0}
		\textpattern{}{1,2}{1/0}
		\textpattern{}{1,2}{1/0,0/0}
		\textpattern{}{1,2}{1/1}
		\textpattern{}{1,2}{0/0,1/1}
		\textpattern{}{1,2}{0/1,1/1}
		\textpattern{}{1,2}{0/1,0/0,1/1}\\
		\textpattern{}{1,2}{1/0,1/1}
		\textpattern{}{1,2}{1/0,0/0,1/1}
		\textpattern{}{1,2}{1/2}
		\textpattern{}{1,2}{1/2,0/0}
		\textpattern{}{1,2}{0/1,1/2}
		\textpattern{}{1,2}{0/1,1/2,0/0}
		\textpattern{}{1,2}{1/2,1/0}
		\textpattern{}{1,2}{1/2,1/0,0/0}
		\textpattern{}{1,2}{1/2,1/1}
		\textpattern{}{1,2}{1/2,0/0,1/1}\\
		\textpattern{}{1,2}{0/1,1/2,1/1}
		\textpattern{}{1,2}{0/1,1/2,0/0,1/1}
		\textpattern{}{1,2}{1/2,1/0,1/1}
		\textpattern{}{1,2}{1/2,1/0,0/0,1/1}
		\textpattern{}{1,2}{2/1}
		\textpattern{}{1,2}{0/0,2/1}
		\textpattern{}{1,2}{0/1,2/1}
		\textpattern{}{1,2}{0/1,0/0,2/1}
		\textpattern{}{1,2}{1/0,2/1}
		\textpattern{}{1,2}{1/0,0/0,2/1}\\
		\textpattern{}{1,2}{1/1,2/1}
		\textpattern{}{1,2}{0/0,1/1,2/1}
		\textpattern{}{1,2}{0/1,1/1,2/1}
		\textpattern{}{1,2}{0/1,0/0,1/1,2/1}
		\textpattern{}{1,2}{1/0,1/1,2/1}
		\textpattern{}{1,2}{1/0,0/0,1/1,2/1}
		\textpattern{}{1,2}{2/2}
		\textpattern{}{1,2}{0/0,2/2}
		\textpattern{}{1,2}{0/1,2/2}
		\textpattern{}{1,2}{0/1,0/0,2/2}\\
		\textpattern{}{1,2}{1/0,2/2}
		\textpattern{}{1,2}{1/0,0/0,2/2}
		\textpattern{}{1,2}{1/1,2/2}
		\textpattern{}{1,2}{0/1,1/1,2/2}
		\textpattern{}{1,2}{1/0,1/1,2/2}
		\textpattern{}{1,2}{1/2,2/2}
		\textpattern{}{1,2}{1/2,0/0,2/2}
		\textpattern{}{1,2}{0/1,1/2,2/2}
		\textpattern{}{1,2}{0/1,1/2,0/0,2/2}
		\textpattern{}{1,2}{1/2,1/0,2/2}\\
		\textpattern{}{1,2}{1/2,1/0,0/0,2/2}
		\textpattern{}{1,2}{1/2,1/1,2/2}
		\textpattern{}{1,2}{0/1,1/2,1/1,2/2}
		\textpattern{}{1,2}{1/2,1/0,1/1,2/2}
		\textpattern{}{1,2}{2/1,2/2}
		\textpattern{}{1,2}{0/0,2/1,2/2}
		\textpattern{}{1,2}{0/1,2/1,2/2}
		\textpattern{}{1,2}{0/1,0/0,2/1,2/2}
		\textpattern{}{1,2}{1/0,2/1,2/2}
		\textpattern{}{1,2}{1/0,0/0,2/1,2/2}\\
		\textpattern{}{1,2}{1/1,2/1,2/2}
		\textpattern{}{1,2}{0/1,1/1,2/1,2/2}
		\textpattern{}{1,2}{1/0,1/1,2/1,2/2}
	\end{gathered}
	\label{cls:0}
\end{equation}

\begin{equation}
	\begin{gathered}
		\textpattern{}{1,2}{0/2}
		\textpattern{}{1,2}{0/0,0/2}
		\textpattern{}{1,2}{0/1,0/2}
		\textpattern{}{1,2}{0/1,0/0,0/2}
		\textpattern{}{1,2}{1/0,0/2}
		\textpattern{}{1,2}{1/0,0/0,0/2}
		\textpattern{}{1,2}{0/2,1/1}
		\textpattern{}{1,2}{0/0,0/2,1/1}
		\textpattern{}{1,2}{0/1,0/2,1/1}
		\textpattern{}{1,2}{0/1,0/0,0/2,1/1}\\
		\textpattern{}{1,2}{1/0,0/2,1/1}
		\textpattern{}{1,2}{1/2,0/2}
		\textpattern{}{1,2}{1/2,0/0,0/2}
		\textpattern{}{1,2}{0/1,1/2,0/2}
		\textpattern{}{1,2}{0/1,1/2,0/0,0/2}
		\textpattern{}{1,2}{1/2,1/0,0/2}
		\textpattern{}{1,2}{1/2,1/0,0/0,0/2}
		\textpattern{}{1,2}{1/2,0/2,1/1}
		\textpattern{}{1,2}{1/2,0/0,0/2,1/1}
		\textpattern{}{1,2}{0/1,1/2,0/2,1/1}\\
		\textpattern{}{1,2}{0/1,1/2,0/0,0/2,1/1}
		\textpattern{}{1,2}{1/2,1/0,0/2,1/1}
		\textpattern{}{1,2}{0/2,2/1}
		\textpattern{}{1,2}{0/0,0/2,2/1}
		\textpattern{}{1,2}{0/1,0/2,2/1}
		\textpattern{}{1,2}{0/1,0/0,0/2,2/1}
		\textpattern{}{1,2}{0/2,2/1,1/1}
		\textpattern{}{1,2}{0/0,0/2,2/1,1/1}
		\textpattern{}{1,2}{0/1,0/2,2/1,1/1}
		\textpattern{}{1,2}{0/1,0/0,0/2,2/1,1/1}\\
		\textpattern{}{1,2}{0/2,2/2}
		\textpattern{}{1,2}{0/0,0/2,2/2}
		\textpattern{}{1,2}{0/1,0/2,2/2}
		\textpattern{}{1,2}{0/1,0/0,0/2,2/2}
		\textpattern{}{1,2}{1/0,0/2,2/2}
		\textpattern{}{1,2}{1/0,0/0,0/2,2/2}
		\textpattern{}{1,2}{0/2,1/1,2/2}
		\textpattern{}{1,2}{0/1,0/2,1/1,2/2}
		\textpattern{}{1,2}{1/0,0/2,1/1,2/2}
		\textpattern{}{1,2}{1/2,0/2,2/2}\\
		\textpattern{}{1,2}{1/2,0/0,0/2,2/2}
		\textpattern{}{1,2}{0/1,1/2,0/2,2/2}
		\textpattern{}{1,2}{0/1,1/2,0/0,0/2,2/2}
		\textpattern{}{1,2}{1/2,1/0,0/2,2/2}
		\textpattern{}{1,2}{1/2,1/0,0/0,0/2,2/2}
		\textpattern{}{1,2}{1/2,0/2,1/1,2/2}
		\textpattern{}{1,2}{0/1,1/2,0/2,1/1,2/2}
		\textpattern{}{1,2}{1/2,1/0,0/2,1/1,2/2}
		\textpattern{}{1,2}{0/2,2/1,2/2}
		\textpattern{}{1,2}{0/0,0/2,2/1,2/2}\\
		\textpattern{}{1,2}{0/1,0/2,2/1,2/2}
		\textpattern{}{1,2}{0/1,0/0,0/2,2/1,2/2}
	\end{gathered}
	\label{cls:1}
\end{equation}
\noindent\begin{longtabu}[l]{XX}
\begin{equation}
	\begin{gathered}
		\textpattern{}{1,2}{0/1,1/0}
		\textpattern{}{1,2}{0/1,1/0,1/1}
	\end{gathered}
	\label{cls:2}
\end{equation}
    &
\begin{equation}
	\begin{gathered}
		\textpattern{}{1,2}{0/1,1/0,0/0}
		\textpattern{}{1,2}{0/1,1/0,0/0,1/1}
	\end{gathered}
	\label{cls:3}
\end{equation}\\
\begin{equation}
	\begin{gathered}
		\textpattern{}{1,2}{0/1,1/0,0/2}
		\textpattern{}{1,2}{0/1,1/0,0/2,1/1}
		\textpattern{}{1,2}{0/1,1/2,1/0,0/2}
		\textpattern{}{1,2}{0/1,1/2,1/0,0/2,1/1}
	\end{gathered}
	\label{cls:4}
\end{equation}
    &
\begin{equation}
	\begin{gathered}
		\textpattern{}{1,2}{0/1,1/0,0/0,0/2}
		\textpattern{}{1,2}{0/1,1/0,0/0,0/2,1/1}
		\textpattern{}{1,2}{0/1,1/2,1/0,0/0,0/2}
		\textpattern{}{1,2}{0/1,1/2,0/0,1/0,0/2,1/1}
	\end{gathered}
	\label{cls:5}
\end{equation}
\end{longtabu}

\noindent\begin{longtabu}[l]{XXX}
\begin{equation}
	\begin{gathered}
		\textpattern{}{1,2}{1/0,0/0,0/2,1/1}
		\textpattern{}{1,2}{1/2,1/0,0/0,0/2,1/1}
	\end{gathered}
	\label{cls:6}
\end{equation}
    &
\begin{equation}
	\begin{gathered}
		\textpattern{}{1,2}{0/1,1/2,1/0}
		\textpattern{}{1,2}{0/1,1/2,1/0,1/1}
	\end{gathered}
	\label{cls:7}
\end{equation}
    &
\begin{equation}
	\begin{gathered}
		\textpattern{}{1,2}{0/1,1/2,1/0,0/0}
		\textpattern{}{1,2}{0/1,1/2,1/0,0/0,1/1}
	\end{gathered}
	\label{cls:8}
\end{equation}
\end{longtabu}

\begin{equation}
	\begin{gathered}
		\textpattern{}{1,2}{2/0}
		\textpattern{}{1,2}{2/0,0/0}
		\textpattern{}{1,2}{0/1,2/0}
		\textpattern{}{1,2}{0/1,2/0,0/0}
		\textpattern{}{1,2}{2/0,1/0}
		\textpattern{}{1,2}{2/0,1/0,0/0}
		\textpattern{}{1,2}{2/0,1/1}
		\textpattern{}{1,2}{2/0,0/0,1/1}
		\textpattern{}{1,2}{0/1,2/0,1/1}
		\textpattern{}{1,2}{2/0,1/0,1/1}\\
		\textpattern{}{1,2}{2/0,1/0,0/0,1/1}
		\textpattern{}{1,2}{1/2,2/0}
		\textpattern{}{1,2}{1/2,2/0,0/0}
		\textpattern{}{1,2}{1/2,2/0,1/0}
		\textpattern{}{1,2}{1/2,2/0,1/0,0/0}
		\textpattern{}{1,2}{1/2,2/0,1/1}
		\textpattern{}{1,2}{1/2,2/0,0/0,1/1}
		\textpattern{}{1,2}{1/2,2/0,1/0,1/1}
		\textpattern{}{1,2}{1/2,2/0,1/0,0/0,1/1}
		\textpattern{}{1,2}{2/0,2/1}\\
		\textpattern{}{1,2}{2/0,0/0,2/1}
		\textpattern{}{1,2}{0/1,2/0,2/1}
		\textpattern{}{1,2}{0/1,2/0,0/0,2/1}
		\textpattern{}{1,2}{2/0,1/0,2/1}
		\textpattern{}{1,2}{2/0,1/0,0/0,2/1}
		\textpattern{}{1,2}{2/0,1/1,2/1}
		\textpattern{}{1,2}{2/0,0/0,1/1,2/1}
		\textpattern{}{1,2}{0/1,2/0,1/1,2/1}
		\textpattern{}{1,2}{2/0,1/0,1/1,2/1}
		\textpattern{}{1,2}{2/0,1/0,0/0,1/1,2/1}\\
		\textpattern{}{1,2}{2/0,2/2}
		\textpattern{}{1,2}{2/0,0/0,2/2}
		\textpattern{}{1,2}{0/1,2/0,2/2}
		\textpattern{}{1,2}{0/1,2/0,0/0,2/2}
		\textpattern{}{1,2}{2/0,1/0,2/2}
		\textpattern{}{1,2}{2/0,1/0,0/0,2/2}
		\textpattern{}{1,2}{2/0,1/1,2/2}
		\textpattern{}{1,2}{0/1,2/0,1/1,2/2}
		\textpattern{}{1,2}{2/0,1/0,1/1,2/2}
		\textpattern{}{1,2}{1/2,2/0,2/2}\\
		\textpattern{}{1,2}{1/2,2/0,0/0,2/2}
		\textpattern{}{1,2}{1/2,2/0,1/0,2/2}
		\textpattern{}{1,2}{1/2,2/0,1/0,0/0,2/2}
		\textpattern{}{1,2}{2/0,2/1,2/2}
		\textpattern{}{1,2}{2/0,0/0,2/1,2/2}
		\textpattern{}{1,2}{0/1,2/0,2/1,2/2}
		\textpattern{}{1,2}{0/1,2/0,0/0,2/1,2/2}
		\textpattern{}{1,2}{2/0,1/0,2/1,2/2}
		\textpattern{}{1,2}{2/0,1/0,0/0,2/1,2/2}
		\textpattern{}{1,2}{2/0,1/1,2/1,2/2}\\
		\textpattern{}{1,2}{0/1,2/0,1/1,2/1,2/2}
		\textpattern{}{1,2}{2/0,1/0,1/1,2/1,2/2}
	\end{gathered}
	\label{cls:9}
\end{equation}

\begin{equation}
	\begin{gathered}
		\textpattern{}{1,2}{2/0,0/2}
		\textpattern{}{1,2}{2/0,0/0,0/2}
		\textpattern{}{1,2}{2/0,0/2,2/2}
		\textpattern{}{1,2}{2/0,0/0,0/2,2/2}
	\end{gathered}
	\label{cls:10}
\end{equation}

\begin{equation}
	\begin{gathered}
		\textpattern{}{1,2}{0/1,2/0,0/2}
		\textpattern{}{1,2}{0/1,2/0,0/0,0/2}
		\textpattern{}{1,2}{2/0,0/2,2/1}
		\textpattern{}{1,2}{2/0,0/0,0/2,2/1}
		\textpattern{}{1,2}{0/1,2/0,0/2,2/1}
		\textpattern{}{1,2}{0/1,2/0,0/0,0/2,2/1}
		\textpattern{}{1,2}{0/1,2/0,0/2,2/2}
		\textpattern{}{1,2}{0/1,2/0,0/0,0/2,2/2}
		\textpattern{}{1,2}{2/0,0/2,2/1,2/2}
		\textpattern{}{1,2}{2/0,0/0,0/2,2/1,2/2}\\
		\textpattern{}{1,2}{0/1,2/0,0/2,2/1,2/2}
		\textpattern{}{1,2}{0/1,2/0,0/0,2/2,0/2,2/1}
	\end{gathered}
	\label{cls:11}
\end{equation}

\noindent\begin{longtabu}[l]{XX}
\begin{equation}
	\begin{gathered}
		\textpattern{}{1,2}{0/1,2/0,1/0}
		\textpattern{}{1,2}{0/1,2/0,1/0,1/1}
		\textpattern{}{1,2}{0/1,2/0,1/0,2/1}
		\textpattern{}{1,2}{0/1,2/0,1/0,1/1,2/1}
	\end{gathered}
	\label{cls:12}
\end{equation}
    &
\begin{equation}
	\begin{gathered}
		\textpattern{}{1,2}{0/1,2/0,1/0,0/0}
		\textpattern{}{1,2}{0/1,2/0,1/0,0/0,1/1}
		\textpattern{}{1,2}{0/1,2/0,1/0,0/0,2/1}
		\textpattern{}{1,2}{0/1,2/0,0/0,1/0,1/1,2/1}
	\end{gathered}
	\label{cls:13}
\end{equation}
\end{longtabu}

\begin{equation}
	\begin{gathered}
		\textpattern{}{1,2}{2/0,1/0,0/2}
		\textpattern{}{1,2}{2/0,1/0,0/0,0/2}
		\textpattern{}{1,2}{1/2,2/0,0/2}
		\textpattern{}{1,2}{1/2,2/0,0/0,0/2}
		\textpattern{}{1,2}{1/2,2/0,1/0,0/2}
		\textpattern{}{1,2}{1/2,2/0,1/0,0/0,0/2}
		\textpattern{}{1,2}{2/0,1/0,0/2,2/2}
		\textpattern{}{1,2}{2/0,1/0,0/0,0/2,2/2}
		\textpattern{}{1,2}{1/2,2/0,0/2,2/2}
		\textpattern{}{1,2}{1/2,2/0,0/0,0/2,2/2}\\
		\textpattern{}{1,2}{1/2,2/0,1/0,0/2,2/2}
		\textpattern{}{1,2}{1/2,2/0,0/0,2/2,1/0,0/2}
	\end{gathered}
	\label{cls:14}
\end{equation}

\begin{equation}
	\begin{gathered}
		\textpattern{}{1,2}{0/1,2/0,1/0,0/2}
		\textpattern{}{1,2}{0/1,2/0,1/0,0/2,1/1}
		\textpattern{}{1,2}{0/1,1/2,1/0,0/2,2/0}
		\textpattern{}{1,2}{0/1,1/2,2/0,1/0,1/1,0/2}
		\textpattern{}{1,2}{0/1,2/0,1/0,0/2,2/1}
		\textpattern{}{1,2}{0/1,2/0,1/0,0/2,2/1,1/1}
		\textpattern{}{1,2}{0/1,2/0,1/0,0/2,2/2}
		\textpattern{}{1,2}{0/1,1/2,2/0,1/0,2/2,0/2}
		\textpattern{}{1,2}{0/1,2/0,2/2,1/0,0/2,2/1}
	\end{gathered}
	\label{cls:15}
\end{equation}

\begin{equation}
	\begin{gathered}
		\textpattern{}{1,2}{0/1,2/0,1/0,0/0,0/2}
		\textpattern{}{1,2}{0/1,2/0,0/0,1/0,0/2,1/1}
		\textpattern{}{1,2}{0/1,1/2,0/0,2/0,1/0,0/2}
		\textpattern{}{1,2}{0/1,1/2,0/0,2/0,1/0,0/2,1/1}
		\textpattern{}{1,2}{0/1,2/0,0/0,1/0,0/2,2/1}
		\textpattern{}{1,2}{0/1,2/0,0/0,1/0,0/2,2/1,1/1}
		\textpattern{}{1,2}{0/1,2/0,0/0,2/2,1/0,0/2}
		\textpattern{}{1,2}{0/1,1/2,0/0,2/0,1/0,2/2,0/2}
		\textpattern{}{1,2}{0/1,2/0,0/0,2/2,1/0,0/2,2/1}
	\end{gathered}
	\label{cls:16}
\end{equation}
\noindent\begin{longtabu}[l]{XXX}
\begin{equation}
	\begin{gathered}
		\textpattern{}{1,2}{0/1,2/0,0/0,1/1}
		\textpattern{}{1,2}{0/1,2/0,0/0,1/1,2/1}
	\end{gathered}
	\label{cls:17}
\end{equation}
    &
\begin{equation}
	\begin{gathered}
		\textpattern{}{1,2}{2/0,0/2,1/1}
	\end{gathered}
	\label{cls:18}
\end{equation}
    &
\begin{equation}
	\begin{gathered}
		\textpattern{}{1,2}{2/0,0/0,0/2,1/1}
	\end{gathered}
	\label{cls:19}
\end{equation}
\\
\begin{equation}
	\begin{gathered}
		\textpattern{}{1,2}{0/1,2/0,0/2,1/1}
		\textpattern{}{1,2}{2/0,0/2,2/1,1/1}
		\textpattern{}{1,2}{0/1,2/0,0/2,2/1,1/1}
	\end{gathered}
	\label{cls:20}
\end{equation}
    &
\begin{equation}
	\begin{gathered}
		\textpattern{}{1,2}{0/1,2/0,0/0,0/2,1/1}
		\textpattern{}{1,2}{0/1,2/0,0/0,0/2,2/1,1/1}
	\end{gathered}
	\label{cls:21}
\end{equation}
    &
\begin{equation}
	\begin{gathered}
		\textpattern{}{1,2}{2/0,1/0,0/2,1/1}
		\textpattern{}{1,2}{1/2,2/0,0/2,1/1}
		\textpattern{}{1,2}{1/2,2/0,1/0,0/2,1/1}
	\end{gathered}
	\label{cls:22}
\end{equation}
\\
\begin{equation}
	\begin{gathered}
		\textpattern{}{1,2}{2/0,1/0,0/0,0/2,1/1}
		\textpattern{}{1,2}{1/2,2/0,0/0,1/0,0/2,1/1}
	\end{gathered}
	\label{cls:23}
\end{equation}
    &
\begin{equation}
	\begin{gathered}
		\textpattern{}{1,2}{0/1,1/2,2/0}
	\end{gathered}
	\label{cls:24}
\end{equation}
    &
\begin{equation}
	\begin{gathered}
		\textpattern{}{1,2}{0/1,1/2,0/0,2/0}
	\end{gathered}
	\label{cls:25}
\end{equation}
\end{longtabu}

\noindent\begin{longtabu}[l]{XX}
\begin{equation}
	\begin{gathered}
		\textpattern{}{1,2}{0/1,1/2,0/2,2/0}
		\textpattern{}{1,2}{0/1,1/2,0/0,0/2,2/0}
		\textpattern{}{1,2}{0/1,1/2,2/0,0/2,2/2}
		\textpattern{}{1,2}{0/1,1/2,0/0,2/0,2/2,0/2}
	\end{gathered}
	\label{cls:26}
\end{equation}
    &
\begin{equation}
	\begin{gathered}
		\textpattern{}{1,2}{0/1,1/2,1/0,2/0}
		\textpattern{}{1,2}{0/1,1/2,1/0,1/1,2/0}
	\end{gathered}
	\label{cls:27}
\end{equation}
\end{longtabu}

\noindent\begin{longtabu}[l]{XXX}
\begin{equation}
	\begin{gathered}
		\textpattern{}{1,2}{0/1,1/2,1/0,0/0,2/0}
		\textpattern{}{1,2}{0/1,1/2,0/0,2/0,1/0,1/1}
	\end{gathered}
	\label{cls:28}
\end{equation}
    &
\begin{equation}
	\begin{gathered}
		\textpattern{}{1,2}{0/1,1/2,1/1,2/0}
	\end{gathered}
	\label{cls:29}
\end{equation}
    &
\begin{equation}
	\begin{gathered}
		\textpattern{}{1,2}{0/1,1/2,0/0,1/1,2/0}
	\end{gathered}
	\label{cls:30}
\end{equation}
\\
\begin{equation}
	\begin{gathered}
		\textpattern{}{1,2}{1/2,2/0,0/0,0/2,1/1}
	\end{gathered}
	\label{cls:31}
\end{equation}
    &
\begin{equation}
	\begin{gathered}
		\textpattern{}{1,2}{0/1,1/2,0/0,2/0,1/1,0/2}
	\end{gathered}
	\label{cls:32}
\end{equation}
    &
\begin{equation}
	\begin{gathered}
		\textpattern{}{1,2}{0/1,1/0,2/1}
		\textpattern{}{1,2}{0/1,1/0,1/1,2/1}
	\end{gathered}
	\label{cls:33}
\end{equation}
\\
\begin{equation}
	\begin{gathered}
		\textpattern{}{1,2}{0/1,1/0,0/0,2/1}
		\textpattern{}{1,2}{0/1,1/0,0/0,1/1,2/1}
	\end{gathered}
	\label{cls:34}
\end{equation}
    &
\begin{equation}
	\begin{gathered}
		\textpattern{}{1,2}{1/0,0/2,2/1}
	\end{gathered}
	\label{cls:35}
\end{equation}
    &
\begin{equation}
	\begin{gathered}
		\textpattern{}{1,2}{1/0,0/0,0/2,2/1}
	\end{gathered}
	\label{cls:36}
\end{equation}
\\
\begin{equation}
	\begin{gathered}
		\textpattern{}{1,2}{0/1,1/0,0/2,2/1}
		\textpattern{}{1,2}{0/1,1/0,0/2,2/1,1/1}
	\end{gathered}
	\label{cls:37}
\end{equation}
    &
\begin{equation}
	\begin{gathered}
		\textpattern{}{1,2}{0/1,1/0,0/0,0/2,2/1}
		\textpattern{}{1,2}{0/1,0/0,1/0,0/2,2/1,1/1}
	\end{gathered}
	\label{cls:38}
\end{equation}
    &
\begin{equation}
	\begin{gathered}
		\textpattern{}{1,2}{1/0,0/2,2/1,1/1}
	\end{gathered}
	\label{cls:39}
\end{equation}
\\
\begin{equation}
	\begin{gathered}
		\textpattern{}{1,2}{1/0,0/0,0/2,2/1,1/1}
	\end{gathered}
	\label{cls:40}
\end{equation}
    &
\begin{equation}
	\begin{gathered}
		\textpattern{}{1,2}{1/2,2/1}
		\textpattern{}{1,2}{1/2,1/1,2/1}
	\end{gathered}
	\label{cls:41}
\end{equation}
    &
\begin{equation}
	\begin{gathered}
		\textpattern{}{1,2}{1/2,0/0,2/1}
	\end{gathered}
	\label{cls:42}
\end{equation}
\\
\begin{equation}
	\begin{gathered}
		\textpattern{}{1,2}{0/1,1/2,2/1}
		\textpattern{}{1,2}{0/1,1/2,1/1,2/1}
	\end{gathered}
	\label{cls:43}
\end{equation}
    &
\begin{equation}
	\begin{gathered}
		\textpattern{}{1,2}{0/1,1/2,0/0,2/1}
	\end{gathered}
	\label{cls:44}
\end{equation}
    &
\begin{equation}
	\begin{gathered}
		\textpattern{}{1,2}{1/2,0/2,2/1}
		\textpattern{}{1,2}{0/1,1/2,0/2,2/1}
		\textpattern{}{1,2}{1/2,0/2,2/1,1/1}
		\textpattern{}{1,2}{0/1,1/2,0/2,2/1,1/1}
	\end{gathered}
	\label{cls:45}
\end{equation}
\\
\begin{equation}
	\begin{gathered}
		\textpattern{}{1,2}{1/2,0/0,0/2,2/1}
		\textpattern{}{1,2}{0/1,1/2,0/0,0/2,2/1}
	\end{gathered}
	\label{cls:46}
\end{equation}
    &
\begin{equation}
	\begin{gathered}
		\textpattern{}{1,2}{1/2,1/0,2/1}
		\textpattern{}{1,2}{1/2,1/0,1/1,2/1}
	\end{gathered}
	\label{cls:47}
\end{equation}
    &
\begin{equation}
	\begin{gathered}
		\textpattern{}{1,2}{1/2,1/0,0/0,2/1}
	\end{gathered}
	\label{cls:48}
\end{equation}
\\
\begin{equation}
	\begin{gathered}
		\textpattern{}{1,2}{0/1,1/2,1/0,2/1}
	\end{gathered}
	\label{cls:49}
\end{equation}
    &
\begin{equation}
	\begin{gathered}
		\textpattern{}{1,2}{0/1,1/2,1/0,0/0,2/1}
	\end{gathered}
	\label{cls:50}
\end{equation}
    &
\begin{equation}
	\begin{gathered}
		\textpattern{}{1,2}{1/2,1/0,0/2,2/1}
		\textpattern{}{1,2}{1/2,1/0,0/2,2/1,1/1}
	\end{gathered}
	\label{cls:51}
\end{equation}
\\
\begin{equation}
	\begin{gathered}
		\textpattern{}{1,2}{1/2,1/0,0/0,0/2,2/1}
	\end{gathered}
	\label{cls:52}
\end{equation}
    &
\begin{equation}
	\begin{gathered}
		\textpattern{}{1,2}{0/1,1/2,1/0,0/2,2/1}
	\end{gathered}
	\label{cls:53}
\end{equation}
    &
\begin{equation}
	\begin{gathered}
		\textpattern{}{1,2}{0/1,1/2,0/0,1/0,0/2,2/1}
	\end{gathered}
	\label{cls:54}
\end{equation}
\\
\begin{equation}
	\begin{gathered}
		\textpattern{}{1,2}{1/2,0/0,1/1,2/1}
	\end{gathered}
	\label{cls:55}
\end{equation}
    &
\begin{equation}
	\begin{gathered}
		\textpattern{}{1,2}{0/1,1/2,0/0,1/1,2/1}
	\end{gathered}
	\label{cls:56}
\end{equation}
    &
\begin{equation}
	\begin{gathered}
		\textpattern{}{1,2}{1/2,0/0,0/2,2/1,1/1}
	\end{gathered}
	\label{cls:57}
\end{equation}
\\
\begin{equation}
	\begin{gathered}
		\textpattern{}{1,2}{0/1,1/2,0/0,0/2,2/1,1/1}
	\end{gathered}
	\label{cls:58}
\end{equation}
&
\begin{equation}
	\begin{gathered}
		\textpattern{}{1,2}{1/2,1/0,0/0,1/1,2/1}
	\end{gathered}
	\label{cls:59}
\end{equation}
&
\begin{equation}
	\begin{gathered}
		\textpattern{}{1,2}{0/1,1/2,1/0,1/1,2/1}
	\end{gathered}
	\label{cls:60}
\end{equation}
\\
\begin{equation}
	\begin{gathered}
		\textpattern{}{1,2}{0/1,1/2,0/0,1/0,1/1,2/1}
	\end{gathered}
	\label{cls:61}
\end{equation}
&
\begin{equation}
	\begin{gathered}
		\textpattern{}{1,2}{1/2,0/0,1/0,0/2,2/1,1/1}
	\end{gathered}
	\label{cls:62}
\end{equation}
&
\begin{equation}
	\begin{gathered}
		\textpattern{}{1,2}{0/1,1/2,1/0,0/2,2/1,1/1}
	\end{gathered}
	\label{cls:63}
\end{equation}\\
\begin{equation}
	\begin{gathered}
		\textpattern{}{1,2}{0/1,1/2,0/0,1/0,0/2,2/1,1/1}
	\end{gathered}
	\label{cls:64}
\end{equation}
&
\begin{equation}
	\begin{gathered}
		\textpattern{}{1,2}{2/0,1/0,0/2,2/1}
		\textpattern{}{1,2}{2/0,1/0,0/0,0/2,2/1}
		\textpattern{}{1,2}{2/0,1/0,0/2,2/1,2/2}
		\textpattern{}{1,2}{2/0,0/0,2/2,1/0,0/2,2/1}
	\end{gathered}
	\label{cls:65}
\end{equation}
&
\begin{equation}
	\begin{gathered}
		\textpattern{}{1,2}{2/0,0/0,0/2,2/1,1/1}
	\end{gathered}
	\label{cls:66}
\end{equation}
\\
\begin{equation}
	\begin{gathered}
		\textpattern{}{1,2}{0/1,1/2,0/2,1/1,2/0}
		\textpattern{}{1,2}{2/0,1/0,0/2,2/1,1/1}
	\end{gathered}
	\label{cls:67}
\end{equation}
    &
\begin{equation}
	\begin{gathered}
		\textpattern{}{1,2}{2/0,0/0,1/0,0/2,2/1,1/1}
	\end{gathered}
	\label{cls:68}
\end{equation}
    &
\begin{equation}
	\begin{gathered}
		\textpattern{}{1,2}{1/2,2/0,2/1}
		\textpattern{}{1,2}{1/2,2/0,1/0,2/1}
		\textpattern{}{1,2}{1/2,2/0,1/1,2/1}
		\textpattern{}{1,2}{1/2,2/0,1/0,1/1,2/1}
	\end{gathered}
	\label{cls:69}
\end{equation}
\\
\begin{equation}
	\begin{gathered}
		\textpattern{}{1,2}{1/2,2/0,0/0,2/1}
		\textpattern{}{1,2}{1/2,2/0,1/0,0/0,2/1}
	\end{gathered}
	\label{cls:70}
\end{equation}
    &
\begin{equation}
	\begin{gathered}
		\textpattern{}{1,2}{0/1,1/2,2/1,2/0}
		\textpattern{}{1,2}{0/1,1/2,1/1,2/1,2/0}
	\end{gathered}
	\label{cls:71}
\end{equation}
&
\begin{equation}
	\begin{gathered}
		\textpattern{}{1,2}{0/1,1/2,0/0,2/1,2/0}
	\end{gathered}
	\label{cls:72}
\end{equation}
\end{longtabu}

\begin{equation}
	\begin{gathered}
		\textpattern{}{1,2}{1/2,2/0,0/2,2/1}
		\textpattern{}{1,2}{1/2,2/0,0/0,0/2,2/1}
		\textpattern{}{1,2}{0/1,1/2,0/2,2/1,2/0}
		\textpattern{}{1,2}{0/1,1/2,0/0,2/0,0/2,2/1}
		\textpattern{}{1,2}{1/2,2/0,1/0,0/2,2/1}
		\textpattern{}{1,2}{1/2,2/0,0/0,1/0,0/2,2/1}
		\textpattern{}{1,2}{1/2,2/0,0/2,2/1,1/1}
		\textpattern{}{1,2}{0/1,1/2,2/0,1/1,2/1,0/2}
		\textpattern{}{1,2}{1/2,2/0,1/0,0/2,2/1,1/1}
	\end{gathered}
	\label{cls:73}
\end{equation}

\begin{longtabu}[l]{XXX}
\begin{equation}
	\begin{gathered}
		\textpattern{}{1,2}{0/1,1/2,1/0,2/1,2/0}
	\end{gathered}
	\label{cls:74}
\end{equation}
    &
\begin{equation}
	\begin{gathered}
		\textpattern{}{1,2}{0/1,1/2,0/0,2/0,1/0,2/1}
	\end{gathered}
	\label{cls:75}
\end{equation}
    &
\begin{equation}
	\begin{gathered}
		\textpattern{}{1,2}{0/1,1/2,2/0,1/0,0/2,2/1}
	\end{gathered}
	\label{cls:76}
\end{equation}
\\
\begin{equation}
	\begin{gathered}
		\textpattern{}{1,2}{0/1,1/2,0/0,2/0,1/0,0/2,2/1}
	\end{gathered}
	\label{cls:77}
\end{equation}
    &
\begin{equation}
	\begin{gathered}
		\textpattern{}{1,2}{1/2,2/0,0/0,1/1,2/1}
	\end{gathered}
	\label{cls:78}
\end{equation}
    &
\begin{equation}
	\begin{gathered}
		\textpattern{}{1,2}{0/1,1/2,0/0,2/0,1/1,2/1}
	\end{gathered}
	\label{cls:79}
\end{equation}
\\
\begin{equation}
	\begin{gathered}
		\textpattern{}{1,2}{1/2,2/0,0/0,0/2,2/1,1/1}
	\end{gathered}
	\label{cls:80}
\end{equation}
    &
\begin{equation}
	\begin{gathered}
		\textpattern{}{1,2}{0/1,1/2,0/0,2/0,0/2,2/1,1/1}
	\end{gathered}
	\label{cls:81}
\end{equation}
    &
\begin{equation}
	\begin{gathered}
		\textpattern{}{1,2}{1/2,2/0,0/0,1/0,1/1,2/1}
	\end{gathered}
	\label{cls:82}
\end{equation}
\\
\begin{equation}
	\begin{gathered}
		\textpattern{}{1,2}{0/1,1/2,2/0,1/0,1/1,2/1}
	\end{gathered}
	\label{cls:83}
\end{equation}
&
\begin{equation}
	\begin{gathered}
		\textpattern{}{1,2}{0/1,1/2,0/0,2/0,1/0,1/1,2/1}
	\end{gathered}
	\label{cls:84}
\end{equation}
&
\begin{equation}
	\begin{gathered}
		\textpattern{}{1,2}{1/2,2/0,0/0,1/0,0/2,2/1,1/1}
	\end{gathered}
	\label{cls:85}
\end{equation}
\\
\begin{equation}
	\begin{gathered}
		\textpattern{}{1,2}{0/1,1/2,2/0,1/0,0/2,2/1,1/1}
	\end{gathered}
	\label{cls:86}
\end{equation}
&
\begin{equation}
	\begin{gathered}
		\textpattern{}{1,2}{0/1,1/2,0/0,2/1,0/2,2/0,1/0,1/1}
	\end{gathered}
	\label{cls:87}
\end{equation}
&
\begin{equation}
	\begin{gathered}
		\textpattern{}{1,2}{0/1,1/0,2/2}
	\end{gathered}
	\label{cls:88}
\end{equation}
\\
\begin{equation}
	\begin{gathered}
		\textpattern{}{1,2}{0/1,1/0,0/0,2/2}
	\end{gathered}
	\label{cls:89}
\end{equation}
    &
\begin{equation}
	\begin{gathered}
		\textpattern{}{1,2}{0/1,1/0,0/2,2/2}
		\textpattern{}{1,2}{0/1,1/2,1/0,0/2,2/2}
	\end{gathered}
	\label{cls:90}
\end{equation}
    &
\begin{equation}
	\begin{gathered}
		\textpattern{}{1,2}{0/1,1/0,0/0,0/2,2/2}
		\textpattern{}{1,2}{0/1,1/2,0/0,2/2,1/0,0/2}
	\end{gathered}
	\label{cls:91}
\end{equation}
\\
\begin{equation}
	\begin{gathered}
		\textpattern{}{1,2}{0/0,1/1,2/2}
	\end{gathered}
	\label{cls:92}
\end{equation}
    &
\begin{equation}
	\begin{gathered}
		\textpattern{}{1,2}{0/1,0/0,1/1,2/2}
	\end{gathered}
	\label{cls:93}
\end{equation}
    &
\begin{equation}
	\begin{gathered}
		\textpattern{}{1,2}{0/0,0/2,1/1,2/2}
	\end{gathered}
	\label{cls:94}
\end{equation}
\\
\begin{equation}
	\begin{gathered}
		\textpattern{}{1,2}{0/1,0/0,0/2,1/1,2/2}
	\end{gathered}
	\label{cls:95}
\end{equation}
&
\begin{equation}
	\begin{gathered}
		\textpattern{}{1,2}{1/0,0/0,1/1,2/2}
	\end{gathered}
	\label{cls:96}
\end{equation}
&
\begin{equation}
	\begin{gathered}
		\textpattern{}{1,2}{0/1,1/0,1/1,2/2}
	\end{gathered}
	\label{cls:97}
\end{equation}
\\
\begin{equation}
	\begin{gathered}
		\textpattern{}{1,2}{0/1,1/0,0/0,1/1,2/2}
	\end{gathered}
	\label{cls:98}
\end{equation}
&
\begin{equation}
	\begin{gathered}
		\textpattern{}{1,2}{1/0,0/0,0/2,1/1,2/2}
	\end{gathered}
	\label{cls:99}
\end{equation}
&
\begin{equation}
	\begin{gathered}
		\textpattern{}{1,2}{0/1,1/0,0/2,1/1,2/2}
	\end{gathered}
	\label{cls:100}
\end{equation}
\\
\begin{equation}
	\begin{gathered}
		\textpattern{}{1,2}{0/1,0/0,2/2,1/0,1/1,0/2}
	\end{gathered}
	\label{cls:101}
\end{equation}
&
\begin{equation}
	\begin{gathered}
		\textpattern{}{1,2}{0/1,1/2,1/0,2/2}
	\end{gathered}
	\label{cls:102}
\end{equation}
&
\begin{equation}
	\begin{gathered}
		\textpattern{}{1,2}{0/1,1/2,1/0,0/0,2/2}
	\end{gathered}
	\label{cls:103}
\end{equation}
\\
\begin{equation}
	\begin{gathered}
		\textpattern{}{1,2}{1/2,0/0,1/1,2/2}
	\end{gathered}
	\label{cls:104}
\end{equation}
&
\begin{equation}
	\begin{gathered}
		\textpattern{}{1,2}{0/1,1/2,0/0,1/1,2/2}
	\end{gathered}
	\label{cls:105}
\end{equation}
&
\begin{equation}
	\begin{gathered}
		\textpattern{}{1,2}{1/2,0/0,0/2,1/1,2/2}
	\end{gathered}
	\label{cls:106}
\end{equation}
\\
\begin{equation}
	\begin{gathered}
		\textpattern{}{1,2}{0/1,1/2,0/0,2/2,1/1,0/2}
	\end{gathered}
	\label{cls:107}
\end{equation}
&
\begin{equation}
	\begin{gathered}
		\textpattern{}{1,2}{1/2,1/0,0/0,1/1,2/2}
	\end{gathered}
	\label{cls:108}
\end{equation}
&
\begin{equation}
	\begin{gathered}
		\textpattern{}{1,2}{0/1,1/2,1/0,1/1,2/2}
	\end{gathered}
	\label{cls:109}
\end{equation}
\\
\begin{equation}
	\begin{gathered}
		\textpattern{}{1,2}{0/1,1/2,0/0,2/2,1/0,1/1}
	\end{gathered}
	\label{cls:110}
\end{equation}
&
\begin{equation}
	\begin{gathered}
		\textpattern{}{1,2}{1/2,0/0,2/2,1/0,1/1,0/2}
	\end{gathered}
	\label{cls:111}
\end{equation}
&
\begin{equation}
	\begin{gathered}
		\textpattern{}{1,2}{0/1,1/2,2/2,1/0,1/1,0/2}
	\end{gathered}
	\label{cls:112}
\end{equation}
\\
\begin{equation}
	\begin{gathered}
		\textpattern{}{1,2}{0/1,1/2,0/0,2/2,1/0,0/2,1/1}
	\end{gathered}
	\label{cls:113}
\end{equation}
    &
\begin{equation}
	\begin{gathered}
		\textpattern{}{1,2}{0/1,2/0,1/0,2/2}
		\textpattern{}{1,2}{0/1,2/0,1/0,2/1,2/2}
	\end{gathered}
	\label{cls:114}
\end{equation}
    &
\begin{equation}
	\begin{gathered}
		\textpattern{}{1,2}{0/1,2/0,1/0,0/0,2/2}
		\textpattern{}{1,2}{0/1,2/0,0/0,2/2,1/0,2/1}
	\end{gathered}
	\label{cls:115}
\end{equation}
\\
\begin{equation}
	\begin{gathered}
		\textpattern{}{1,2}{2/0,0/0,1/1,2/2}
	\end{gathered}
	\label{cls:116}
\end{equation}
    &
\begin{equation}
	\begin{gathered}
		\textpattern{}{1,2}{0/1,2/0,0/0,1/1,2/2}
	\end{gathered}
	\label{cls:117}
\end{equation}
    &
\begin{equation}
	\begin{gathered}
		\textpattern{}{1,2}{2/0,0/2,1/1,2/2}
	\end{gathered}
	\label{cls:118}
\end{equation}
\\
\begin{equation}
	\begin{gathered}
		\textpattern{}{1,2}{2/0,0/0,0/2,1/1,2/2}
	\end{gathered}
	\label{cls:119}
\end{equation}
    &
\begin{equation}
	\begin{gathered}
		\textpattern{}{1,2}{0/1,2/0,0/2,1/1,2/2}
	\end{gathered}
	\label{cls:120}
\end{equation}
    &
\begin{equation}
	\begin{gathered}
		\textpattern{}{1,2}{0/1,2/0,0/0,2/2,1/1,0/2}
	\end{gathered}
	\label{cls:121}
\end{equation}
\\
\begin{equation}
	\begin{gathered}
		\textpattern{}{1,2}{2/0,1/0,0/0,1/1,2/2}
	\end{gathered}
	\label{cls:122}
\end{equation}
&
\begin{equation}
	\begin{gathered}
		\textpattern{}{1,2}{0/1,2/0,1/0,1/1,2/2}
	\end{gathered}
	\label{cls:123}
\end{equation}
&
\begin{equation}
	\begin{gathered}
		\textpattern{}{1,2}{0/1,2/0,0/0,2/2,1/0,1/1}
	\end{gathered}
	\label{cls:124}
\end{equation}
\\
\begin{equation}
	\begin{gathered}
		\textpattern{}{1,2}{2/0,1/0,0/2,1/1,2/2}
	\end{gathered}
	\label{cls:125}
\end{equation}
&
\begin{equation}
	\begin{gathered}
		\textpattern{}{1,2}{2/0,0/0,2/2,1/0,1/1,0/2}
	\end{gathered}
	\label{cls:126}
\end{equation}
&
\begin{equation}
	\begin{gathered}
		\textpattern{}{1,2}{0/1,2/0,2/2,1/0,1/1,0/2}
	\end{gathered}
	\label{cls:127}
\end{equation}
\\
\begin{equation}
	\begin{gathered}
		\textpattern{}{1,2}{0/1,2/0,0/0,2/2,1/0,0/2,1/1}
	\end{gathered}
	\label{cls:128}
\end{equation}
&
\begin{equation}
	\begin{gathered}
		\textpattern{}{1,2}{0/1,1/2,2/0,2/2}
	\end{gathered}
	\label{cls:129}
\end{equation}
&
\begin{equation}
	\begin{gathered}
		\textpattern{}{1,2}{0/1,1/2,0/0,2/0,2/2}
	\end{gathered}
	\label{cls:130}
\end{equation}
\\
\begin{equation}
	\begin{gathered}
		\textpattern{}{1,2}{0/1,1/2,1/0,2/0,2/2}
	\end{gathered}
	\label{cls:131}
\end{equation}
&
\begin{equation}
	\begin{gathered}
		\textpattern{}{1,2}{0/1,1/2,0/0,2/0,1/0,2/2}
	\end{gathered}
	\label{cls:132}
\end{equation}
&
\begin{equation}
	\begin{gathered}
		\textpattern{}{1,2}{1/2,2/0,1/1,2/2}
		\textpattern{}{1,2}{1/2,2/0,1/0,1/1,2/2}
	\end{gathered}
	\label{cls:133}
\end{equation}
\\
\begin{equation}
	\begin{gathered}
		\textpattern{}{1,2}{1/2,2/0,0/0,1/1,2/2}
	\end{gathered}
	\label{cls:134}
\end{equation}
&
\begin{equation}
	\begin{gathered}
		\textpattern{}{1,2}{0/1,1/2,2/0,1/1,2/2}
	\end{gathered}
	\label{cls:135}
\end{equation}
&
\begin{equation}
	\begin{gathered}
		\textpattern{}{1,2}{0/1,1/2,0/0,2/0,2/2,1/1}
	\end{gathered}
	\label{cls:136}
\end{equation}
\\
\begin{equation}
	\begin{gathered}
		\textpattern{}{1,2}{1/2,2/0,0/2,1/1,2/2}
		\textpattern{}{1,2}{1/2,2/0,2/2,1/0,1/1,0/2}
	\end{gathered}
	\label{cls:137}
\end{equation}
&
\begin{equation}
	\begin{gathered}
		\textpattern{}{1,2}{1/2,2/0,0/0,2/2,1/1,0/2}
	\end{gathered}
	\label{cls:138}
\end{equation}
&
\begin{equation}
	\begin{gathered}
		\textpattern{}{1,2}{0/1,1/2,2/0,2/2,1/1,0/2}
	\end{gathered}
	\label{cls:139}
\end{equation}
\\
\begin{equation}
	\begin{gathered}
		\textpattern{}{1,2}{0/1,1/2,0/0,2/0,2/2,0/2,1/1}
	\end{gathered}
	\label{cls:140}
\end{equation}
    &
\begin{equation}
	\begin{gathered}
		\textpattern{}{1,2}{1/2,2/0,0/0,2/2,1/0,1/1}
	\end{gathered}
	\label{cls:141}
\end{equation}
    &
\begin{equation}
	\begin{gathered}
		\textpattern{}{1,2}{0/1,1/2,2/0,1/0,2/2,1/1}
	\end{gathered}
	\label{cls:142}
\end{equation}
\\
\begin{equation}
	\begin{gathered}
		\textpattern{}{1,2}{0/1,1/2,0/0,2/0,1/0,2/2,1/1}
	\end{gathered}
	\label{cls:143}
\end{equation}
    &
\begin{equation}
	\begin{gathered}
		\textpattern{}{1,2}{1/2,2/0,0/0,2/2,1/0,0/2,1/1}
	\end{gathered}
	\label{cls:144}
\end{equation}
    &
\begin{equation}
	\begin{gathered}
		\textpattern{}{1,2}{0/1,1/2,2/0,1/0,2/2,0/2,1/1}
	\end{gathered}
	\label{cls:145}
\end{equation}
\\
\begin{equation}
	\begin{gathered}
		\textpattern{}{1,2}{0/1,1/2,0/0,0/2,2/0,2/2,1/0,1/1}
	\end{gathered}
	\label{cls:146}
\end{equation}
    &
\begin{equation}
	\begin{gathered}
		\textpattern{}{1,2}{0/1,1/0,2/1,2/2}
	\end{gathered}
	\label{cls:147}
\end{equation}
    &
\begin{equation}
	\begin{gathered}
		\textpattern{}{1,2}{0/1,1/0,0/0,2/1,2/2}
	\end{gathered}
	\label{cls:148}
\end{equation}
\\
\begin{equation}
	\begin{gathered}
		\textpattern{}{1,2}{1/0,0/2,2/1,2/2}
	\end{gathered}
	\label{cls:149}
\end{equation}
&
\begin{equation}
	\begin{gathered}
		\textpattern{}{1,2}{1/0,0/0,0/2,2/1,2/2}
	\end{gathered}
	\label{cls:150}
\end{equation}
&
\begin{equation}
	\begin{gathered}
		\textpattern{}{1,2}{0/1,1/0,0/2,2/1,2/2}
	\end{gathered}
	\label{cls:151}
\end{equation}
\\
\begin{equation}
	\begin{gathered}
		\textpattern{}{1,2}{0/1,0/0,2/2,1/0,0/2,2/1}
	\end{gathered}
	\label{cls:152}
\end{equation}
&
\begin{equation}
	\begin{gathered}
		\textpattern{}{1,2}{0/0,1/1,2/1,2/2}
	\end{gathered}
	\label{cls:153}
\end{equation}
&
\begin{equation}
	\begin{gathered}
		\textpattern{}{1,2}{0/1,0/0,1/1,2/1,2/2}
	\end{gathered}
	\label{cls:154}
\end{equation}
\\
\begin{equation}
	\begin{gathered}
		\textpattern{}{1,2}{0/2,1/1,2/1,2/2}
		\textpattern{}{1,2}{0/1,0/2,1/1,2/1,2/2}
	\end{gathered}
	\label{cls:155}
\end{equation}
&
\begin{equation}
	\begin{gathered}
		\textpattern{}{1,2}{0/0,0/2,1/1,2/1,2/2}
	\end{gathered}
	\label{cls:156}
\end{equation}
&
\begin{equation}
	\begin{gathered}
		\textpattern{}{1,2}{0/1,0/0,2/2,1/1,2/1,0/2}
	\end{gathered}
	\label{cls:157}
\end{equation}
\\
\begin{equation}
	\begin{gathered}
		\textpattern{}{1,2}{1/0,0/0,1/1,2/1,2/2}
	\end{gathered}
	\label{cls:158}
\end{equation}
&
\begin{equation}
	\begin{gathered}
		\textpattern{}{1,2}{0/1,1/0,1/1,2/1,2/2}
	\end{gathered}
	\label{cls:159}
\end{equation}
&
\begin{equation}
	\begin{gathered}
		\textpattern{}{1,2}{0/1,0/0,2/2,1/0,1/1,2/1}
	\end{gathered}
	\label{cls:160}
\end{equation}
\\
\begin{equation}
	\begin{gathered}
		\textpattern{}{1,2}{1/0,0/2,1/1,2/1,2/2}
	\end{gathered}
	\label{cls:161}
\end{equation}
&
\begin{equation}
	\begin{gathered}
		\textpattern{}{1,2}{0/0,2/2,1/0,1/1,2/1,0/2}
	\end{gathered}
	\label{cls:162}
\end{equation}
&
\begin{equation}
	\begin{gathered}
		\textpattern{}{1,2}{0/1,2/2,1/0,1/1,2/1,0/2}
	\end{gathered}
	\label{cls:163}
\end{equation}
\\
\begin{equation}
	\begin{gathered}
		\textpattern{}{1,2}{0/1,0/0,2/2,1/0,0/2,2/1,1/1}
	\end{gathered}
	\label{cls:164}
\end{equation}
&
\begin{equation}
	\begin{gathered}
		\textpattern{}{1,2}{1/2,2/1,2/2}
		\textpattern{}{1,2}{1/2,1/1,2/1,2/2}
	\end{gathered}
	\label{cls:165}
\end{equation}
&
\begin{equation}
	\begin{gathered}
		\textpattern{}{1,2}{1/2,0/0,2/1,2/2}
	\end{gathered}
	\label{cls:166}
\end{equation}
\\
\begin{equation}
	\begin{gathered}
		\textpattern{}{1,2}{0/1,1/2,2/1,2/2}
		\textpattern{}{1,2}{0/1,1/2,1/1,2/1,2/2}
	\end{gathered}
	\label{cls:167}
\end{equation}
    &
\begin{equation}
	\begin{gathered}
		\textpattern{}{1,2}{0/1,1/2,0/0,2/1,2/2}
	\end{gathered}
	\label{cls:168}
\end{equation}
    &
\begin{equation}
	\begin{gathered}
		\textpattern{}{1,2}{1/2,0/2,2/1,2/2}
		\textpattern{}{1,2}{0/1,1/2,0/2,2/1,2/2}
		\textpattern{}{1,2}{1/2,0/2,1/1,2/1,2/2}
		\textpattern{}{1,2}{0/1,1/2,2/2,1/1,2/1,0/2}
	\end{gathered}
	\label{cls:169}
\end{equation}
  \\
\begin{equation}
	\begin{gathered}
		\textpattern{}{1,2}{1/2,0/0,0/2,2/1,2/2}
		\textpattern{}{1,2}{0/1,1/2,0/0,2/2,0/2,2/1}
	\end{gathered}
	\label{cls:170}
\end{equation}
    &
\begin{equation}
	\begin{gathered}
		\textpattern{}{1,2}{1/2,1/0,2/1,2/2}
		\textpattern{}{1,2}{1/2,1/0,1/1,2/1,2/2}
	\end{gathered}
	\label{cls:171}
\end{equation}
    &
\begin{equation}
	\begin{gathered}
		\textpattern{}{1,2}{1/2,1/0,0/0,2/1,2/2}
	\end{gathered}
	\label{cls:172}
\end{equation}
\\
\begin{equation}
	\begin{gathered}
		\textpattern{}{1,2}{0/1,1/2,1/0,2/1,2/2}
	\end{gathered}
	\label{cls:173}
\end{equation}
    &
\begin{equation}
	\begin{gathered}
		\textpattern{}{1,2}{0/1,1/2,0/0,2/2,1/0,2/1}
	\end{gathered}
	\label{cls:174}
\end{equation}
&
\begin{equation}
	\begin{gathered}
		\textpattern{}{1,2}{1/2,1/0,0/2,2/1,2/2}
		\textpattern{}{1,2}{1/2,2/2,1/0,1/1,2/1,0/2}
	\end{gathered}
	\label{cls:175}
\end{equation}
\\
\begin{equation}
	\begin{gathered}
		\textpattern{}{1,2}{1/2,0/0,2/2,1/0,0/2,2/1}
	\end{gathered}
	\label{cls:176}
\end{equation}
&
\begin{equation}
	\begin{gathered}
		\textpattern{}{1,2}{0/1,1/2,2/2,1/0,0/2,2/1}
	\end{gathered}
	\label{cls:177}
\end{equation}
&
\begin{equation}
	\begin{gathered}
		\textpattern{}{1,2}{0/1,1/2,0/0,2/2,1/0,0/2,2/1}
	\end{gathered}
	\label{cls:178}
\end{equation}
\\
\begin{equation}
	\begin{gathered}
		\textpattern{}{1,2}{1/2,0/0,1/1,2/1,2/2}
	\end{gathered}
	\label{cls:179}
\end{equation}
&
\begin{equation}
	\begin{gathered}
		\textpattern{}{1,2}{0/1,1/2,0/0,2/2,1/1,2/1}
	\end{gathered}
	\label{cls:180}
\end{equation}
&
\begin{equation}
	\begin{gathered}
		\textpattern{}{1,2}{1/2,0/0,2/2,1/1,2/1,0/2}
	\end{gathered}
	\label{cls:181}
\end{equation}
\\
\begin{equation}
	\begin{gathered}
		\textpattern{}{1,2}{0/1,1/2,0/0,2/2,0/2,2/1,1/1}
	\end{gathered}
	\label{cls:182}
\end{equation}
&
\begin{equation}
	\begin{gathered}
		\textpattern{}{1,2}{1/2,0/0,2/2,1/0,1/1,2/1}
	\end{gathered}
	\label{cls:183}
\end{equation}
&
\begin{equation}
	\begin{gathered}
		\textpattern{}{1,2}{0/1,1/2,2/2,1/0,1/1,2/1}
	\end{gathered}
	\label{cls:184}
\end{equation}
\\
\begin{equation}
	\begin{gathered}
		\textpattern{}{1,2}{0/1,1/2,0/0,2/2,1/0,1/1,2/1}
	\end{gathered}
	\label{cls:185}
\end{equation}
&
\begin{equation}
	\begin{gathered}
		\textpattern{}{1,2}{1/2,0/0,2/2,1/0,0/2,2/1,1/1}
	\end{gathered}
	\label{cls:186}
\end{equation}
&
\begin{equation}
	\begin{gathered}
		\textpattern{}{1,2}{0/1,1/2,2/2,1/0,0/2,2/1,1/1}
	\end{gathered}
	\label{cls:187}
\end{equation}
\\
\begin{equation}
	\begin{gathered}
		\textpattern{}{1,2}{0/1,1/2,0/0,2/1,0/2,2/2,1/0,1/1}
	\end{gathered}
	\label{cls:188}
\end{equation}
&
\begin{equation}
	\begin{gathered}
		\textpattern{}{1,2}{2/0,0/0,1/1,2/1,2/2}
	\end{gathered}
	\label{cls:189}
\end{equation}
&
\begin{equation}
	\begin{gathered}
		\textpattern{}{1,2}{0/1,2/0,0/0,2/2,1/1,2/1}
	\end{gathered}
	\label{cls:190}
\end{equation}
\\
\begin{equation}
	\begin{gathered}
		\textpattern{}{1,2}{2/0,0/2,1/1,2/1,2/2}
		\textpattern{}{1,2}{0/1,2/0,2/2,1/1,2/1,0/2}
	\end{gathered}
	\label{cls:191}
\end{equation}
    &
\begin{equation}
	\begin{gathered}
		\textpattern{}{1,2}{2/0,0/0,2/2,1/1,2/1,0/2}
	\end{gathered}
	\label{cls:192}
\end{equation}
    &
\begin{equation}
	\begin{gathered}
		\textpattern{}{1,2}{0/1,2/0,0/0,2/2,0/2,2/1,1/1}
	\end{gathered}
	\label{cls:193}
\end{equation}
\\
\begin{equation}
	\begin{gathered}
		\textpattern{}{1,2}{2/0,0/0,2/2,1/0,1/1,2/1}
	\end{gathered}
	\label{cls:194}
\end{equation}
    &
\begin{equation}
	\begin{gathered}
		\textpattern{}{1,2}{0/1,2/0,2/2,1/0,1/1,2/1}
	\end{gathered}
	\label{cls:195}
\end{equation}
    &
\begin{equation}
	\begin{gathered}
		\textpattern{}{1,2}{0/1,2/0,0/0,2/2,1/0,1/1,2/1}
	\end{gathered}
	\label{cls:196}
\end{equation}
\\
\begin{equation}
	\begin{gathered}
		\textpattern{}{1,2}{2/0,2/2,1/0,1/1,2/1,0/2}
	\end{gathered}
	\label{cls:197}
\end{equation}
    &
\begin{equation}
	\begin{gathered}
		\textpattern{}{1,2}{2/0,0/0,2/2,1/0,0/2,2/1,1/1}
	\end{gathered}
	\label{cls:198}
\end{equation}
    &
\begin{equation}
	\begin{gathered}
		\textpattern{}{1,2}{0/1,2/0,2/2,1/0,0/2,2/1,1/1}
	\end{gathered}
	\label{cls:199}
\end{equation}
\\
\begin{equation}
	\begin{gathered}
		\textpattern{}{1,2}{0/1,0/0,2/1,0/2,2/0,2/2,1/0,1/1}
	\end{gathered}
	\label{cls:200}
\end{equation}
&
\begin{equation}
	\begin{gathered}
		\textpattern{}{1,2}{1/2,2/0,2/1,2/2}
		\textpattern{}{1,2}{1/2,2/0,1/0,2/1,2/2}
		\textpattern{}{1,2}{1/2,2/0,1/1,2/1,2/2}
		\textpattern{}{1,2}{1/2,2/0,2/2,1/0,1/1,2/1}
	\end{gathered}
	\label{cls:201}
\end{equation}
&
\begin{equation}
	\begin{gathered}
		\textpattern{}{1,2}{1/2,2/0,0/0,2/1,2/2}
		\textpattern{}{1,2}{1/2,2/0,0/0,2/2,1/0,2/1}
	\end{gathered}
	\label{cls:202}
\end{equation}
\end{longtabu}

\begin{longtabu}[l]{XX}
\begin{equation}
	\begin{gathered}
		\textpattern{}{1,2}{0/1,1/2,2/0,2/1,2/2}
		\textpattern{}{1,2}{0/1,1/2,2/0,2/2,1/1,2/1}
	\end{gathered}
	\label{cls:203}
\end{equation}
    &
\begin{equation}
	\begin{gathered}
		\textpattern{}{1,2}{0/1,1/2,0/0,2/0,2/2,2/1}
	\end{gathered}
	\label{cls:204}
\end{equation}
\end{longtabu}

\begin{equation}
	\begin{gathered}
		\textpattern{}{1,2}{1/2,2/0,0/2,2/1,2/2}
		\textpattern{}{1,2}{1/2,2/0,0/0,2/2,0/2,2/1}
		\textpattern{}{1,2}{0/1,1/2,2/0,2/2,0/2,2/1}
		\textpattern{}{1,2}{0/1,1/2,0/0,2/0,2/2,0/2,2/1}
		\textpattern{}{1,2}{1/2,2/0,2/2,1/0,0/2,2/1}
		\textpattern{}{1,2}{1/2,2/0,0/0,2/2,1/0,0/2,2/1}
		\textpattern{}{1,2}{1/2,2/0,2/2,1/1,2/1,0/2}
		\textpattern{}{1,2}{0/1,1/2,2/0,2/2,0/2,2/1,1/1}
		\textpattern{}{1,2}{1/2,2/0,2/2,1/0,0/2,2/1,1/1}
	\end{gathered}
	\label{cls:205}
\end{equation}

\begin{longtabu}[l]{XXX}
\begin{equation}
	\begin{gathered}
		\textpattern{}{1,2}{0/1,1/2,2/0,1/0,2/2,2/1}
	\end{gathered}
	\label{cls:206}
\end{equation}
    &
\begin{equation}
	\begin{gathered}
		\textpattern{}{1,2}{0/1,1/2,0/0,2/0,1/0,2/2,2/1}
	\end{gathered}
	\label{cls:207}
\end{equation}
    &
\begin{equation}
	\begin{gathered}
		\textpattern{}{1,2}{0/1,1/2,2/0,1/0,2/2,0/2,2/1}
	\end{gathered}
	\label{cls:208}
\end{equation}
\\
\begin{equation}
	\begin{gathered}
		\textpattern{}{1,2}{0/1,1/2,0/0,2/1,2/0,2/2,1/0,0/2}
	\end{gathered}
	\label{cls:209}
\end{equation}
    &
\begin{equation}
	\begin{gathered}
		\textpattern{}{1,2}{1/2,2/0,0/0,2/2,1/1,2/1}
	\end{gathered}
	\label{cls:210}
\end{equation}
    &
\begin{equation}
	\begin{gathered}
		\textpattern{}{1,2}{0/1,1/2,0/0,2/0,2/2,1/1,2/1}
	\end{gathered}
	\label{cls:211}
\end{equation}
\\
\begin{equation}
	\begin{gathered}
		\textpattern{}{1,2}{1/2,2/0,0/0,2/2,0/2,2/1,1/1}
	\end{gathered}
	\label{cls:212}
\end{equation}
    &
\begin{equation}
	\begin{gathered}
		\textpattern{}{1,2}{0/1,1/2,0/0,2/1,0/2,2/0,2/2,1/1}
	\end{gathered}
	\label{cls:213}
\end{equation}
    &
\begin{equation}
	\begin{gathered}
		\textpattern{}{1,2}{1/2,2/0,0/0,2/2,1/0,1/1,2/1}
	\end{gathered}
	\label{cls:214}
\end{equation}
\\
\begin{equation}
	\begin{gathered}
		\textpattern{}{1,2}{0/1,1/2,2/0,1/0,2/2,1/1,2/1}
	\end{gathered}
	\label{cls:215}
\end{equation}
    &
\begin{equation}
	\begin{gathered}
		\textpattern{}{1,2}{0/1,1/2,0/0,2/1,2/0,2/2,1/0,1/1}
	\end{gathered}
	\label{cls:216}
\end{equation}
    &
\begin{equation}
	\begin{gathered}
		\textpattern{}{1,2}{1/2,0/0,2/1,0/2,2/0,2/2,1/0,1/1}
	\end{gathered}
	\label{cls:217}
\end{equation}
\end{longtabu}
\begin{longtabu}[l]{XX}
\begin{equation}
	\begin{gathered}
		\textpattern{}{1,2}{0/1,1/2,2/1,0/2,2/0,2/2,1/0,1/1}
	\end{gathered}
	\label{cls:218}
\end{equation}
    &
\begin{equation}
	\begin{gathered}
		\textpattern{}{1,2}{0/1,1/2,0/0,2/1,0/2,2/0,2/2,1/0,1/1}
	\end{gathered}
	\label{cls:219}
\end{equation}
\end{longtabu}
\end{center}

\section{Consolidation of classes by Dominating Pattern rules}
Each of the lines in the following tables are the sets of classes
that are obtained by successive application of each of the Dominating
Rules, only those coincidences that are not already calculated are shown.
\subsection{First Dominating Rule}
\label{sub:First Dominating Rule}
\subsubsection{Dominating pattern 321}
\begin{center}
\begin{longtabu}{|X[l]|X[2,l]|}
    \hline
    \multicolumn{2}{|c|}{Mesh pattern family}\\
    \hline
    \multicolumn{1}{|c|}{12} & \multicolumn{1}{|c|}{21} \\
    \hline
    &1, 93, 94, 97, 105, 106, 109, 154, 155, 159\\
    &2, 7, 95, 96, 100, 107, 108, 112, 156, 157, 158\\
    &3, 4, 89, 90, 98, 99\\
    &5, 6, 91, 92, 101, 102, 113, 114\\
    &8, 9, 103, 104, 110, 111\\
    &10, 18, 117, 118, 123, 134, 135, 142, 190, 191, 195\\
    &11, 19, 20, 119, 120\\
    &12, 21, 22, 67, 121, 122, 192, 193, 194\\
    &13, 14, 115, 116, 124, 125, 196, 197\\
    &15, 23, 24, 32, 126, 127, 138, 139, 145\\
    &16, 17, 128, 129, 146, 147, 200, 201\\
    &25, 26, 30, 31, 130, 131, 136, 137\\
    &27, 33, 66, 68, 69, 140, 141, 198, 199\\
    &28, 29, 132, 133, 143, 144\\
    &34, 35, 148, 149, 160, 161\\
    &36, 37, 40, 41, 150, 151, 162, 163\\
    &38, 39, 152, 153, 164, 165\\
    &42, 43, 56, 166, 167, 180\\
    &44, 45, 57, 168, 169, 181\\
    &46, 47, 58, 59, 170, 171, 182, 183\\
    &48, 49, 60, 172, 173, 184\\
    &50, 51, 61, 62, 174, 175, 185, 186\\
    &52, 53, 63, 176, 177, 187\\
    &54, 55, 64, 65, 178, 179, 188, 189\\
    &70, 71, 79, 83, 202, 203, 211, 215\\
    &72, 73, 80, 204, 205, 212\\
    &74, 81, 82, 86, 206, 213, 214, 218\\
    &75, 76, 84, 85, 207, 208, 216, 217\\
    &77, 78, 87, 88, 209, 210, 219, 220\\
    \hline
\end{longtabu}
\end{center}
\newpage
\subsubsection{Dominating pattern 231}
\begin{center}
\begin{longtabu}[c]{|X|X|}
    \hline
    \multicolumn{2}{|c|}{Mesh pattern family}\\
    \hline
    \multicolumn{1}{|c|}{12} & \multicolumn{1}{|c|}{21} \\
    \hline
    0, 9, 17, 24, 25, 29, 30, 129, 130, 133, 135&0, 41, 42, 43, 44, 47, 48, 55, 56, 59, 165, 166, 167, 168, 171, 172\\
    1, 10, 11, 14, 18, 19, 20, 21, 22, 26, 31, 32, 39, 66, 67, 118, 120, 125, 137, 139&1, 35, 36, 39, 45, 46, 51, 52, 57, 58, 149, 150, 155, 161, 169, 170, 175, 176\\
    2, 12, 33&2, 7, 33, 49, 60\\
    3, 13, 34&3, 8, 34, 50, 61\\
    4, 15, 37, 90, 151&4, 37, 53, 63\\
    5, 16, 38, 91, 152&5, 38, 54, 64\\
    6, 23&6, 40, 62\\
    7, 27&9, 24, 25, 29, 69, 70, 71, 72, 78, 82, 129, 130, 133, 135, 201, 202, 203, 204\\
    8, 28&10, 11, 14, 18, 20, 22, 26, 65, 67, 73, 118, 120, 125, 137, 139, 191, 197, 205\\
    35, 36, 65, 149, 150&12, 27, 74, 83\\
    40, 68&13, 28, 75, 84\\
    41, 47, 69&15, 76, 86, 208\\
    42, 48, 70&16, 77, 87, 209\\
    43, 71&17, 30, 79\\
    44, 72&19, 31, 66, 80\\
    45, 46, 51, 52, 73&21, 32, 81\\
    49, 74&23, 68, 85\\
    50, 75&88, 102, 147, 173\\
    53, 76&89, 103, 148, 174\\
    54, 77&90, 151, 177\\
    55, 78&91, 152, 178\\
    56, 79&92, 104, 153, 179\\
    57, 80&93, 105, 154, 180\\
    58, 81&94, 106, 156, 181\\
    59, 82&95, 107, 157, 182\\
    60, 83&96, 108, 158, 183\\
    61, 84&97, 109, 159, 184\\
    62, 85&98, 110, 160, 185\\
    63, 86&99, 111, 162, 186\\
    64, 87&100, 112, 163, 187\\
    88, 114, 147&101, 113, 164, 188\\
    89, 115, 148&114, 131, 206\\
    92, 116&115, 132, 207\\
    93, 117&116, 134, 189, 210\\
    94, 119&117, 136, 190, 211\\
    95, 121&119, 138, 192, 212\\
    96, 122&121, 140, 193, 213\\
    97, 123&122, 141, 194, 214\\
    98, 124&123, 142, 195, 215\\
    99, 126&124, 143, 196, 216\\
    100, 127&126, 144, 198, 217\\
    101, 128&127, 145, 199, 218\\
    102, 131&128, 146, 200, 219\\
    103, 132&\\
    104, 134&\\
    105, 136&\\
    106, 138&\\
    107, 140&\\
    108, 141&\\
    109, 142&\\
    110, 143&\\
    111, 144&\\
    112, 145&\\
    113, 146&\\
    153, 189&\\
    154, 190&\\
    155, 191&\\
    156, 192&\\
    157, 193&\\
    158, 194&\\
    159, 195&\\
    160, 196&\\
    161, 197&\\
    162, 198&\\
    163, 199&\\
    164, 200&\\
    165, 171, 201&\\
    166, 172, 202&\\
    167, 203&\\
    168, 204&\\
    169, 170, 175, 176, 205&\\
    173, 206&\\
    174, 207&\\
    177, 208&\\
    178, 209&\\
    179, 210&\\
    180, 211&\\
    181, 212&\\
    182, 213&\\
    183, 214&\\
    184, 215&\\
    185, 216&\\
    186, 217&\\
    187, 218&\\
    188, 219&\\
\hline
\end{longtabu}
\end{center}

\subsection{Second Dominating Rule}
\subsubsection{Dominating pattern 321}
There are no new coincidences when the dominating pattern is \(\perm{3,2,1}\)
when applying the second dominating rule.
\subsubsection{Dominating pattern 231}
\begin{center}
\begin{longtabu}{|X[l]|X[l]|}
    \hline
    \multicolumn{2}{|c|}{Mesh pattern family}\\
    \hline
    \multicolumn{1}{|c|}{12} & \multicolumn{1}{|c|}{21} \\
    \hline
    2, 7&1, 6, 155\\
    3, 8&9, 17, 133\\
    12, 27&10, 18, 19, 118\\
    13, 28&11, 20, 21, 66, 120, 191\\
    41, 43&14, 22, 23, 31, 125, 137\\
    42, 44&15, 127, 145, 199\\
    55, 56&25, 30\\
    57, 58&26, 32, 65, 67, 68, 139, 197\\
    69, 71&36, 40\\
    70, 72&52, 62\\
    78, 79&72, 79\\
    80, 81&73, 80, 81, 85, 205\\
    88, 102&76, 86, 208, 218\\
    89, 103&88, 97, 102, 109, 147, 159\\
    92, 93, 104, 105&90, 100, 112, 151, 163\\
    94, 95, 106, 107&114, 123, 131, 142, 195\\
    96, 108&173, 184\\
    97, 109&177, 187\\
    98, 110&206, 215\\
    99, 111&\\
    100, 112&\\
    101, 113&\\
    114, 131&\\
    115, 132&\\
    122, 141&\\
    123, 142&\\
    124, 143&\\
    126, 144&\\
    127, 145&\\
    128, 146&\\
    153, 154&\\
    156, 157&\\
    165, 167&\\
    166, 168&\\
    179, 180&\\
    181, 182&\\
    189, 190&\\
    192, 193&\\
    201, 203&\\
    202, 204&\\
    210, 211&\\
    212, 213&\\
    \hline
\end{longtabu}
\end{center}

\chapter{Wilf-equivalence data}
\section{Sequences with underlying pattern 231}
\begin{longtabu}{|r|l|c|}
    \hline
    Sequence & Related OEIS entry & Number of patterns in class\\
    \hline
    \endfirsthead
    \hline
    \endhead
    \hline
    \endfoot
    \(1,    1,    1,    1,    1,    1,    1,    1,    1,    1,     1\)  & A000012 &\(210\)\\
    \(1,    1,    1,    2,    4,    9,   21,   51,  127,  323,   835\)  & A001006 &\(32\)\\
    \(1,    1,    1,    2,    5,   13,   35,   97,  275,  794,  2327\)  & A086581 &\(8\)\\
    \(1,    1,    1,    3,    6,   17,   43,  123,  343, 1004,  2938\)  & A143363 &\(2\)\\
    \(1,    1,    1,    2,    5,   14,   42,  132,  429, 1430,  4862\)  & A000108 &\(314\)\\
    \(1,    1,    1,    3,    7,   19,   53,  153,  453, 1367,  4191\)  & A078481 &\(2\)\\
    \(1,    1,    1,    2,    6,   19,   61,  200,  670, 2286,  7918\)  & A035929 &\(32\)\\
    \(1,    1,    1,    2,    6,   20,   68,  233,  805, 2807,  9879\)  & A014138 &\(36\)\\
    \(1,    1,    1,    3,    8,   24,   75,  243,  808, 2742,  9458\)  & A000958 &\(64\)\\
    \(1,    1,    1,    2,    7,   25,   85,  285,  964, 3310, 11527\)  &  &\(4\)\\
    \(1,    1,    1,    3,    9,   28,   90,  297, 1001, 3432, 11934\)  & A000245 &\(176\)\\
    \(1,    1,    1,    4,   10,   31,   97,  316, 1054, 3586, 12394\)  &  &\(2\)\\
    \(1,    1,    1,    3,    9,   29,   95,  317, 1075, 3699, 12891\)  &  &\(4\)\\
    \(1,    1,    1,    3,   10,   31,   98,  321, 1078, 3686, 12789\)  & A114487 &\(4\)\\
    \(1,    1,    1,    2,    7,   26,   93,  325, 1129, 3935, 13813\)  & A014140 &\(8\)\\
    \(1,    1,    1,    4,   11,   33,  105,  343, 1148, 3916, 13563\)  & A127154 &\(2\)\\
    \(1,    1,    1,    4,   11,   34,  108,  354, 1187, 4054, 14054\)  & A000958 &\(8\)\\
    \(1,    1,    1,    3,   10,   33,  109,  364, 1233, 4236, 14740\)  & A014137 &\(38\)\\
    \(1,    1,    1,    4,   12,   37,  118,  387, 1298, 4433, 15366\)  & A00108 &\(46\)\\
    \(1,    1,    1,    2,    8,   32,  117,  408, 1402, 4826, 16751\)  & A000217 &\(2\)\\
    \(1,    1,    1,    3,   11,   38,  127,  423, 1423, 4854, 16787\)  & A000027 &\(6\)\\
    \(1,    1,    1,    4,   13,   41,  131,  428, 1429, 4861, 16795\)  & A001453 &\(18\)\\
    \(1,    1,    1,    5,   14,   42,  132,  429, 1430, 4862, 16796\)  & A000108 &\(6\)\\
    \hline
\end{longtabu}

\section{Sequences with underlying pattern 321}
\begin{longtabu}{|r|l|c|}
    \hline
    Sequence & Related OEIS entry & Number of patterns in class\\
    \hline
    \endfirsthead
    \hline
    \endhead
    \hline
    \endfoot
    \(    1,    1,    1,    0,    0,    0,    0,    0,    0,    0,    0\) & 	& \(63\)\\
    \(    1,    1,    1,    1,    1,    1,    1,    1,    1,    1,    1\) & A000012 & \(180\)\\
    \(    1,    1,    1,    2,    3,    4,    5,    6,    7,    8,    9\) &	A000027 & \(5\)\\
    \(    1,    1,    1,    1,    2,    3,    6,   11,   22,   44,   90\) &	A007477 & \(8\)\\
    \(    1,    1,    1,    2,    4,    8,   16,   32,   64,  128,  256\) &	A000079 & \(30\)\\
    \(    1,    1,    1,    1,    3,    6,   13,   28,   60,  129,  277\) &	A002478 & \(2\)\\
    \(    1,    1,    1,    2,    3,    9,   16,   48,  102,  289,  693\) &  & \(1\)\\
    \(    1,    1,    1,    2,    4,    9,   21,   51,  127,  323,  835\) &	A001006 & \(17\)\\
    \(    1,    1,    1,    1,    3,    8,   21,   55,  144,  377,  987\) &	A001906 & \(4\)\\
    \(    1,    1,    1,    1,    3,    7,   19,   53,  153,  453, 1367\) &	A078481 & \(2\)\\
    \(    1,    1,    1,    1,    2,    5,   14,   42,  132,  429, 1430\) &	A000108 & \(12\)\\
    \(    1,    1,    1,    1,    3,   10,   30,   84,  227,  603, 1589\) & 	& \(2\)\\
    \(    1,    1,    1,    2,    5,   13,   34,   89,  233,  610, 1597\) &	A001519 & \(8\)\\
    \(    1,    1,    1,    2,    3,    7,   19,   56,  174,  561, 1859\) &	A167422 & \(2\)\\
    \(    1,    1,    1,    2,    5,   13,   36,  103,  303,  910, 2779\) &	 & \(8\)\\
    \(    1,    1,    1,    1,    3,    9,   28,   90,  297, 1001, 3432\) &	A000245 & \(8\)\\
    \(    1,    1,    1,    3,    6,   18,   47,  139,  405, 1225, 3740\) &  & \(2\)\\
    \(    1,    1,    1,    3,    7,   19,   53,  153,  453, 1367, 4191\) & A078481 & \(2\)\\
    \(    1,    1,    1,    2,    4,   11,   34,  110,  365, 1234, 4237\) & 	& \(4\)\\
    \(    1,    1,    1,    1,    3,   10,   33,  111,  379, 1312, 4596\) &	A001558 & \(2\)\\
    \(    1,    1,    1,    2,    5,   14,   42,  132,  429, 1430, 4862\) &	A000108 & \(170\)\\
    \(    1,    1,    1,    1,    4,   12,   39,  129,  436, 1498, 5218\) &	A122920 & \(4\)\\
    \(    1,    1,    1,    2,    5,   15,   48,  159,  538, 1850, 6446\) & 	& \(8\)\\
    \(    1,    1,    1,    1,    4,   14,   48,  165,  572, 2002, 7072\) &	A002057 & \(6\)\\
    \(    1,    1,    1,    2,    6,   18,   57,  186,  622, 2120, 7338\) &	A000957 & \(4\)\\
    \(    1,    1,    1,    2,    6,   19,   61,  200,  670, 2286, 7918\) &	A035929 & \(32\)\\
    \(    1,    1,    1,    1,    5,   17,   57,  193,  662, 2299, 8073\) & 	& \(4\)\\
    \(    1,    1,    1,    2,    6,   19,   61,  202,  683, 2349, 8191\) &	 & \(2\)\\
    \(    1,    1,    1,    2,    6,   19,   62,  207,  704, 2431, 8502\) &	A026012 & \(8\)\\
    \(    1,    1,    1,    3,    8,   24,   75,  243,  808, 2742, 9458\) &	A000958 & \(18\)\\
    \(    1,    1,    1,    2,    7,   22,   71,  235,  794, 2728, 9503\) &	 & \(2\)\\
    \(    1,    1,    1,    2,    6,   20,   68,  233,  805, 2807, 9879\) & A014138 & \(12\)\\
    \(    1,    1,    1,    3,    8,   25,   80,  264,  890, 3053,10622\) &	 & \(4\)\\
    \(    1,    1,    1,    1,    4,   16,   63,  239,  880, 3184,11431\) &	 & \(2\)\\
    \(    1,    1,    1,    2,    8,   26,   85,  283,  959, 3300,11505\) &	 & \(4\)\\
    \(    1,    1,    1,    1,    5,   20,   74,  265,  937, 3304,11678\) &	 & \(4\)\\
    \(    1,    1,    1,    2,    6,   21,   75,  266,  938, 3305,11679\) &	& \(8\)\\
    \(    1,    1,    1,    2,    7,   25,   86,  292,  995, 3425,11926\) &	& \(4\)\\
    \(    1,    1,    1,    3,    9,   28,   90,  297, 1001, 3432,11934\) &	A000245 & \(86\)\\
    \(    1,    1,    1,    4,   10,   31,   97,  316, 1054, 3586,12394\) &	& \(1\)\\
    \(    1,    1,    1,    3,   10,   31,   98,  321, 1078, 3686,12789\) &	A114487 & \(8\)\\
    \(    1,    1,    1,    4,   11,   33,  105,  343, 1148, 3916,13563\) &	A127154 & \(1\)\\
    \(    1,    1,    1,    2,    7,   26,   93,  325, 1129, 3935,13813\) &	A014140 & \(16\)\\
    \(    1,    1,    1,    4,   11,   34,  108,  354, 1187, 4054,14054\) &	A000958 & \(2\)\\
    \(    1,    1,    1,    3,    9,   31,  105,  355, 1210, 4171,14543\) &	& \(2\)\\
    \(    1,    1,    1,    3,   10,   33,  109,  364, 1233, 4236,14740\) &	A014137& \(36\)\\
    \(    1,    1,    1,    2,    9,   33,  113,  381, 1291, 4425,15357\) &	A192480& \(1\)\\
    \(    1,    1,    1,    3,   10,   34,  114,  382, 1292, 4426,15358\) &	& \(2\)\\
    \(    1,    1,    1,    4,   12,   37,  118,  387, 1298, 4433,15366\) &	A000108 & \(19\)\\
    \(    1,    1,    1,    3,    9,   30,  104,  365, 1286, 4542,16092\) &	A045623 & \(8\)\\
    \(    1,    1,    1,    1,    6,   22,   91,  349, 1277, 4570,16235\) &	& \(1\)\\
    \(    1,    1,    1,    2,    6,   25,   96,  357, 1289, 4587,16258\) &	& \(2\)\\
    \(    1,    1,    1,    3,    7,   28,  101,  365, 1301, 4604,16281\) &	& \(1\)\\
    \(    1,    1,    1,    1,    7,   25,  102,  377, 1339, 4699,16496\) &	& \(4\)\\
    \(    1,    1,    1,    2,    7,   28,  106,  382, 1345, 4706,16504\) &	A132109 & \(4\)\\
    \(    1,    1,    1,    1,    8,   28,  108,  387, 1354, 4720,16524\) &	& \(2\)\\
    \(    1,    1,    1,    2,    7,   29,  109,  388, 1355, 4721,16525\) & A081494	& \(4\)\\
    \(    1,    1,    1,    2,    8,   31,  112,  392, 1360, 4727,16532\) &	A006127 & \(4\)\\
    \(    1,    1,    1,    3,    9,   32,  113,  393, 1361, 4728,16533\) &	A132736 & \(4\)\\
    \(    1,    1,    1,    3,   10,   34,  116,  397, 1366, 4734,16540\) &	A000079 & \(50\)\\
    \(    1,    1,    1,    1,    8,   31,  116,  407, 1401, 4825,16750\) &	& \(2\)\\
    \(    1,    1,    1,    2,    8,   32,  117,  408, 1402, 4826,16751\) &	A000217 & \(2\)\\
    \(    1,    1,    1,    2,    9,   34,  122,  417, 1416, 4846,16778\) & A003265	& \(2\)\\
    \(    1,    1,    1,    2,   10,   37,  126,  422, 1422, 4853,16786\) &	& \(2\)\\
    \(    1,    1,    1,    3,   10,   37,  126,  422, 1422, 4853,16786\) &	& \(4\)\\
    \(    1,    1,    1,    2,   11,   37,  126,  422, 1422, 4853,16786\) &	& \(4\)\\
    \(    1,    1,    1,    3,   11,   38,  127,  423, 1423, 4854,16787\) &	A000027 & \(22\)\\
    \(    1,    1,    1,    3,   12,   40,  130,  427, 1428, 4860,16794\) &	& \(4\)\\
    \(    1,    1,    1,    4,   12,   40,  130,  427, 1428, 4860,16794\) &	& \(2\)\\
    \(    1,    1,    1,    3,   13,   40,  130,  427, 1428, 4860,16794\) &	& \(1\)\\
    \(    1,    1,    1,    4,   13,   41,  131,  428, 1429, 4861,16795\) &	A001453 & \(54\)\\
    \(    1,    1,    1,    5,   14,   42,  132,  429, 1430, 4862,16796\) &	A000108 & \(9\)\\
    \hline
\end{longtabu}

\chapter{Code}\label{cha:code}
% \begin{minted}[linenos, breaklines]{python}
% def boring(args=None):
%     pass
% \end{minted}
\inputminted[linenos, breaklines, frame=lines, stripall]{python}{src/main.py}
% \lstinputlisting[language=Python]{src/main.py}
% You can put code in your document using the listings package, which is
% loaded by default in \path{custom.tex}.  Be aware that the listings
% package does not put code in your document if you are in draft mode
% unless you set the \texttt{forcegraphics} option.
%
% There is an example java (Listing~\ref{src:Data_Bus.java}) and XML
% file (Listing~\ref{src:AndroidManifest.xml}).  Thanks to the
% \texttt{url} package, you can typeset OSX and unix paths like this:
% \path{/afs/rnd.ru.is/project/thesis-template}.  Windows paths:
% \path{C:\windows\temp\ }.  You can also typeset them using the menukey
% package, but it tends to delete the last separator and has other
% complications.\footnote{The menukey package has issues with biblatex,
%   read \path{custom.tex} for more information.}
%
% If you are trying to include multiple different languages, you should
% go read the documentation and set these up in \path{custom.tex}.  You
% will save yourself a lot of effort, especially if you have to fix
% anything.
%
% %I have put the source code in the \directory{src/} folder.
% \lstinputlisting[language=Java, firstline=1,
% lastline=40, caption={Data\_Bus.java: Setting up the class.},
% label={src:Data_Bus.java}]{src/Data_Bus.java}
%
% \lstinputlisting[language={[android]XML}, firstline=1, lastline=20,
% caption={AndroidManifest.xml: Configuration for the Android UI.},
% label={src:AndroidManifest.xml}]{src/AndroidManifest.xml}

%%% Local Variables:
%%% mode: latex
%%% TeX-master: "msc-tannock-2016"
%%% End:

% \chapter{Code}\label{cha:code}
% \begin{minted}[linenos, breaklines]{python}
% def boring(args=None):
%     pass
% \end{minted}
\inputminted[linenos, breaklines, frame=lines, stripall]{python}{src/main.py}
% \lstinputlisting[language=Python]{src/main.py}
% You can put code in your document using the listings package, which is
% loaded by default in \path{custom.tex}.  Be aware that the listings
% package does not put code in your document if you are in draft mode
% unless you set the \texttt{forcegraphics} option.
%
% There is an example java (Listing~\ref{src:Data_Bus.java}) and XML
% file (Listing~\ref{src:AndroidManifest.xml}).  Thanks to the
% \texttt{url} package, you can typeset OSX and unix paths like this:
% \path{/afs/rnd.ru.is/project/thesis-template}.  Windows paths:
% \path{C:\windows\temp\ }.  You can also typeset them using the menukey
% package, but it tends to delete the last separator and has other
% complications.\footnote{The menukey package has issues with biblatex,
%   read \path{custom.tex} for more information.}
%
% If you are trying to include multiple different languages, you should
% go read the documentation and set these up in \path{custom.tex}.  You
% will save yourself a lot of effort, especially if you have to fix
% anything.
%
% %I have put the source code in the \directory{src/} folder.
% \lstinputlisting[language=Java, firstline=1,
% lastline=40, caption={Data\_Bus.java: Setting up the class.},
% label={src:Data_Bus.java}]{src/Data_Bus.java}
%
% \lstinputlisting[language={[android]XML}, firstline=1, lastline=20,
% caption={AndroidManifest.xml: Configuration for the Android UI.},
% label={src:AndroidManifest.xml}]{src/AndroidManifest.xml}

%%% Local Variables:
%%% mode: latex
%%% TeX-master: "msc-tannock-2016"
%%% End:
 % as an example, perhaps some of your code

%\backmatter{} % Sections after this don't get numbers
%% We prefer that all elements be numbered

%%%%%%%%%%%%% SHOW GLOSSARY %%%%%%%%%%%%%%%%%
%% Glossary, optional.  A good idea on longer documents
%% Remember to run "makeglossaries <DEGREE-NAME-YEAR>"
%%   to update the glossary
% \printglossary{}%%RUM: Not mentioned
% \printacronyms{}%%RUM: Not mentioned

%% Glossary debugging code
%% Uncomment this to figure out of one of the glossary entries is broken
%% by putting all of the entries in the glossary.
%\ifthenelse{\boolean{debug}}{%
%  \glsaddall{} %put all entries in
%  \printglossaries{}}{}% dump them all
%% This command dumps all of the glossaries

%%%%%%%%%%%%% SHOW INDEX %%%%%%%%%%%%%%%%%%
%% Index, optional.  A good idea on longer documents

% You can put instructions at the beginning of the index:
%\renewcommand{\preindexhook}{%
%  The first page number is usually, but not always,
%  the primary reference to the indexed topic.\vskip\onelineskip}

%% You may have to run "makeindex <FILENAME>" to have it be generated
%% Depending upon which package you chose.
%%
\clearforchapter{}
\printindex{}%%RUM: Not mentioned

\backcover{}%%RUM: "Back cover (standard)
\end{document}

%% ---------------------------------------------------------------

%%% This is information for the emacs editor
%%% Local Variables:
%%% mode: latex
%%% TeX-master: t
%%% End:

%%% You can change TeX-engine if you want to use xetex for extra font powers
